\documentclass[12pt,envcountsame]{llncs}

\pagestyle{plain}

% ESOP 2020 submission - 25.10.2019
% ESOP 2020 final version - 21.02.2020

%% BIBLIOGRAPHY STYLE
\bibliographystyle{splncs04}

\usepackage{charter,eulervm}

%% PACKAGES
\usepackage[T1]{fontenc}
\usepackage[utf8]{inputenc}

\usepackage{etoolbox}
\makeatletter
\let\llncs@addcontentsline\addcontentsline
\patchcmd{\maketitle}{\addcontentsline}{\llncs@addcontentsline}{}{}
\patchcmd{\maketitle}{\addcontentsline}{\llncs@addcontentsline}{}{}
\patchcmd{\maketitle}{\addcontentsline}{\llncs@addcontentsline}{}{}
\setcounter{tocdepth}{4}
\makeatother

\usepackage{amsmath}
\usepackage{amssymb}
\usepackage{bbold} % alternative blackboard letters (amssymb already provides a set of them)
\usepackage{graphicx}
\usepackage{color}
\usepackage{listings}
\usepackage{mathabx}
\usepackage{mathpartir} % inference rules
\usepackage{stmaryrd}
\usepackage{mathtools} % to typeset equations in the appendix
\usepackage{upquote}
\usepackage[dvipsnames]{xcolor}
\usepackage[all]{xy}

% Fix the definition environment (whoever had the bright idea to make it cursive?).
\let\definition\relax % undefine the environment
\spnewtheorem{definition}{Definition}{\bfseries}{\rmfamily}

%\usepackage{amsmath,amsthm,amscd,amssymb}
\usepackage[colorlinks=true
,breaklinks=true
,urlcolor=blue
,anchorcolor=blue
,citecolor=blue
,filecolor=blue
,linkcolor=blue
,menucolor=blue
,linktocpage=true]{hyperref}
\hypersetup{
bookmarksopen=true,
bookmarksnumbered=true,
bookmarksopenlevel=10
}
\usepackage[capitalise,nameinlink]{cleveref} % customisable references

%\usepackage[noBBpl,sc]{mathpazo}
%\linespread{1.05}
%\usepackage[scaled=.92]{helvet}
%\DeclareMathAlphabet      {\mathsf}{OT1}{phv}{m}{n}

\usepackage[papersize={6.9in, 10.0in}, left=.5in, right=.5in, top=1in, bottom=.9in]{geometry}
\tolerance=2000
\hbadness=2000

% these include amsmath and that can cause trouble in older docs.
\makeatletter
\@ifpackageloaded{amsmath}{}{\RequirePackage{amsmath}}

\DeclareFontFamily{U}  {cmex}{}
\DeclareSymbolFont{Csymbols}       {U}  {cmex}{m}{n}
\DeclareFontShape{U}{cmex}{m}{n}{
    <-6>  cmex5
   <6-7>  cmex6
   <7-8>  cmex6
   <8-9>  cmex7
   <9-10> cmex8
  <10-12> cmex9
  <12->   cmex10}{}

\def\Set@Mn@Sym#1{\@tempcnta #1\relax}
\def\Next@Mn@Sym{\advance\@tempcnta 1\relax}
\def\Prev@Mn@Sym{\advance\@tempcnta-1\relax}
\def\@Decl@Mn@Sym#1#2#3#4{\DeclareMathSymbol{#2}{#3}{#4}{#1}}
\def\Decl@Mn@Sym#1#2#3{%
  \if\relax\noexpand#1%
    \let#1\undefined
  \fi
  \expandafter\@Decl@Mn@Sym\expandafter{\the\@tempcnta}{#1}{#3}{#2}%
  \Next@Mn@Sym}
\def\Decl@Mn@Alias#1#2#3{\Prev@Mn@Sym\Decl@Mn@Sym{#1}{#2}{#3}}
\let\Decl@Mn@Char\Decl@Mn@Sym
\def\Decl@Mn@Op#1#2#3{\def#1{\DOTSB#3\slimits@}}
\def\Decl@Mn@Int#1#2#3{\def#1{\DOTSI#3\ilimits@}}

\let\sum\undefined
\DeclareMathSymbol{\tsum}{\mathop}{Csymbols}{"50}
\DeclareMathSymbol{\dsum}{\mathop}{Csymbols}{"51}

\Decl@Mn@Op\sum\dsum\tsum

\makeatother

\makeatletter
\@ifpackageloaded{amsmath}{}{\RequirePackage{amsmath}}

\DeclareFontFamily{OMX}{MnSymbolE}{}
\DeclareSymbolFont{largesymbolsX}{OMX}{MnSymbolE}{m}{n}
\DeclareFontShape{OMX}{MnSymbolE}{m}{n}{
    <-6>  MnSymbolE5
   <6-7>  MnSymbolE6
   <7-8>  MnSymbolE7
   <8-9>  MnSymbolE8
   <9-10> MnSymbolE9
  <10-12> MnSymbolE10
  <12->   MnSymbolE12}{}

\DeclareMathSymbol{\downbrace}    {\mathord}{largesymbolsX}{'251}
\DeclareMathSymbol{\downbraceg}   {\mathord}{largesymbolsX}{'252}
\DeclareMathSymbol{\downbracegg}  {\mathord}{largesymbolsX}{'253}
\DeclareMathSymbol{\downbraceggg} {\mathord}{largesymbolsX}{'254}
\DeclareMathSymbol{\downbracegggg}{\mathord}{largesymbolsX}{'255}
\DeclareMathSymbol{\upbrace}      {\mathord}{largesymbolsX}{'256}
\DeclareMathSymbol{\upbraceg}     {\mathord}{largesymbolsX}{'257}
\DeclareMathSymbol{\upbracegg}    {\mathord}{largesymbolsX}{'260}
\DeclareMathSymbol{\upbraceggg}   {\mathord}{largesymbolsX}{'261}
\DeclareMathSymbol{\upbracegggg}  {\mathord}{largesymbolsX}{'262}
\DeclareMathSymbol{\braceld}      {\mathord}{largesymbolsX}{'263}
\DeclareMathSymbol{\bracelu}      {\mathord}{largesymbolsX}{'264}
\DeclareMathSymbol{\bracerd}      {\mathord}{largesymbolsX}{'265}
\DeclareMathSymbol{\braceru}      {\mathord}{largesymbolsX}{'266}
\DeclareMathSymbol{\bracemd}      {\mathord}{largesymbolsX}{'267}
\DeclareMathSymbol{\bracemu}      {\mathord}{largesymbolsX}{'270}
\DeclareMathSymbol{\bracemid}     {\mathord}{largesymbolsX}{'271}

\def\horiz@expandable#1#2#3#4#5#6#7#8{%
  \@mathmeasure\z@#7{#8}%
  \@tempdima=\wd\z@
  \@mathmeasure\z@#7{#1}%
  \ifdim\noexpand\wd\z@>\@tempdima
    $\m@th#7#1$%
  \else
    \@mathmeasure\z@#7{#2}%
    \ifdim\noexpand\wd\z@>\@tempdima
      $\m@th#7#2$%
    \else
      \@mathmeasure\z@#7{#3}%
      \ifdim\noexpand\wd\z@>\@tempdima
        $\m@th#7#3$%
      \else
        \@mathmeasure\z@#7{#4}%
        \ifdim\noexpand\wd\z@>\@tempdima
          $\m@th#7#4$%
        \else
          \@mathmeasure\z@#7{#5}%
          \ifdim\noexpand\wd\z@>\@tempdima
            $\m@th#7#5$%
          \else
           #6#7%
          \fi
        \fi
      \fi
    \fi
  \fi}

\def\overbrace@expandable#1#2#3{\vbox{\m@th\ialign{##\crcr
  #1#2{#3}\crcr\noalign{\kern2\p@\nointerlineskip}%
  $\m@th\hfil#2#3\hfil$\crcr}}}
\def\underbrace@expandable#1#2#3{\vtop{\m@th\ialign{##\crcr
  $\m@th\hfil#2#3\hfil$\crcr
  \noalign{\kern2\p@\nointerlineskip}%
  #1#2{#3}\crcr}}}

\def\overbrace@#1#2#3{\vbox{\m@th\ialign{##\crcr
  #1#2\crcr\noalign{\kern2\p@\nointerlineskip}%
  $\m@th\hfil#2#3\hfil$\crcr}}}
\def\underbrace@#1#2#3{\vtop{\m@th\ialign{##\crcr
  $\m@th\hfil#2#3\hfil$\crcr
  \noalign{\kern2\p@\nointerlineskip}%
  #1#2\crcr}}}

\def\bracefill@#1#2#3#4#5{$\m@th#5#1\leaders\hbox{$#4$}\hfill#2\leaders\hbox{$#4$}\hfill#3$}

\def\downbracefill@{\bracefill@\braceld\bracemd\bracerd\bracemid}
\def\upbracefill@{\bracefill@\bracelu\bracemu\braceru\bracemid}

\DeclareRobustCommand{\downbracefill}{\downbracefill@\textstyle}
\DeclareRobustCommand{\upbracefill}{\upbracefill@\textstyle}

\def\upbrace@expandable{%
  \horiz@expandable
    \upbrace
    \upbraceg
    \upbracegg
    \upbraceggg
    \upbracegggg
    \upbracefill@}
\def\downbrace@expandable{%
  \horiz@expandable
    \downbrace
    \downbraceg
    \downbracegg
    \downbraceggg
    \downbracegggg
    \downbracefill@}

\DeclareRobustCommand{\overbrace}[1]{\mathop{\mathpalette{\overbrace@expandable\downbrace@expandable}{#1}}\limits}
\DeclareRobustCommand{\underbrace}[1]{\mathop{\mathpalette{\underbrace@expandable\upbrace@expandable}{#1}}\limits}

\makeatother


%\usepackage[dotinlabels]{titletoc}
%\titlelabel{{\thetitle}.\quad}
%\usepackage{titletoc}
\usepackage[small]{titlesec}

\titleformat{\section}[block]
  {\fillast\medskip}
  {\bfseries{\thesection. }}
  {1ex minus .1ex}
  {\bfseries}
 
\titleformat*{\subsection}{\itshape}
\titleformat*{\subsubsection}{\itshape}

\setcounter{tocdepth}{2}

\titlecontents{section}
              [2.3em] 
              {\bigskip}
              {{\contentslabel{2.3em}}}
              {\hspace*{-2.3em}}
              {\titlerule*[1pc]{}\contentspage}
              
\titlecontents{subsection}
              [4.7em] 
              {}
              {{\contentslabel{2.3em}}}
              {\hspace*{-2.3em}}
              {\titlerule*[.5pc]{}\contentspage}

% hopefully not used.           
\titlecontents{subsubsection}
              [7.9em]
              {}
              {{\contentslabel{3.3em}}}
              {\hspace*{-3.3em}}
              {\titlerule*[.5pc]{}\contentspage}
%\makeatletter
\renewcommand\tableofcontents{%
    \section*{\contentsname
        \@mkboth{%
           \MakeLowercase\contentsname}{\MakeLowercase\contentsname}}%
    \@starttoc{toc}%
    }
\def\@oddhead{{\scshape\rightmark}\hfil{\small\scshape\thepage}}%
\def\sectionmark#1{%
      \markright{\MakeLowercase{%
        \ifnum \c@secnumdepth >\m@ne
          \thesection\quad
        \fi
        #1}}}
        
\makeatother

%\makeatletter

 \def\small{%
  \@setfontsize\small\@xipt{13pt}%
  \abovedisplayskip 8\p@ \@plus3\p@ \@minus6\p@
  \belowdisplayskip \abovedisplayskip
  \abovedisplayshortskip \z@ \@plus3\p@
  \belowdisplayshortskip 6.5\p@ \@plus3.5\p@ \@minus3\p@
  \def\@listi{%
    \leftmargin\leftmargini
    \topsep 9\p@ \@plus3\p@ \@minus5\p@
    \parsep 4.5\p@ \@plus2\p@ \@minus\p@
    \itemsep \parsep
  }%
}%
 \def\footnotesize{%
  \@setfontsize\footnotesize\@xpt{12pt}%
  \abovedisplayskip 10\p@ \@plus2\p@ \@minus5\p@
  \belowdisplayskip \abovedisplayskip
  \abovedisplayshortskip \z@ \@plus3\p@
  \belowdisplayshortskip 6\p@ \@plus3\p@ \@minus3\p@
  \def\@listi{%
    \leftmargin\leftmargini
    \topsep 6\p@ \@plus2\p@ \@minus2\p@
    \parsep 3\p@ \@plus2\p@ \@minus\p@
    \itemsep \parsep
  }%
}%
\def\open@column@one#1{%
 \ltxgrid@info@sw{\class@info{\string\open@column@one\string#1}}{}%
 \unvbox\pagesofar
 \@ifvoid{\footsofar}{}{%
  \insert\footins\bgroup\unvbox\footsofar\egroup
  \penalty\z@
 }%
 \gdef\thepagegrid{one}%
 \global\pagegrid@col#1%
 \global\pagegrid@cur\@ne
 \global\count\footins\@m
 \set@column@hsize\pagegrid@col
 \set@colht
}%

\def\frontmatter@abstractheading{%
\bigskip
 \begingroup
  \centering\large
  \abstractname
  \par\bigskip
 \endgroup
}%

\makeatother

%\DeclareSymbolFont{CMlargesymbols}{OMX}{cmex}{m}{n}
%\DeclareMathSymbol{\sum}{\mathop}{CMlargesymbols}{"50}


\begin{document}

%%%% MACROS FOR NOTATION %%%%
% Use these for any notation where there are multiple options.

%%% Notes and exercise sections
\makeatletter
\newcommand{\sectionNotes}{\phantomsection\section*{Notes}\addcontentsline{toc}{section}{Notes}\markright{\textsc{\@chapapp{} \thechapter{} Notes}}}
\newcommand{\sectionExercises}[1]{\ifdef{\OPTexerciseperpage}{\newpage}{}\phantomsection\section*{Exercises}\addcontentsline{toc}{section}{Exercises}\markright{\textsc{\@chapapp{} \thechapter{} Exercises}}}
\makeatother

%%% Definitional equality (used infix) %%%
\newcommand{\jdeq}{\equiv}      % An equality judgment
\let\judgeq\jdeq
%\newcommand{\defeq}{\coloneqq}  % An equality currently being defined
\newcommand{\defeq}{\vcentcolon\equiv}  % A judgmental equality currently being defined

%%% Term being defined
\newcommand{\define}[1]{\textbf{#1}}

%%% Vec (for example)

\newcommand{\Vect}{\ensuremath{\mathsf{Vec}}}
\newcommand{\Fin}{\ensuremath{\mathsf{Fin}}}
\newcommand{\fmax}{\ensuremath{\mathsf{fmax}}}
\newcommand{\seq}[1]{\langle #1\rangle}

%%% Dependent products %%%
\def\prdsym{\textstyle\prod}
%% Call the macro like \prd{x,y:A}{p:x=y} with any number of
%% arguments.  Make sure that whatever comes *after* the call doesn't
%% begin with an open-brace, or it will be parsed as another argument.
\makeatletter
% Currently the macro is configured to produce
%     {\textstyle\prod}(x:A) \; {\textstyle\prod}(y:B),{\ }
% in display-math mode, and
%     \prod_{(x:A)} \prod_{y:B}
% in text-math mode.
% \def\prd#1{\@ifnextchar\bgroup{\prd@parens{#1}}{%
%     \@ifnextchar\sm{\prd@parens{#1}\@eatsm}{%
%         \prd@noparens{#1}}}}
\def\prd#1{\@ifnextchar\bgroup{\prd@parens{#1}}{%
    \@ifnextchar\sm{\prd@parens{#1}\@eatsm}{%
    \@ifnextchar\prd{\prd@parens{#1}\@eatprd}{%
    \@ifnextchar\;{\prd@parens{#1}\@eatsemicolonspace}{%
    \@ifnextchar\\{\prd@parens{#1}\@eatlinebreak}{%
    \@ifnextchar\narrowbreak{\prd@parens{#1}\@eatnarrowbreak}{%
      \prd@noparens{#1}}}}}}}}
\def\prd@parens#1{\@ifnextchar\bgroup%
  {\mathchoice{\@dprd{#1}}{\@tprd{#1}}{\@tprd{#1}}{\@tprd{#1}}\prd@parens}%
  {\@ifnextchar\sm%
    {\mathchoice{\@dprd{#1}}{\@tprd{#1}}{\@tprd{#1}}{\@tprd{#1}}\@eatsm}%
    {\mathchoice{\@dprd{#1}}{\@tprd{#1}}{\@tprd{#1}}{\@tprd{#1}}}}}
\def\@eatsm\sm{\sm@parens}
\def\prd@noparens#1{\mathchoice{\@dprd@noparens{#1}}{\@tprd{#1}}{\@tprd{#1}}{\@tprd{#1}}}
% Helper macros for three styles
\def\lprd#1{\@ifnextchar\bgroup{\@lprd{#1}\lprd}{\@@lprd{#1}}}
\def\@lprd#1{\mathchoice{{\textstyle\prod}}{\prod}{\prod}{\prod}({\textstyle #1})\;}
\def\@@lprd#1{\mathchoice{{\textstyle\prod}}{\prod}{\prod}{\prod}({\textstyle #1}),\ }
\def\tprd#1{\@tprd{#1}\@ifnextchar\bgroup{\tprd}{}}
\def\@tprd#1{\mathchoice{{\textstyle\prod_{(#1)}}}{\prod_{(#1)}}{\prod_{(#1)}}{\prod_{(#1)}}}
\def\dprd#1{\@dprd{#1}\@ifnextchar\bgroup{\dprd}{}}
\def\@dprd#1{\prod_{(#1)}\,}
\def\@dprd@noparens#1{\prod_{#1}\,}

% Look through spaces and linebreaks
\def\@eatnarrowbreak\narrowbreak{%
  \@ifnextchar\prd{\narrowbreak\@eatprd}{%
    \@ifnextchar\sm{\narrowbreak\@eatsm}{%
      \narrowbreak}}}
\def\@eatlinebreak\\{%
  \@ifnextchar\prd{\\\@eatprd}{%
    \@ifnextchar\sm{\\\@eatsm}{%
      \\}}}
\def\@eatsemicolonspace\;{%
  \@ifnextchar\prd{\;\@eatprd}{%
    \@ifnextchar\sm{\;\@eatsm}{%
      \;}}}

%%% Lambda abstractions.
% Each variable being abstracted over is a separate argument.  If
% there is more than one such argument, they *must* be enclosed in
% braces.  Arguments can be untyped, as in \lam{x}{y}, or typed with a
% colon, as in \lam{x:A}{y:B}. In the latter case, the colons are
% automatically noticed and (with current implementation) the space
% around the colon is reduced.  You can even give more than one variable
% the same type, as in \lam{x,y:A}.
\def\lam#1{{\lambda}\@lamarg#1:\@endlamarg\@ifnextchar\bgroup{.\,\lam}{.\,}}
\def\@lamarg#1:#2\@endlamarg{\if\relax\detokenize{#2}\relax #1\else\@lamvar{\@lameatcolon#2},#1\@endlamvar\fi}
\def\@lamvar#1,#2\@endlamvar{(#2\,{:}\,#1)}
% \def\@lamvar#1,#2{{#2}^{#1}\@ifnextchar,{.\,{\lambda}\@lamvar{#1}}{\let\@endlamvar\relax}}
\def\@lameatcolon#1:{#1}
\let\lamt\lam
% This version silently eats any typing annotation.
\def\lamu#1{{\lambda}\@lamuarg#1:\@endlamuarg\@ifnextchar\bgroup{.\,\lamu}{.\,}}
\def\@lamuarg#1:#2\@endlamuarg{#1}

%%% Dependent products written with \forall, in the same style
\def\fall#1{\forall (#1)\@ifnextchar\bgroup{.\,\fall}{.\,}}

%%% Existential quantifier %%%
\def\exis#1{\exists (#1)\@ifnextchar\bgroup{.\,\exis}{.\,}}

%%% Dependent sums %%%
\def\smsym{\textstyle\sum}
% Use in the same way as \prd
\def\sm#1{\@ifnextchar\bgroup{\sm@parens{#1}}{%
    \@ifnextchar\prd{\sm@parens{#1}\@eatprd}{%
    \@ifnextchar\sm{\sm@parens{#1}\@eatsm}{%
    \@ifnextchar\;{\sm@parens{#1}\@eatsemicolonspace}{%
    \@ifnextchar\\{\sm@parens{#1}\@eatlinebreak}{%
    \@ifnextchar\narrowbreak{\sm@parens{#1}\@eatnarrowbreak}{%
        \sm@noparens{#1}}}}}}}}
\def\sm@parens#1{\@ifnextchar\bgroup%
  {\mathchoice{\@dsm{#1}}{\@tsm{#1}}{\@tsm{#1}}{\@tsm{#1}}\sm@parens}%
  {\@ifnextchar\prd%
    {\mathchoice{\@dsm{#1}}{\@tsm{#1}}{\@tsm{#1}}{\@tsm{#1}}\@eatprd}%
    {\mathchoice{\@dsm{#1}}{\@tsm{#1}}{\@tsm{#1}}{\@tsm{#1}}}}}
\def\@eatprd\prd{\prd@parens}
\def\sm@noparens#1{\mathchoice{\@dsm@noparens{#1}}{\@tsm{#1}}{\@tsm{#1}}{\@tsm{#1}}}
\def\lsm#1{\@ifnextchar\bgroup{\@lsm{#1}\lsm}{\@@lsm{#1}}}
\def\@lsm#1{\mathchoice{{\textstyle\sum}}{\sum}{\sum}{\sum}({\textstyle #1})\;}
\def\@@lsm#1{\mathchoice{{\textstyle\sum}}{\sum}{\sum}{\sum}({\textstyle #1}),\ }
\def\tsm#1{\@tsm{#1}\@ifnextchar\bgroup{\tsm}{}}
\def\@tsm#1{\mathchoice{{\textstyle\sum_{(#1)}}}{\sum_{(#1)}}{\sum_{(#1)}}{\sum_{(#1)}}}
\def\dsm#1{\@dsm{#1}\@ifnextchar\bgroup{\dsm}{}}
\def\@dsm#1{\sum_{(#1)}\,}
\def\@dsm@noparens#1{\sum_{#1}\,}

%%% W-types
\def\wtypesym{{\mathsf{W}}}
\def\wtype#1{\@ifnextchar\bgroup%
  {\mathchoice{\@twtype{#1}}{\@twtype{#1}}{\@twtype{#1}}{\@twtype{#1}}\wtype}%
  {\mathchoice{\@twtype{#1}}{\@twtype{#1}}{\@twtype{#1}}{\@twtype{#1}}}}
\def\lwtype#1{\@ifnextchar\bgroup{\@lwtype{#1}\lwtype}{\@@lwtype{#1}}}
\def\@lwtype#1{\mathchoice{{\textstyle\mathsf{W}}}{\mathsf{W}}{\mathsf{W}}{\mathsf{W}}({\textstyle #1})\;}
\def\@@lwtype#1{\mathchoice{{\textstyle\mathsf{W}}}{\mathsf{W}}{\mathsf{W}}{\mathsf{W}}({\textstyle #1}),\ }
\def\twtype#1{\@twtype{#1}\@ifnextchar\bgroup{\twtype}{}}
\def\@twtype#1{\mathchoice{{\textstyle\mathsf{W}_{(#1)}}}{\mathsf{W}_{(#1)}}{\mathsf{W}_{(#1)}}{\mathsf{W}_{(#1)}}}
\def\dwtype#1{\@dwtype{#1}\@ifnextchar\bgroup{\dwtype}{}}
\def\@dwtype#1{\mathsf{W}_{(#1)}\,}

\newcommand{\suppsym}{{\mathsf{sup}}}
\newcommand{\supp}{\ensuremath\suppsym\xspace}

\def\wtypeh#1{\@ifnextchar\bgroup%
  {\mathchoice{\@lwtypeh{#1}}{\@twtypeh{#1}}{\@twtypeh{#1}}{\@twtypeh{#1}}\wtypeh}%
  {\mathchoice{\@@lwtypeh{#1}}{\@twtypeh{#1}}{\@twtypeh{#1}}{\@twtypeh{#1}}}}
\def\lwtypeh#1{\@ifnextchar\bgroup{\@lwtypeh{#1}\lwtypeh}{\@@lwtypeh{#1}}}
\def\@lwtypeh#1{\mathchoice{{\textstyle\mathsf{W}^h}}{\mathsf{W}^h}{\mathsf{W}^h}{\mathsf{W}^h}({\textstyle #1})\;}
\def\@@lwtypeh#1{\mathchoice{{\textstyle\mathsf{W}^h}}{\mathsf{W}^h}{\mathsf{W}^h}{\mathsf{W}^h}({\textstyle #1}),\ }
\def\twtypeh#1{\@twtypeh{#1}\@ifnextchar\bgroup{\twtypeh}{}}
\def\@twtypeh#1{\mathchoice{{\textstyle\mathsf{W}^h_{(#1)}}}{\mathsf{W}^h_{(#1)}}{\mathsf{W}^h_{(#1)}}{\mathsf{W}^h_{(#1)}}}
\def\dwtypeh#1{\@dwtypeh{#1}\@ifnextchar\bgroup{\dwtypeh}{}}
\def\@dwtypeh#1{\mathsf{W}^h_{(#1)}\,}


\makeatother

% Other notations related to dependent sums
\let\setof\Set    % from package 'braket', write \setof{ x:A | P(x) }.
\newcommand{\pair}{\ensuremath{\mathsf{pair}}\xspace}
\newcommand{\tup}[2]{(#1,#2)}
\newcommand{\proj}[1]{\ensuremath{\mathsf{pr}_{#1}}\xspace}
\newcommand{\fst}{\ensuremath{\proj1}\xspace}
\newcommand{\snd}{\ensuremath{\proj2}\xspace}
\newcommand{\ac}{\ensuremath{\mathsf{ac}}\xspace} % not needed in symbol index

%%% recursor and induction
\newcommand{\rec}[1]{\mathsf{rec}_{#1}}
\newcommand{\ind}[1]{\mathsf{ind}_{#1}}
\newcommand{\indid}[1]{\ind{=_{#1}}} % (Martin-Lof) path induction principle for identity types
\newcommand{\indidb}[1]{\ind{=_{#1}}'} % (Paulin-Mohring) based path induction principle for identity types

%%% Uniqueness principles
\newcommand{\uniq}[1]{\mathsf{uniq}_{#1}}

% Paths in pairs
\newcommand{\pairpath}{\ensuremath{\mathsf{pair}^{\mathord{=}}}\xspace}
% \newcommand{\projpath}[1]{\proj{#1}^{\mathord{=}}}
\newcommand{\projpath}[1]{\ensuremath{\apfunc{\proj{#1}}}\xspace}
\newcommand{\pairct}{\ensuremath{\mathsf{pair}^{\mathord{\ct}}}\xspace}

%%% For quotients %%%
%\newcommand{\pairr}[1]{{\langle #1\rangle}}
\newcommand{\pairr}[1]{{\mathopen{}(#1)\mathclose{}}}
\newcommand{\Pairr}[1]{{\mathopen{}\left(#1\right)\mathclose{}}}

% \newcommand{\type}{\ensuremath{\mathsf{Type}}} % this command is overridden below, so it's commented out
\newcommand{\im}{\ensuremath{\mathsf{im}}} % the image

%%% 2D path operations
\newcommand{\leftwhisker}{\mathbin{{\ct}_{\mathsf{l}}}}  % was \ell
\newcommand{\rightwhisker}{\mathbin{{\ct}_{\mathsf{r}}}} % was r
\newcommand{\hct}{\star}

%%% modalities %%%
\newcommand{\modal}{\ensuremath{\ocircle}}
\let\reflect\modal
\newcommand{\modaltype}{\ensuremath{\type_\modal}}
% \newcommand{\ism}[1]{\ensuremath{\mathsf{is}_{#1}}}
% \newcommand{\ismodal}{\ism{\modal}}
% \newcommand{\existsmodal}{\ensuremath{{\exists}_{\modal}}}
% \newcommand{\existsmodalunique}{\ensuremath{{\exists!}_{\modal}}}
% \newcommand{\modalfunc}{\textsf{\modal-fun}}
% \newcommand{\Ecirc}{\ensuremath{\mathsf{E}_\modal}}
% \newcommand{\Mcirc}{\ensuremath{\mathsf{M}_\modal}}
\newcommand{\mreturn}{\ensuremath{\eta}}
\let\project\mreturn
%\newcommand{\mbind}[1]{\ensuremath{\hat{#1}}}
\newcommand{\ext}{\mathsf{ext}}
%\newcommand{\mmap}[1]{\ensuremath{\bar{#1}}}
%\newcommand{\mjoin}{\ensuremath{\mreturn^{-1}}}
% Subuniverse
\renewcommand{\P}{\ensuremath{\type_{P}}\xspace}

%%% Localizations
% \newcommand{\islocal}[1]{\ensuremath{\mathsf{islocal}_{#1}}\xspace}
% \newcommand{\loc}[1]{\ensuremath{\mathcal{L}_{#1}}\xspace}

%%% Identity types %%%
\newcommand{\idsym}{{=}}
\newcommand{\id}[3][]{\ensuremath{#2 =_{#1} #3}\xspace}
\newcommand{\idtype}[3][]{\ensuremath{\mathsf{Id}_{#1}(#2,#3)}\xspace}
\newcommand{\idtypevar}[1]{\ensuremath{\mathsf{Id}_{#1}}\xspace}
% A propositional equality currently being defined
\newcommand{\defid}{\coloneqq}

%%% Dependent paths
\newcommand{\dpath}[4]{#3 =^{#1}_{#2} #4}

%%% singleton
% \newcommand{\sgl}{\ensuremath{\mathsf{sgl}}\xspace}
% \newcommand{\sctr}{\ensuremath{\mathsf{sctr}}\xspace}

%%% Reflexivity terms %%%
% \newcommand{\reflsym}{{\mathsf{refl}}}
\newcommand{\refl}[1]{\ensuremath{\mathsf{refl}_{#1}}\xspace}

%%% Path concatenation (used infix, in diagrammatic order) %%%
\newcommand{\ct}{%
  \mathchoice{\mathbin{\raisebox{0.5ex}{$\displaystyle\centerdot$}}}%
             {\mathbin{\raisebox{0.5ex}{$\centerdot$}}}%
             {\mathbin{\raisebox{0.25ex}{$\scriptstyle\,\centerdot\,$}}}%
             {\mathbin{\raisebox{0.1ex}{$\scriptscriptstyle\,\centerdot\,$}}}
}

%%% Path reversal %%%
\newcommand{\opp}[1]{\mathord{{#1}^{-1}}}
\let\rev\opp

%%% Coherence paths %%%
\newcommand{\ctassoc}{\mathsf{assoc}} % associativity law

%%% Transport (covariant) %%%
\newcommand{\trans}[2]{\ensuremath{{#1}_{*}\mathopen{}\left({#2}\right)\mathclose{}}\xspace}
\let\Trans\trans
%\newcommand{\Trans}[2]{\ensuremath{{#1}_{*}\left({#2}\right)}\xspace}
\newcommand{\transf}[1]{\ensuremath{{#1}_{*}}\xspace} % Without argument
%\newcommand{\transport}[2]{\ensuremath{\mathsf{transport}_{*} \: {#2}\xspace}}
\newcommand{\transfib}[3]{\ensuremath{\mathsf{transport}^{#1}(#2,#3)\xspace}}
\newcommand{\Transfib}[3]{\ensuremath{\mathsf{transport}^{#1}\Big(#2,\, #3\Big)\xspace}}
\newcommand{\transfibf}[1]{\ensuremath{\mathsf{transport}^{#1}\xspace}}

%%% 2D transport
\newcommand{\transtwo}[2]{\ensuremath{\mathsf{transport}^2\mathopen{}\left({#1},{#2}\right)\mathclose{}}\xspace}

%%% Constant transport
\newcommand{\transconst}[3]{\ensuremath{\mathsf{transportconst}}^{#1}_{#2}(#3)\xspace}
\newcommand{\transconstf}{\ensuremath{\mathsf{transportconst}}\xspace}

%%% Map on paths %%%
\newcommand{\mapfunc}[1]{\ensuremath{\mathsf{ap}_{#1}}\xspace} % Without argument
\newcommand{\map}[2]{\ensuremath{{#1}\mathopen{}\left({#2}\right)\mathclose{}}\xspace}
\let\Ap\map
%\newcommand{\Ap}[2]{\ensuremath{{#1}\left({#2}\right)}\xspace}
\newcommand{\mapdepfunc}[1]{\ensuremath{\mathsf{apd}_{#1}}\xspace} % Without argument
% \newcommand{\mapdep}[2]{\ensuremath{{#1}\llparenthesis{#2}\rrparenthesis}\xspace}
\newcommand{\mapdep}[2]{\ensuremath{\mapdepfunc{#1}\mathopen{}\left(#2\right)\mathclose{}}\xspace}
\let\apfunc\mapfunc
\let\ap\map
\let\apdfunc\mapdepfunc
\let\apd\mapdep

%%% 2D map on paths
\newcommand{\aptwofunc}[1]{\ensuremath{\mathsf{ap}^2_{#1}}\xspace}
\newcommand{\aptwo}[2]{\ensuremath{\aptwofunc{#1}\mathopen{}\left({#2}\right)\mathclose{}}\xspace}
\newcommand{\apdtwofunc}[1]{\ensuremath{\mathsf{apd}^2_{#1}}\xspace}
\newcommand{\apdtwo}[2]{\ensuremath{\apdtwofunc{#1}\mathopen{}\left(#2\right)\mathclose{}}\xspace}

%%% Identity functions %%%
\newcommand{\idfunc}[1][]{\ensuremath{\mathsf{id}_{#1}}\xspace}

%%% Homotopies (written infix) %%%
\newcommand{\htpy}{\sim}

%%% Other meanings of \sim
\newcommand{\bisim}{\sim}       % bisimulation
\newcommand{\eqr}{\sim}         % an equivalence relation

%%% Equivalence types %%%
\newcommand{\eqv}[2]{\ensuremath{#1 \simeq #2}\xspace}
\newcommand{\eqvspaced}[2]{\ensuremath{#1 \;\simeq\; #2}\xspace}
\newcommand{\eqvsym}{\simeq}    % infix symbol
\newcommand{\texteqv}[2]{\ensuremath{\mathsf{Equiv}(#1,#2)}\xspace}
\newcommand{\isequiv}{\ensuremath{\mathsf{isequiv}}}
\newcommand{\qinv}{\ensuremath{\mathsf{qinv}}}
\newcommand{\ishae}{\ensuremath{\mathsf{ishae}}}
\newcommand{\linv}{\ensuremath{\mathsf{linv}}}
\newcommand{\rinv}{\ensuremath{\mathsf{rinv}}}
\newcommand{\biinv}{\ensuremath{\mathsf{biinv}}}
\newcommand{\lcoh}[3]{\mathsf{lcoh}_{#1}(#2,#3)}
\newcommand{\rcoh}[3]{\mathsf{rcoh}_{#1}(#2,#3)}
\newcommand{\hfib}[2]{{\mathsf{fib}}_{#1}(#2)}

%%% Map on total spaces %%%
\newcommand{\total}[1]{\ensuremath{\mathsf{total}(#1)}}

%%% Universe types %%%
%\newcommand{\type}{\ensuremath{\mathsf{Type}}\xspace}
\newcommand{\UU}{\ensuremath{\mathcal{U}}\xspace}
\let\bbU\UU
\let\type\UU
% Universes of truncated types
\newcommand{\typele}[1]{\ensuremath{{#1}\text-\mathsf{Type}}\xspace}
\newcommand{\typeleU}[1]{\ensuremath{{#1}\text-\mathsf{Type}_\UU}\xspace}
\newcommand{\typelep}[1]{\ensuremath{{(#1)}\text-\mathsf{Type}}\xspace}
\newcommand{\typelepU}[1]{\ensuremath{{(#1)}\text-\mathsf{Type}_\UU}\xspace}
\let\ntype\typele
\let\ntypeU\typeleU
\let\ntypep\typelep
\let\ntypepU\typelepU
\renewcommand{\set}{\ensuremath{\mathsf{Set}}\xspace}
\newcommand{\setU}{\ensuremath{\mathsf{Set}_\UU}\xspace}
\newcommand{\prop}{\ensuremath{\mathsf{Prop}}\xspace}
\newcommand{\propU}{\ensuremath{\mathsf{Prop}_\UU}\xspace}
%Pointed types
\newcommand{\pointed}[1]{\ensuremath{#1_\bullet}}

%%% Ordinals and cardinals
\newcommand{\card}{\ensuremath{\mathsf{Card}}\xspace}
\newcommand{\ord}{\ensuremath{\mathsf{Ord}}\xspace}
\newcommand{\ordsl}[2]{{#1}_{/#2}}

%%% Univalence
\newcommand{\ua}{\ensuremath{\mathsf{ua}}\xspace} % the inverse of idtoeqv
\newcommand{\idtoeqv}{\ensuremath{\mathsf{idtoeqv}}\xspace}
\newcommand{\univalence}{\ensuremath{\mathsf{univalence}}\xspace} % the full axiom

%%% Truncation levels
\newcommand{\iscontr}{\ensuremath{\mathsf{isContr}}}
\newcommand{\contr}{\ensuremath{\mathsf{contr}}} % The path to the center of contraction
\newcommand{\isset}{\ensuremath{\mathsf{isSet}}}
\newcommand{\isprop}{\ensuremath{\mathsf{isProp}}}
% h-propositions
% \newcommand{\anhprop}{a mere proposition\xspace}
% \newcommand{\hprops}{mere propositions\xspace}

%%% Homotopy fibers %%%
%\newcommand{\hfiber}[2]{\ensuremath{\mathsf{hFiber}(#1,#2)}\xspace}
\let\hfiber\hfib

%%% Bracket/squash/truncation types %%%
% \newcommand{\brck}[1]{\textsf{mere}(#1)}
% \newcommand{\Brck}[1]{\textsf{mere}\Big(#1\Big)}
% \newcommand{\trunc}[2]{\tau_{#1}(#2)}
% \newcommand{\Trunc}[2]{\tau_{#1}\Big(#2\Big)}
% \newcommand{\truncf}[1]{\tau_{#1}}
%\newcommand{\trunc}[2]{\Vert #2\Vert_{#1}}
\newcommand{\trunc}[2]{\mathopen{}\left\Vert #2\right\Vert_{#1}\mathclose{}}
\newcommand{\ttrunc}[2]{\bigl\Vert #2\bigr\Vert_{#1}}
\newcommand{\Trunc}[2]{\Bigl\Vert #2\Bigr\Vert_{#1}}
\newcommand{\truncf}[1]{\Vert \blank \Vert_{#1}}
\newcommand{\tproj}[3][]{\mathopen{}\left|#3\right|_{#2}^{#1}\mathclose{}}
\newcommand{\tprojf}[2][]{|\blank|_{#2}^{#1}}
\def\pizero{\trunc0}
%\newcommand{\brck}[1]{\trunc{-1}{#1}}
%\newcommand{\Brck}[1]{\Trunc{-1}{#1}}
%\newcommand{\bproj}[1]{\tproj{-1}{#1}}
%\newcommand{\bprojf}{\tprojf{-1}}

\newcommand{\brck}[1]{\trunc{}{#1}}
\newcommand{\bbrck}[1]{\ttrunc{}{#1}}
\newcommand{\Brck}[1]{\Trunc{}{#1}}
\newcommand{\bproj}[1]{\tproj{}{#1}}
\newcommand{\bprojf}{\tprojf{}}

% Big parentheses
\newcommand{\Parens}[1]{\Bigl(#1\Bigr)}

% Projection and extension for truncations
\let\extendsmb\ext
\newcommand{\extend}[1]{\extendsmb(#1)}

%
%%% The empty type
\newcommand{\emptyt}{\ensuremath{\mathbf{0}}\xspace}

%%% The unit type
\newcommand{\unit}{\ensuremath{\mathbf{1}}\xspace}
\newcommand{\ttt}{\ensuremath{\star}\xspace}

%%% The two-element type
\newcommand{\bool}{\ensuremath{\mathbf{2}}\xspace}
\newcommand{\btrue}{{1_{\bool}}}
\newcommand{\bfalse}{{0_{\bool}}}

%%% Injections into binary sums and pushouts
\newcommand{\inlsym}{{\mathsf{inl}}}
\newcommand{\inrsym}{{\mathsf{inr}}}
\newcommand{\inl}{\ensuremath\inlsym\xspace}
\newcommand{\inr}{\ensuremath\inrsym\xspace}

%%% The segment of the interval
\newcommand{\seg}{\ensuremath{\mathsf{seg}}\xspace}

%%% Free groups
\newcommand{\freegroup}[1]{F(#1)}
\newcommand{\freegroupx}[1]{F'(#1)} % the "other" free group

%%% Glue of a pushout
\newcommand{\glue}{\mathsf{glue}}

%%% Colimits
\newcommand{\colim}{\mathsf{colim}}
\newcommand{\inc}{\mathsf{inc}}
\newcommand{\cmp}{\mathsf{cmp}}

%%% Circles and spheres
\newcommand{\Sn}{\mathbb{S}}
\newcommand{\base}{\ensuremath{\mathsf{base}}\xspace}
\newcommand{\lloop}{\ensuremath{\mathsf{loop}}\xspace}
\newcommand{\surf}{\ensuremath{\mathsf{surf}}\xspace}

%%% Suspension
\newcommand{\susp}{\Sigma}
\newcommand{\north}{\mathsf{N}}
\newcommand{\south}{\mathsf{S}}
\newcommand{\merid}{\mathsf{merid}}

%%% Blanks (shorthand for lambda abstractions)
\newcommand{\blank}{\mathord{\hspace{1pt}\text{--}\hspace{1pt}}}

%%% Nameless objects
\newcommand{\nameless}{\mathord{\hspace{1pt}\underline{\hspace{1ex}}\hspace{1pt}}}

%%% Some decorations
%\newcommand{\bbU}{\ensuremath{\mathbb{U}}\xspace}
% \newcommand{\bbB}{\ensuremath{\mathbb{B}}\xspace}
\newcommand{\bbP}{\ensuremath{\mathbb{P}}\xspace}

%%% Some categories
\newcommand{\uset}{\ensuremath{\mathcal{S}et}\xspace}
\newcommand{\ucat}{\ensuremath{{\mathcal{C}at}}\xspace}
\newcommand{\urel}{\ensuremath{\mathcal{R}el}\xspace}
\newcommand{\uhilb}{\ensuremath{\mathcal{H}ilb}\xspace}
\newcommand{\utype}{\ensuremath{\mathcal{T}\!ype}\xspace}

% Pullback corner
\newbox\pbbox
\setbox\pbbox=\hbox{\xy \POS(65,0)\ar@{-} (0,0) \ar@{-} (65,65)\endxy}
\def\pb{\save[]+<3.5mm,-3.5mm>*{\copy\pbbox} \restore}

% Macros for the categories chapter
\newcommand{\inv}[1]{{#1}^{-1}}
\newcommand{\idtoiso}{\ensuremath{\mathsf{idtoiso}}\xspace}
\newcommand{\isotoid}{\ensuremath{\mathsf{isotoid}}\xspace}
\newcommand{\op}{^{\mathrm{op}}}
\newcommand{\y}{\ensuremath{\mathbf{y}}\xspace}
\newcommand{\dgr}[1]{{#1}^{\dagger}}
\newcommand{\unitaryiso}{\mathrel{\cong^\dagger}}
\newcommand{\cteqv}[2]{\ensuremath{#1 \simeq #2}\xspace}
\newcommand{\cteqvsym}{\simeq}     % Symbol for equivalence of categories

%%% Natural numbers
\newcommand{\N}{\ensuremath{\mathbb{N}}\xspace}
%\newcommand{\N}{\textbf{N}}
\let\nat\N
\newcommand{\natp}{\ensuremath{\nat'}\xspace} % alternative nat in induction chapter

\newcommand{\zerop}{\ensuremath{0'}\xspace}   % alternative zero in induction chapter
\newcommand{\suc}{\mathsf{succ}}
\newcommand{\sucp}{\ensuremath{\suc'}\xspace} % alternative suc in induction chapter
\newcommand{\add}{\mathsf{add}}
\newcommand{\ack}{\mathsf{ack}}
\newcommand{\ite}{\mathsf{iter}}
\newcommand{\assoc}{\mathsf{assoc}}
\newcommand{\dbl}{\ensuremath{\mathsf{double}}}
\newcommand{\dblp}{\ensuremath{\dbl'}\xspace} % alternative double in induction chapter


%%% Lists
\newcommand{\lst}[1]{\mathsf{List}(#1)}
\newcommand{\nil}{\mathsf{nil}}
\newcommand{\cons}{\mathsf{cons}}
\newcommand{\lost}[1]{\mathsf{Lost}(#1)}

%%% Vectors of given length, used in induction chapter
\newcommand{\vect}[2]{\ensuremath{\mathsf{Vec}_{#1}(#2)}\xspace}

%%% Integers
\newcommand{\Z}{\ensuremath{\mathbb{Z}}\xspace}
\newcommand{\Zsuc}{\mathsf{succ}}
\newcommand{\Zpred}{\mathsf{pred}}

%%% Rationals
\newcommand{\Q}{\ensuremath{\mathbb{Q}}\xspace}

%%% Function extensionality
\newcommand{\funext}{\mathsf{funext}}
\newcommand{\happly}{\mathsf{happly}}

%%% A naturality lemma
\newcommand{\com}[3]{\mathsf{swap}_{#1,#2}(#3)}

%%% Code/encode/decode
\newcommand{\code}{\ensuremath{\mathsf{code}}\xspace}
\newcommand{\encode}{\ensuremath{\mathsf{encode}}\xspace}
\newcommand{\decode}{\ensuremath{\mathsf{decode}}\xspace}

% Function definition with domain and codomain
\newcommand{\function}[4]{\left\{\begin{array}{rcl}#1 &
      \longrightarrow & #2 \\ #3 & \longmapsto & #4 \end{array}\right.}

%%% Cones and cocones
\newcommand{\cone}[2]{\mathsf{cone}_{#1}(#2)}
\newcommand{\cocone}[2]{\mathsf{cocone}_{#1}(#2)}
% Apply a function to a cocone
\newcommand{\composecocone}[2]{#1\circ#2}
\newcommand{\composecone}[2]{#2\circ#1}
%%% Diagrams
\newcommand{\Ddiag}{\mathscr{D}}

%%% (pointed) mapping spaces
\newcommand{\Map}{\mathsf{Map}}

%%% The interval
\newcommand{\interval}{\ensuremath{I}\xspace}
\newcommand{\izero}{\ensuremath{0_{\interval}}\xspace}
\newcommand{\ione}{\ensuremath{1_{\interval}}\xspace}

%%% Arrows
\newcommand{\epi}{\ensuremath{\twoheadrightarrow}}
\newcommand{\mono}{\ensuremath{\rightarrowtail}}

%%% Sets
\newcommand{\bin}{\ensuremath{\mathrel{\widetilde{\in}}}}

%%% Semigroup structure
\newcommand{\semigroupstrsym}{\ensuremath{\mathsf{SemigroupStr}}}
\newcommand{\semigroupstr}[1]{\ensuremath{\mathsf{SemigroupStr}}(#1)}
\newcommand{\semigroup}[0]{\ensuremath{\mathsf{Semigroup}}}

%%% Macros for the formal type theory
\newcommand{\emptyctx}{\ensuremath{\cdot}}
\newcommand{\production}{\vcentcolon\vcentcolon=}
\newcommand{\conv}{\downarrow}
\newcommand{\ctx}{\ensuremath{\mathsf{ctx}}}
\newcommand{\wfctx}[1]{#1\ \ctx}
\newcommand{\oftp}[3]{#1 \vdash #2 : #3}
\newcommand{\jdeqtp}[4]{#1 \vdash #2 \jdeq #3 : #4}
\newcommand{\judg}[2]{#1 \vdash #2}
\newcommand{\tmtp}[2]{#1 \mathord{:} #2}

% rule names
\newcommand{\rform}{\textsc{form}}
\newcommand{\rintro}{\textsc{intro}}
\newcommand{\relim}{\textsc{elim}}
\newcommand{\rcomp}{\textsc{comp}}
\newcommand{\runiq}{\textsc{uniq}}
\newcommand{\Weak}{\mathsf{Wkg}}
\newcommand{\Vble}{\mathsf{Vble}}
\newcommand{\Exch}{\mathsf{Exch}}
\newcommand{\Subst}{\mathsf{Subst}}

%%% Macros for HITs
\newcommand{\cc}{\mathsf{c}}
\newcommand{\pp}{\mathsf{p}}
\newcommand{\cct}{\widetilde{\mathsf{c}}}
\newcommand{\ppt}{\widetilde{\mathsf{p}}}
\newcommand{\Wtil}{\ensuremath{\widetilde{W}}\xspace}

%%% Macros for n-types
\newcommand{\istype}[1]{\mathsf{is}\mbox{-}{#1}\mbox{-}\mathsf{type}}
\newcommand{\nplusone}{\ensuremath{(n+1)}}
\newcommand{\nminusone}{\ensuremath{(n-1)}}
\newcommand{\fact}{\mathsf{fact}}

%%% Macros for homotopy
\newcommand{\kbar}{\overline{k}} % Used in van Kampen's theorem

%%% Macros for induction
\newcommand{\natw}{\ensuremath{\mathbf{N^w}}\xspace}
\newcommand{\zerow}{\ensuremath{0^\mathbf{w}}\xspace}
\newcommand{\sucw}{\ensuremath{\mathsf{succ}^{\mathbf{w}}}\xspace}
\newcommand{\nalg}{\nat\mathsf{Alg}}
\newcommand{\nhom}{\nat\mathsf{Hom}}
\newcommand{\ishinitw}{\mathsf{isHinit}_{\mathsf{W}}}
\newcommand{\ishinitn}{\mathsf{isHinit}_\nat}
\newcommand{\w}{\mathsf{W}}
\newcommand{\walg}{\w\mathsf{Alg}}
\newcommand{\whom}{\w\mathsf{Hom}}

%%% Macros for real numbers
\newcommand{\RC}{\ensuremath{\mathbb{R}_\mathsf{c}}\xspace} % Cauchy
\newcommand{\RD}{\ensuremath{\mathbb{R}_\mathsf{d}}\xspace} % Dedekind
\newcommand{\R}{\ensuremath{\mathbb{R}}\xspace}           % Either
\newcommand{\barRD}{\ensuremath{\bar{\mathbb{R}}_\mathsf{d}}\xspace} % Dedekind completion of Dedekind

\newcommand{\close}[1]{\sim_{#1}} % Relation of closeness
\newcommand{\closesym}{\mathord\sim}
\newcommand{\rclim}{\mathsf{lim}} % HIT constructor for Cauchy reals
\newcommand{\rcrat}{\mathsf{rat}} % Embedding of rationals into Cauchy reals
\newcommand{\rceq}{\mathsf{eq}_{\RC}} % HIT path constructor
\newcommand{\CAP}{\mathcal{C}}    % The type of Cauchy approximations
\newcommand{\Qp}{\Q_{+}}
\newcommand{\apart}{\mathrel{\#}}  % apartness
\newcommand{\dcut}{\mathsf{isCut}}  % Dedekind cut
\newcommand{\cover}{\triangleleft} % inductive cover
\newcommand{\intfam}[3]{(#2, \lam{#1} #3)} % family of rational intervals

% Macros for the Cauchy reals construction
\newcommand{\bsim}{\frown}
\newcommand{\bbsim}{\smile}

\newcommand{\hapx}{\diamondsuit\approx}
\newcommand{\hapname}{\diamondsuit}
\newcommand{\hapxb}{\heartsuit\approx}
\newcommand{\hapbname}{\heartsuit}
\newcommand{\tap}[1]{\bullet\approx_{#1}\triangle}
\newcommand{\tapname}{\triangle}
\newcommand{\tapb}[1]{\bullet\approx_{#1}\square}
\newcommand{\tapbname}{\square}

%%% Macros for surreals
\newcommand{\NO}{\ensuremath{\mathsf{No}}\xspace}
\newcommand{\surr}[2]{\{\,#1\,\big|\,#2\,\}}
\newcommand{\LL}{\mathcal{L}}
\newcommand{\RR}{\mathcal{R}}
\newcommand{\noeq}{\mathsf{eq}_{\NO}} % HIT path constructor

\newcommand{\ble}{\trianglelefteqslant}
\newcommand{\blt}{\vartriangleleft}
\newcommand{\bble}{\sqsubseteq}
\newcommand{\bblt}{\sqsubset}

\newcommand{\hle}{\diamondsuit\preceq}
\newcommand{\hlt}{\diamondsuit\prec}
\newcommand{\hlname}{\diamondsuit}
\newcommand{\hleb}{\heartsuit\preceq}
\newcommand{\hltb}{\heartsuit\prec}
\newcommand{\hlbname}{\heartsuit}
% \newcommand{\tle}{(\bullet\preceq\triangle)}
% \newcommand{\tlt}{(\bullet\prec\triangle)}
\newcommand{\tle}{\triangle\preceq}
\newcommand{\tlt}{\triangle\prec}
\newcommand{\tlname}{\triangle}
% \newcommand{\tleb}{(\bullet\preceq\square)}
% \newcommand{\tltb}{(\bullet\prec\square)}
\newcommand{\tleb}{\square\preceq}
\newcommand{\tltb}{\square\prec}
\newcommand{\tlbname}{\square}

%%% Macros for set theory
\newcommand{\vset}{\mathsf{set}}  % point constructor for cummulative hierarchy V
\def\cd{\tproj0}
\newcommand{\inj}{\ensuremath{\mathsf{inj}}} % type of injections
\newcommand{\acc}{\ensuremath{\mathsf{acc}}} % accessibility

\newcommand{\atMostOne}{\mathsf{atMostOne}}

\newcommand{\power}[1]{\mathcal{P}(#1)} % power set
\newcommand{\powerp}[1]{\mathcal{P}_+(#1)} % inhabited power set

%%%% THEOREM ENVIRONMENTS %%%%

% The cleveref package provides \cref{...} which is like \ref{...}
% except that it automatically inserts the type of the thing you're
% referring to, e.g. it produces "Theorem 3.8" instead of just "3.8"
% (and hyperref makes the whole thing a hyperlink).  This saves a slight amount
% of typing, but more importantly it means that if you decide later on
% that 3.8 should be a Lemma or a Definition instead of a Theorem, you
% don't have to change the name in all the places you referred to it.

% The following hack improves on this by using the same counter for
% all theorem-type environments, so that after Theorem 1.1 comes
% Corollary 1.2 rather than Corollary 1.1.  This makes it much easier
% for the reader to find a particular theorem when flipping through
% the document.
\makeatletter
\def\defthm#1#2#3{%
  %% Ensure all theorem types are numbered with the same counter
  \newaliascnt{#1}{thm}
  \newtheorem{#1}[#1]{#2}
  \aliascntresetthe{#1}
  %% This command tells cleveref's \cref what to call things
  \crefname{#1}{#2}{#3}% following brace must be on separate line to support poorman cleveref sed file
}

% Now define a bunch of theorem-type environments.
\newtheorem{thm}{Theorem}[section]
\crefname{thm}{Theorem}{Theorems}
%\defthm{prop}{Proposition}   % Probably we shouldn't use "Proposition" in this way
\defthm{cor}{Corollary}{Corollaries}
\defthm{lem}{Lemma}{Lemmas}
\defthm{axiom}{Axiom}{Axioms}
% Since definitions and theorems in type theory are synonymous, should
% we actually use the same theoremstyle for them?
\theoremstyle{definition}
\defthm{defn}{Definition}{Definitions}
\theoremstyle{remark}
\defthm{rmk}{Remark}{Remarks}
\defthm{eg}{Example}{Examples}
\defthm{egs}{Examples}{Examples}
\defthm{notes}{Notes}{Notes}
% Number exercises within chapters, with their own counter.
\newtheorem{ex}{Exercise}[chapter]
\crefname{ex}{Exercise}{Exercises}

% Display format for sections
\crefformat{section}{\S#2#1#3}
\Crefformat{section}{Section~#2#1#3}
\crefrangeformat{section}{\S\S#3#1#4--#5#2#6}
\Crefrangeformat{section}{Sections~#3#1#4--#5#2#6}
\crefmultiformat{section}{\S\S#2#1#3}{ and~#2#1#3}{, #2#1#3}{ and~#2#1#3}
\Crefmultiformat{section}{Sections~#2#1#3}{ and~#2#1#3}{, #2#1#3}{ and~#2#1#3}
\crefrangemultiformat{section}{\S\S#3#1#4--#5#2#6}{ and~#3#1#4--#5#2#6}{, #3#1#4--#5#2#6}{ and~#3#1#4--#5#2#6}
\Crefrangemultiformat{section}{Sections~#3#1#4--#5#2#6}{ and~#3#1#4--#5#2#6}{, #3#1#4--#5#2#6}{ and~#3#1#4--#5#2#6}

% Display format for appendices
\crefformat{appendix}{Appendix~#2#1#3}
\Crefformat{appendix}{Appendix~#2#1#3}
\crefrangeformat{appendix}{Appendices~#3#1#4--#5#2#6}
\Crefrangeformat{appendix}{Appendices~#3#1#4--#5#2#6}
\crefmultiformat{appendix}{Appendices~#2#1#3}{ and~#2#1#3}{, #2#1#3}{ and~#2#1#3}
\Crefmultiformat{appendix}{Appendices~#2#1#3}{ and~#2#1#3}{, #2#1#3}{ and~#2#1#3}
\crefrangemultiformat{appendix}{Appendices~#3#1#4--#5#2#6}{ and~#3#1#4--#5#2#6}{, #3#1#4--#5#2#6}{ and~#3#1#4--#5#2#6}
\Crefrangemultiformat{appendix}{Appendices~#3#1#4--#5#2#6}{ and~#3#1#4--#5#2#6}{, #3#1#4--#5#2#6}{ and~#3#1#4--#5#2#6}

\crefname{part}{Part}{Parts}

% Number subsubsections
\setcounter{secnumdepth}{5}

% Display format for figures
\crefname{figure}{Figure}{Figures}

%%%% EQUATION NUMBERING %%%%

% The following hack uses the single theorem counter to number
% equations as well, so that we don't have both Theorem 1.1 and
% equation (1.1).
\let\c@equation\c@thm
\numberwithin{equation}{section}


%%%% ENUMERATE NUMBERING %%%%

% Number the first level of enumerates as (i), (ii), ...
\renewcommand{\theenumi}{(\roman{enumi})}
\renewcommand{\labelenumi}{\theenumi}


%%%% MARGINS %%%%

% This is a matter of personal preference, but I think the left
% margins on enumerates and itemizes are too wide.
\setitemize[1]{leftmargin=2em}
\setenumerate[1]{leftmargin=*}

% Likewise that they are too spaced out.
\setitemize[1]{itemsep=-0.2em}
\setenumerate[1]{itemsep=-0.2em}

%%% Notes %%%
\def\noteson{%
\gdef\note##1{\mbox{}\marginpar{\color{blue}\textasteriskcentered\ ##1}}}
\gdef\notesoff{\gdef\note##1{\null}}
\noteson

\newcommand{\Coq}{\textsc{Coq}\xspace}
\newcommand{\Agda}{\textsc{Agda}\xspace}
\newcommand{\NuPRL}{\textsc{NuPRL}\xspace}

%%%% CITATIONS %%%%

% \let \cite \citep

%%%% INDEX %%%%

\newcommand{\footstyle}[1]{{\hyperpage{#1}}n} % If you index something that is in a footnote
\newcommand{\defstyle}[1]{\textbf{\hyperpage{#1}}}  % Style for pageref to a definition

\newcommand{\indexdef}[1]{\index{#1|defstyle}}   % Index a definition
\newcommand{\indexfoot}[1]{\index{#1|footstyle}} % Index a term in a footnote
\newcommand{\indexsee}[2]{\index{#1|see{#2}}}    % Index "see also"


%%%% Standard phrasing or spelling of common phrases %%%%

\newcommand{\ZF}{Zermelo--Fraenkel}
\newcommand{\CZF}{Constructive \ZF{} Set Theory}

\newcommand{\LEM}[1]{\ensuremath{\mathsf{LEM}_{#1}}\xspace}
\newcommand{\choice}[1]{\ensuremath{\mathsf{AC}_{#1}}\xspace}

%%%% MISC %%%%

\newcommand{\mentalpause}{\medskip} % Use for "mental" pause, instead of \smallskip or \medskip

%% Use \symlabel instead of \label to mark a pageref that you need in the index of symbols
\newcounter{symindex}
\newcommand{\symlabel}[1]{\refstepcounter{symindex}\label{#1}}

% Local Variables:
% mode: latex
% TeX-master: "hott-online"
% End:


%% TITLE
\title{Runners in action}

%% AUTHORS
\author{Danel Ahman \and Andrej Bauer}
\institute{%
  Faculty of Mathematics and Physics\\
  University of Ljubljana, Slovenia}

\maketitle

%% ABSTRACT
\begin{abstract}

Runners of algebraic effects, also known as comodels, provide a mathematical
model of resource management. We show that they also give rise to a programming
concept that models top-level external resources, as well as allows programmers
to modularly define their own intermediate ``virtual machines''.
We capture the core ideas of programming with runners in an equational calculus $\lambdacoop$, which we equip  
with a sound and coherent denotational semantics that guarantees the linear use of
resources and execution of finalisation code. We accompany $\lambdacoop$ with
examples of runners in action, provide a prototype language implementation in \pl{OCaml},
as well as a \pl{Haskell} library based on $\lambdacoop$.
\keywords{Runners, comodels, algebraic effects, resources, finalisation.}
\end{abstract}

%%%%%%%%%%%%%%%%%%%%%%%%%%%%%%%%%%%%%%%%%%%%%%%%%%%%

%% CODE SNIPPETS TYPESETTING

\definecolor{codegreen}{rgb}{0,0.6,0}
\definecolor{codegray}{rgb}{0.5,0.5,0.5}
\definecolor{codepurple}{rgb}{0.58,0,0.82}
\definecolor{backcolour}{rgb}{0.95,0.95,0.92}

\colorlet{keywordColor}{NavyBlue} % the color of language keywords
\colorlet{rulenameColor}{Gray} % the color of rule names

\def\lstlanguagefiles{coop.tex}
\lstset{language=coop,upquote=true}
\let\ls\lstinline

%%%%%%%%%%%%%%%%%%%%%%%%%%%%%%%%%%%%%%%%%%%%%%%%%%%%

%% PAPER CONTENTS

% Customize the display of references to sections, subsections, subsubsections, theorems, and propositions.
\crefformat{section}{\S#2#1#3}
\Crefformat{section}{\S#2#1#3}

\crefformat{subsection}{\S#2#1#3}
\Crefformat{subsection}{\S#2#1#3}

\crefformat{subsubsection}{\S#2#1#3}
\Crefformat{subsubsection}{\S#2#1#3}

\crefformat{theorem}{Thm.~#2#1#3}
\Crefformat{theorem}{Thm.~#2#1#3}

\crefformat{proposition}{Prop.~#2#1#3}
\Crefformat{proposition}{Prop.~#2#1#3}

\crefformat{figure}{Fig.~#2#1#3}
\Crefformat{figure}{Fig.~#2#1#3}

\chapter*{Introduction}
\markboth{\textsc{Introduction}}{}
\addcontentsline{toc}{chapter}{Introduction}
\setcounter{page}{1}
\pagenumbering{arabic}


\emph{Homotopy type theory} is a new branch of mathematics that combines aspects of several different fields in a surprising way. It is based on a recently discovered connection between \emph{homotopy theory} and \emph{type theory}.
Homotopy theory is an outgrowth of algebraic topology and homological algebra, with relationships to higher category theory; while type theory is a branch of mathematical logic and theoretical computer science.
Although the connections between the two are currently the focus of intense investigation, it is increasingly clear that they are just the beginning of a subject that will take more time and more hard work to fully understand.
It touches on topics as seemingly distant as the homotopy groups of spheres, the algorithms for type checking, and the definition of weak $\infty$-groupoids.

Homotopy type theory also brings new ideas into the very foundation of mathematics.
\index{foundations, univalent}%
On the one hand, there is Voevodsky's subtle and beautiful \emph{univalence axiom}. 
\index{univalence axiom}%
The univalence axiom implies, in particular, that isomorphic structures can be identified, a principle that mathematicians have been happily using on workdays, despite its incompatibility with the ``official'' doctrines of conventional foundations.
On the other hand, we have \emph{higher inductive types}, which provide direct, logical descriptions of some of the basic spaces and constructions of homotopy theory: spheres, cylinders, truncations, localizations, etc.
Both ideas are impossible to capture directly in classical set-theoretic foundations, but when combined in homotopy type theory, they permit an entirely new kind of ``logic of homotopy types''.
\index{foundations}%

This suggests a new conception of foundations of mathematics, with intrinsic homotopical content, an ``invariant'' conception of the objects of mathematics --- and convenient machine implementations, which can serve as a practical aid to the working mathematician.
This is the \emph{Univalent Foundations} program.
The present book is intended as a first systematic exposition of the basics of univalent foundations, and a collection of examples of this new style of reasoning --- but without requiring the reader to know or learn any formal logic, or to use any computer proof assistant.

% This enlarges the page by one line in letter format. Use sparringly.
\OPTwidow

We emphasize that homotopy type theory is a young field, and univalent foundations is very much a work in progress. 
This book should be regarded as a ``snapshot'' of just one portion of the field, taken at the time it was written, rather than a polished exposition of a completed edifice. 
As we will discuss briefly later, there are many aspects of homotopy type theory that are not yet fully understood --- and some that are not even touched upon here. 
The ultimate theory will almost certainly not look exactly like the one described in this book, but it will surely be \emph{at least} as capable and powerful; we therefore believe that univalent foundations will eventually become a viable alternative to set theory as the ``implicit foundation'' for the unformalized mathematics done by most mathematicians.

\subsection*{Type theory}

Type theory was originally invented by Bertrand Russell \cite{Russell:1908},\index{Russell, Bertrand} as a device for blocking the paradoxes in the logical foundations of mathematics  that were under investigation at the time.
It was developed further by many people over the next few decades, particularly Church~\cite{Church:1940tu,Church:1941tc} who combined it with his \textit{$\lambda$-calculus}.
Although it is not generally regarded as the foundation for classical mathematics, set theory being more customary, type theory still has numerous applications, especially in computer science and the theory of programming languages~\cite{Pierce-TAPL}.
\index{programming}%
\index{type theory}%
\index{lambda-calculus@$\lambda$-calculus}%
Per Martin-L\"{o}f \cite{Martin-Lof-1972,Martin-Lof-1973,Martin-Lof-1979,martin-lof:bibliopolis}, among others,
developed a ``predicative'' modification of Church's type system, which is now usually called dependent, constructive, intuitionistic, or simply \emph{Martin\--L\"of type theory}. This is the basis of the system that we consider here; it was originally intended as a rigorous framework for the formalization of constructive mathematics.  In what follows, we will often use ``type theory'' to refer specifically to this system and similar ones, although type theory as a subject is much broader (see \cite{somma,kamar} for the history of type theory).

In type theory, unlike set theory, objects are classified using a primitive notion of \emph{type}, similar to the data-types used in programming languages.  These elaborately structured types can be used to express detailed specifications of the objects classified, giving rise to principles of reasoning about these objects.  To take a very simple example, the objects of a product type $A\times B$ are known to be of the form $\pairr{a,b}$, and so one automatically knows how to construct them and how to decompose them. Similarly, an object of function type $A\to B$ can be acquired from an object of type $B$ parametrized by objects of type $A$, and can be evaluated at an argument of type $A$.  This rigidly predictable behavior of all objects (as opposed to set theory's more liberal formation principles, allowing inhomogeneous sets) is one aspect of type theory that has led to its extensive use in verifying the correctness of computer programs.  The clear reasoning principles associated with the construction of types also form the basis of modern \emph{computer proof assistants},%
\index{proof!assistant}%
\indexsee{computer proof assistant}{proof assistant}
\index{mathematics!formalized}%
which are used for formalizing mathematics and verifying the correctness of formalized proofs.  We return to this aspect of type theory below.  

One problem in understanding type theory from a mathematical point of view, however, has always been that the basic concept of \emph{type} is unlike that of \emph{set} in ways that have been hard to make precise.  We believe that the new idea of regarding types, not as strange sets (perhaps constructed without using classical logic), but as spaces, viewed from the perspective of homotopy theory, is a significant step forward.  In particular, it solves the problem of understanding how the notion of equality of elements of a type differs from that of elements of a set.

In homotopy theory one is concerned with spaces
\index{topological!space}%
and continuous mappings between them, 
\index{function!continuous!in classical homotopy theory}%
up to homotopy.  A \emph{homotopy}
\index{homotopy!topological}%
between a pair of continuous maps $f : X \to Y$
and  $g : X\to Y$ is 
a continuous map $H : X \times [0, 1] \to Y$ satisfying
$H(x, 0) = f (x)$  and $H(x, 1) = g(x)$. The homotopy $H$ may be thought of as a ``continuous deformation'' of $f$ into $g$. The spaces $X$ and $Y$ are said to be \emph{homotopy equivalent},
\index{homotopy!equivalence!topological}%
$\eqv X Y$, if there are continuous maps going back and forth, the composites of which are homotopical to the respective identity mappings, i.e., if they are isomorphic ``up to homotopy''.  Homotopy equivalent spaces have the same algebraic invariants (e.g., homology, or the fundamental group), and are said to have the same \emph{homotopy type}.

\subsection*{Homotopy type theory}

Homotopy type theory (HoTT) interprets type theory from a homotopical perspective.
In homotopy type theory, we regard the types as ``spaces'' (as studied in homotopy theory) or higher groupoids, and the logical constructions (such as the product $A\times B$) as homotopy-invariant constructions on these spaces.
In this way, we are able to manipulate spaces directly without first having to develop point-set topology (or any combinatorial replacement for it, such as the theory of simplicial sets).
To briefly explain this perspective, consider first the basic concept of type theory, namely that
the \emph{term} $a$ is of \emph{type} $A$, which is written:
\[ a:A. \]
This expression is traditionally thought of as akin to:
\begin{center}
``$a$ is an element of the set $A$''.
\end{center}
However, in homotopy type theory we think of it instead as:
\begin{center}
``$a$ is a point of the space $A$''.
\end{center}
\index{continuity of functions in type theory@``continuity'' of functions in type theory}%
Similarly, every function $f : A\to B$ in type theory is regarded as a continuous map from the space $A$ to the space $B$.

We should stress that these ``spaces'' are treated purely homotopically, not topologically.
For instance, there is no notion of ``open subset'' of a type or of ``convergence'' of a sequence of elements of a type.
We only have ``homotopical'' notions, such as paths between points and homotopies between paths, which also make sense in other models of homotopy theory (such as simplicial sets).
Thus, it would be more accurate to say that we treat types as \emph{$\infty$-groupoids}\index{.infinity-groupoid@$\infty$-groupoid}; this is a name for the ``invariant objects'' of homotopy theory which can be presented by topological spaces,
\index{topological!space}%
simplicial sets, or any other model for homotopy theory.
However, it is convenient to sometimes use topological words such as ``space'' and ``path'', as long as we remember that other topological concepts are not applicable.

(It is tempting to also use the phrase \emph{homotopy type}
\index{homotopy!type}%
for these objects, suggesting the dual interpretation of ``a type (as in type theory) viewed homotopically'' and ``a space considered from the point of view of homotopy theory''.
The latter is a bit different from the classical meaning of ``homotopy type'' as an \emph{equivalence class} of spaces modulo homotopy equivalence, although it does preserve the meaning of phrases such as ``these two spaces have the same homotopy type''.)

The idea of interpreting types as structured objects, rather than sets, has a long pedigree, and is known to clarify various mysterious aspects of type theory.
For instance, interpreting types as sheaves helps explain the intuitionistic nature of type-theoretic logic, while interpreting them as partial equivalence relations or ``domains'' helps explain its computational aspects.
It also implies that we can use type-theoretic reasoning to study the structured objects, leading to the rich field of categorical logic.
The homotopical interpretation fits this same pattern: it clarifies the nature of \emph{identity} (or equality) in type theory, and allows us to use type-theoretic reasoning in the study of homotopy theory.

The key new idea of the homotopy interpretation is that the logical notion of identity $a = b$ of two objects $a, b: A$ of the same type $A$ can be understood as the existence of a path $p : a \leadsto b$ from point $a$ to point $b$ in the space $A$.
This also means that two functions $f, g: A\to B$ can be identified if they are homotopic, since a homotopy is just a (continuous) family of paths $p_x: f(x) \leadsto g(x)$ in $B$, one for each $x:A$.
In type theory, for every type $A$ there is a (formerly somewhat mysterious) type $\idtypevar{A}$ of identifications of two objects of $A$; in homotopy type theory, this is just the \emph{path space} $A^I$ of all continuous maps $I\to A$ from the unit interval.
\index{unit!interval}%
\index{interval!topological unit}%
\index{path!topological}%
\index{topological!path}%
In this way, a term $p : \idtype[A]{a}{b}$ represents a path $p : a \leadsto b$ in $A$. 

The idea of homotopy type theory arose around 2006 in independent work by Awodey and Warren~\cite{AW} and Voevodsky~\cite{VV}, but it was inspired by 
Hofmann and Streicher's earlier groupoid interpretation~\cite{hs:gpd-typethy}.
Indeed, higher-dimensional category theory (particularly the theory of weak $\infty$-groupoids) is now known to be intimately connected to homotopy theory, as proposed by Grothendieck and now being studied intensely by mathematicians of both sorts.
The original semantic models of Awodey--Warren and Voevodsky use well-known notions and techniques from homotopy theory which are now also in use in higher category theory, such as  Quillen model categories and Kan\index{Kan complex} simplicial sets\index{simplicial!sets}.
\index{Quillen model category}%
\index{model category}%

In particular, Voevodsky constructed an interpretation of type theory in Kan simplicial sets, and recognized that this interpretation satisfied a further crucial property which he dubbed \emph{univalence}.
This had not previously been considered in type theory (although Church's principle of extensionality for propositions turns out to be a very special case of it, and Hofmann and Streicher had considered another special case under the name ``universe extensionality'').
Adding univalence to type theory in the form of a new axiom has far-reaching consequences, many of which are natural, simplifying and compelling.
The univalence axiom also further strengthens the homotopical view of type theory, since it holds in the simplicial model and other related models, while failing under the view of types as sets.

\subsection*{Univalent foundations}

Very briefly, the basic idea of the univalence axiom can be explained as follows.
In type theory, one can have a universe $\UU$, the terms of which are themselves types, $A : \UU$, etc.
Those types that are terms of $\UU$ are commonly called \emph{small} types.
\index{type!small}%
\index{small!type}%
Like any type, $\UU$ has an identity type $\idtypevar{\UU}$, which expresses the identity relation $A = B$ between small types.
Thinking of types as spaces, $\UU$ is a space, the points of which are spaces; to understand its identity type, we must ask, what is a path $p : A \leadsto B$ between spaces in $\UU$?
The univalence axiom says that such paths correspond to homotopy equivalences $\eqv A B$, (roughly) as explained above.
A bit more precisely, given any (small) types $A$ and $B$, in addition to the primitive type $\idtype[\UU]AB$ of identifications of $A$ with $B$, there is the defined type $\texteqv AB$ of equivalences from $A$ to $B$.
Since the identity map on any object is an equivalence, there is a canonical map,
\[\idtype[\UU]AB\to\texteqv AB.\]
The univalence axiom states that this map is itself an equivalence.
At the risk of oversimplifying, we can state this succinctly as follows:

\begin{description}\index{univalence axiom}%
\item[Univalence Axiom:]  $\eqvspaced{(A = B)}{(\eqv A B)}$.
\end{description}
%
In other words, identity is equivalent to equivalence. \index{identity}% 
In particular, one may say that ``equivalent types are identical''.
However, this phrase is somewhat misleading, since it may sound like a sort of ``skeletality'' condition which \emph{collapses} the notion of equivalence to coincide with identity, whereas in fact univalence is about \emph{expanding} the notion of identity so as to coincide with the (unchanged) notion of equivalence.

From the homotopical point of view, univalence implies that spaces of the same homotopy type are connected by a path in the universe $\UU$, in accord with the intuition of a classifying space for (small) spaces.
From the logical point of view, however, it is a radically new idea: it says that isomorphic things can be identified!  Mathematicians are of course used to identifying isomorphic structures in practice, but they generally do so by ``abuse of notation''\index{abuse!of notation}, or some other informal device, knowing that the objects involved are not ``really'' identical.  But in this new foundational scheme, such structures can be formally identified, in the logical sense that every property or construction involving one also applies to the other. Indeed, the identification is now made explicit, and properties and constructions can be systematically transported along it.  Moreover, the different ways in which such identifications may be made themselves form a structure that one can (and should!)\ take into account.

Thus in sum, for points $A$ and $B$ of the universe $\UU$ (i.e., small types), the univalence axiom identifies the following three notions:
\begin{itemize}
\item (logical) an identification $p:A=B$ of $A$ and $B$
\item (topological) a path $p:A \leadsto B$ from $A$ to $B$ in $\UU$
\item (homotopical) an equivalence $p:\eqv A B$ between $A$ and $B$.
\end{itemize}

\subsection*{Higher inductive types}\index{type!higher inductive}%

One of the classical advantages of type theory is its simple and effective techniques for working with inductively defined structures.
The simplest nontrivial inductively defined structure is the natural numbers, which is inductively generated by zero and the successor function.
From this statement one can algorithmically\index{algorithm} extract the principle of mathematical induction, which characterizes the natural numbers.
More general inductive definitions encompass lists and well-founded trees of all sorts, each of which is characterized by a corresponding ``induction principle''.
This includes most data structures used in certain programming languages; hence the usefulness of type theory in formal reasoning about the latter.
If conceived in a very general sense, inductive definitions also include examples such as a disjoint union $A+B$, which may be regarded as ``inductively'' generated by the two injections $A\to A+B$ and $B\to A+B$.
The ``induction principle'' in this case is ``proof by case analysis'', which characterizes the disjoint union.

In homotopy theory, it is natural to consider also ``inductively defined spaces'' which are generated not merely by a collection of \emph{points}, but also by collections of \emph{paths} and higher paths.
Classically, such spaces are called \emph{CW complexes}.
\index{CW complex}%
For instance, the circle $S^1$ is generated by a single point and a single path from that point to itself.
Similarly, the 2-sphere $S^2$ is generated by a single point $b$ and a single two-dimensional path from the constant path at $b$ to itself, while the torus $T^2$ is generated by a single point, two paths $p$ and $q$ from that point to itself, and a two-dimensional path from $p\ct q$ to $q\ct p$.

By using the identification of paths with identities in homotopy type theory, these sort of ``inductively defined spaces'' can be characterized in type theory by ``induction principles'', entirely analogously to classical examples such as the natural numbers and the disjoint union.
The resulting \emph{higher inductive types}
\index{type!higher inductive}%
give a direct ``logical'' way to reason about familiar spaces such as spheres, which (in combination with univalence) can be used to perform familiar arguments from homotopy theory, such as calculating homotopy groups of spheres, in a purely formal way.
The resulting proofs are a marriage of classical homotopy-theoretic ideas with classical type-theoretic ones, yielding new insight into both disciplines.

Moreover, this is only the tip of the iceberg: many abstract constructions from homotopy theory, such as homotopy colimits, suspensions, Postnikov towers, localization, completion, and spectrification, can also be expressed as higher inductive types.
Many of these are classically constructed using Quillen's ``small object argument'', which can be regarded as a finite way of algorithmically describing an infinite CW complex presentation\index{presentation!of a space as a CW complex} of a space, just as ``zero and successor'' is a finite algorithmic\index{algorithm} description of the infinite set of natural numbers.
Spaces produced by the small object argument are infamously complicated and difficult to understand; the type-theoretic approach is potentially much simpler, bypassing the need for any explicit construction by giving direct access to the appropriate ``induction principle''.
Thus, the combination of univalence and higher inductive types suggests the possibility of a revolution, of sorts, in the practice of homotopy theory.


\subsection*{Sets in univalent foundations}

\index{set|(}%

We have claimed that univalent foundations can eventually serve as a foundation for ``all'' of mathematics, but so far we have discussed 
only homotopy theory.  Of course, there are many specific examples of the use of type theory without the new homotopy type theory features to formalize mathematics,
\index{mathematics!formalized}%
\index{theorem!Feit--Thompson}%
\index{theorem!odd-order}%
\index{Feit--Thompson theorem}%
\index{odd-order theorem}%
such as the recent formalization of the Feit--Thompson odd-order theorem in \Coq~\cite{gonthier}.

But the traditional view is that mathematics is founded on set theory, in the sense that all mathematical objects and constructions can be coded into a theory such as Zermelo--Fraenkel set theory (ZF).
\index{set theory!Zermelo--Fraenkel}%
\indexsee{Zermelo-Fraenkel set theory}{set theory}%
\indexsee{ZF}{set theory}%
\indexsee{ZFC}{set theory}%
However, it is well-established by now that for most mathematics outside of set theory proper, the intricate hierarchical membership structure of sets in ZF is really unnecessary: a more ``structural'' theory, such as Lawvere's\index{Lawvere} Elementary Theory of the Category of Sets~\cite{lawvere:etcs-long}, suffices.
\index{Elementary Theory of the Category of Sets}%

In univalent foundations, the basic objects are ``homotopy types'' rather than sets, but we can \emph{define} a class of types which behave like sets.
Homotopically, these can be thought of as spaces in which every connected component is contractible, i.e.\ those which are homotopy equivalent to a discrete space.
\index{discrete!space}%
It is a theorem  that the category of such ``sets'' satisfies Lawvere's\index{Lawvere} axioms (or related ones, depending on the details of the theory).
Thus, any sort of mathematics that can be represented in an ETCS-like theory (which, experience suggests, is essentially all of mathematics) can equally well be represented in univalent foundations.  

This supports the claim that univalent foundations is at least as good as existing foundations of mathematics.
A mathematician working in univalent foundations can build structures out of sets in a familiar way, with more general homotopy types waiting in the foundational background until there is need of them.
For this reason, most of the applications in this book have been chosen to be areas where univalent foundations has something \emph{new} to contribute that distinguishes it from existing foundational systems.

Unsurprisingly, homotopy theory and category theory are two of these, but perhaps less obvious is that univalent foundations has something new and interesting to offer even in subjects such as set theory and real analysis.
For instance, the univalence axiom allows us to identify isomorphic structures, while higher inductive types allow direct descriptions of objects by their universal properties.
Thus we can generally avoid resorting to arbitrarily chosen representatives or transfinite iterative constructions.
In fact, even the objects of study in ZF set theory can be characterized, inside the sets of univalent foundations, by such an inductive universal property.

\index{set|)}%


\subsection*{Informal type theory}

\index{mathematics!formalized|(defstyle}%
\index{informal type theory|(defstyle}%
\index{type theory!informal|(defstyle}%
\index{type theory!formal|(}%
One difficulty often encountered by the classical mathematician when faced with learning about type theory is that it is usually presented as a fully or partially formalized deductive system.
This style, which is very useful for proof-theoretic investigations, is not particularly convenient for use in applied, informal reasoning.
Nor is it even familiar to most working mathematicians, even those who might be interested in foundations of mathematics.
One objective of the present work is to develop an informal style of doing mathematics in univalent foundations that is at once rigorous and precise, but is also closer to the language and style of presentation of everyday mathematics.

In present-day mathematics, one usually constructs and reasons about mathematical objects in a way that could in principle, one presumes, be formalized in a system of elementary set theory, such as ZFC --- at least given enough ingenuity and patience.
For the most part, one does not even need to be aware of this possibility, since it largely coincides with the condition that a proof be ``fully rigorous'' (in the sense that all mathematicians have come to understand intuitively through education and experience).
But one does need to learn to be careful about a few aspects of ``informal set theory'': the use of collections too large or inchoate to be sets; the axiom of choice and its equivalents; even (for undergraduates) the method of proof by contradiction; and so on.
Adopting a new foundational system such as homotopy type theory as the \emph{implicit formal basis} of informal reasoning will require adjusting some of one's instincts and practices.
The present text is intended to serve as an example of this ``new kind of mathematics'', which is still informal, but could now in principle be formalized in homotopy type theory, rather than ZFC, again given enough ingenuity and patience.

It is worth emphasizing that, in this new system, such formalization can have real practical benefits.
The formal system of type theory is suited to computer systems and has been implemented in existing proof assistants.
\index{proof!assistant}%
A proof assistant is a computer program which guides the user in construction of a fully formal proof, only allowing valid steps of reasoning.
It also provides some degree of automation, can search libraries for existing theorems, and can even extract numerical algorithms\index{algorithm} \index{extraction of algorithms} from the resulting (constructive) proofs.

We believe that this aspect of the univalent foundations program distinguishes it from other approaches to foundations, potentially providing a new practical utility for the working mathematician.
Indeed, proof assistants based on older type theories have already been used to formalize substantial mathematical proofs, such as the four-color theorem\index{theorem!four-color} \index{four-color theorem} and the Feit--Thompson theorem.
Computer implementations of univalent foundations are presently works in progress (like the theory itself).
\index{proof!assistant}%
However, even its currently available implementations (which are mostly small modifications to existing proof assistants such as \Coq and 
\Agda) have already demonstrated their worth, not only in the formalization of known proofs, but in the discovery of new ones.
Indeed, many of the proofs described in this book were actually \emph{first} done in a fully formalized form in a proof assistant, and are only now being ``unformalized'' for the first time --- a reversal of the usual relation between formal and informal mathematics.

One can imagine a not-too-distant future when it will be possible for mathematicians to verify the correctness of their own papers by working within the system of univalent foundations, formalized in a proof assistant, and that doing so will become as natural as typesetting their own papers in \TeX.
%(Whether this proves to be the publishers' dream or their nightmare remains to be seen.) 
In principle, this could be equally true for any other foundational system, but we believe it to be more practically attainable using univalent foundations, as witnessed by the present work and its formal counterpart.

\index{type theory!formal|)}%
\index{informal type theory|)}%
\index{type theory!informal|)}%
\index{mathematics!formalized|)}%

\subsection*{Constructivity} 

\index{mathematics!constructive|(}%

One of the most striking differences between classical\index{mathematics!classical} foundations and type theory is the idea of \emph{proof relevance}, according to which mathematical statements, and even their proofs, become first-class mathematical objects.
In type theory, we represent mathematical statements by types, which can be regarded simultaneously as both mathematical constructions and mathematical assertions, a conception also known as \emph{propositions as types}.
\index{proposition!as types}%
Accordingly, we can regard a term $a : A$ as both an element of the type $A$ (or in homotopy type theory, a point of the space $A$), and at the same time, a proof of the proposition $A$.
To take an example, suppose we have sets $A$ and $B$ (discrete spaces),
\index{discrete!space}%
and consider the statement ``$A$ is isomorphic to $B$''.
In type theory, this can be rendered as:
\begin{narrowmultline*}
  \mathsf{Iso}(A,B) \defeq \narrowbreak
  \sm{f : A\to B}{g : B\to A}\Big(\big(\tprd{x:A} g(f(x)) = x\big) \times \big(\tprd{y:B}\, f(g(y)) = y\big)\Big).
\end{narrowmultline*}
%
Reading the type constructors $\Sigma, \Pi, \times$  here  as ``there exists'', ``for all'', and ``and'' respectively yields the usual formulation of ``$A$ and $B$ are isomorphic''; on the other hand, reading them as sums and products yields the \emph{type of all isomorphisms} between $A$ and $B$!  To prove that $A$ and $B$ are isomorphic, one  constructs a proof $p : \mathsf{Iso}(A,B)$, which is therefore the same  as constructing an isomorphism between $A$ and $B$, i.e., exhibiting a pair of functions $f, g$ together with \emph{proofs} that their composites are the respective identity maps.  The latter proofs, in turn, are nothing but homotopies of the appropriate sorts.  In this way, \emph{proving a proposition is the same as constructing an element of some particular type.}
In particular, to prove a statement of the form ``$A$ and $B$'' is just to prove $A$ and to prove $B$, i.e., to give an element of the type $A\times B$.
And to prove that $A$ implies $B$ is just to find an element of $A\to B$, i.e.\ a function from $A$ to $B$ (determining a mapping of proofs of $A$ to proofs of $B$).

The logic of propositions-as-types is flexible and supports many variations, such as using only a subclass of types to represent propositions.
In homotopy type theory, there are natural such subclasses arising from the fact that the system of all types, just like spaces in classical homotopy theory, is ``stratified'' according to the dimensions in which their higher homotopy structure exists or collapses.
In particular, Voevodsky has found a purely type-theoretic definition of \emph{homotopy $n$-types}, corresponding to spaces with no nontrivial homotopy information above dimension $n$.
(The $0$-types are the ``sets'' mentioned previously as satisfying Lawvere's axioms\index{Lawvere}.)
Moreover, with higher inductive types, we can universally ``truncate'' a type into an $n$-type; in classical homotopy theory this would be its $n^{\mathrm{th}}$ Postnikov\index{Postnikov tower} section.\index{n-type@$n$-type}
Particularly important for logic is the case of homotopy $(-1)$-types, which we call \emph{mere propositions}.
Classically, every $(-1)$-type is empty or contractible; we interpret these possibilities as the truth values ``false'' and ``true'' respectively.

Using all types as propositions yields a very ``constructive'' conception of logic; for more on this, see~\cite{kolmogorov,TroelstraI,TroelstraII}.
For instance, every proof that something exists carries with it enough information to actually find such an object; and every proof that ``$A$ or $B$'' holds is either a proof that $A$ holds or a proof that $B$ holds.
Thus, from every proof we can automatically extract an algorithm;\index{algorithm} \index{extraction of algorithms} this can be very useful in applications to computer programming.

On the other hand, however, this logic does diverge from the traditional understanding of existence proofs in mathematics.
In particular, it does not faithfully represent certain important classical principles of reasoning, such as the axiom of choice and the law of excluded middle.
For these we need to use the ``$(-1)$-truncated'' logic, in which only the homotopy $(-1)$-types represent propositions.

\index{axiom!of choice}%
More specifically, consider on one hand the \emph{axiom of choice}: ``if for every $x: A$ there exists a $y:B$ such that $R(x,y)$, there is a function $f : A\to B$ such that for all $x:A$ we have $R(x, f(x))$.''
The pure propositions-as-types notion of ``there exists'' is strong enough to make this statement simply provable --- yet it does not have all the consequences of the usual axiom of choice.
However, in $(-1)$-truncated logic, this statement is not automatically true, but is a strong assumption with the same sorts of consequences as its counterpart in classical\index{mathematics!classical} set theory.

\index{excluded middle}%
\index{univalence axiom}%
On the other hand, consider the \emph{law of excluded middle}: ``for all $A$, either $A$ or not $A$.''
Interpreting this in the pure propositions-as-types logic yields a statement that is inconsistent with the univalence axiom.
For since proving ``$A$'' means exhibiting an element of it, this assumption would give a uniform way of selecting an element from every nonempty type --- a sort of Hilbertian choice operator.
Univalence implies that the element of $A$ selected by such a choice operator must be invariant under all self-equivalences of $A$, since these are identified with self-identities and every operation must respect identity; but clearly some types have automorphisms with no fixed points, e.g.\ we can swap the elements of a two-element type.
\index{automorphism!fixed-point-free}%
However, the ``$(-1)$-truncated law of excluded middle'', though also not automatically true, may consistently be assumed with most of the same consequences as in classical mathematics.

In other words, while the pure propositions-as-types logic is ``constructive'' in the strong algorithmic sense mentioned above, the default $(-1)$-truncated logic is ``constructive'' in a different sense (namely, that of the logic formalized by Heyting under the name ``intuitionistic''); and to the latter we may freely add the axioms of choice and excluded middle to obtain a logic that may be called ``classical''.
Thus, homotopy type theory is compatible with both constructive and classical conceptions of logic, and many more besides.
\index{logic!constructive vs classical}%
Indeed, the homotopical perspective reveals that classical and constructive logic can coexist, as endpoints of a spectrum of different systems, with an infinite number of possibilities in between (the homotopy $n$-types for $-1 < n < \infty$).
We may speak of ``\LEM{n}'' and ``\choice{n}'', with $\choice{\infty}$ being provable and \LEM{\infty} inconsistent with univalence, while $\choice{-1}$ and $\LEM{-1}$ are the versions familiar to classical mathematicians (hence in most cases it is appropriate to assume the subscript $(-1)$ when none is given).  Indeed, one can even have useful systems in which only \emph{certain} types satisfy such further ``classical'' principles, while types in general remain ``constructive''.\index{excluded middle}\index{axiom!of choice}%%

It is worth emphasizing that univalent foundations does not \emph{require} the use of constructive or intuitionistic logic.\index{logic!intuitionistic}\index{logic!constructive} %
Most of classical mathematics which depends on the law of excluded middle and the axiom of choice can be performed in univalent foundations, simply by assuming that these two principles hold (in their proper, $(-1)$-truncated, form).
However, type theory does encourage avoiding these principles when they are unnecessary, for several reasons.

First of all, every mathematician knows that a theorem is more powerful when proven using fewer assumptions, since it applies to more examples.
The situation with \choice{} and \LEM{} is no different:
type theory admits many interesting ``nonstandard'' models, such as in sheaf toposes,\index{topos} where classicality principles such as \choice{} and \LEM{} tend to fail.
Homotopy type theory admits similar models in higher toposes, such as are studied in~\cite{ToenVezzosi02,Rezk05,lurie:higher-topoi}.
Thus, if we avoid using these principles, the theorems we prove will be valid internally to all such models.

Secondly, one of the additional virtues of type theory is its computable character.
In addition to being a foundation for mathematics, type theory is a formal theory of computation, and can be treated as a powerful programming language.
\index{programming}%
From this perspective, the rules of the system cannot be chosen arbitrarily the way set-theoretic axioms can: there must be a harmony between them which allows all proofs to be ``executed'' as programs.
We do not yet fully understand the new principles introduced by homotopy type theory, such as univalence and higher inductive types, from
this point of view, but the basic outlines are emerging; see, for example,~\cite{lh:canonicity}.
It has been known for a long time, however, that principles such as \choice{} and \LEM{} are fundamentally antithetical to computability, since they assert baldly that certain things exist without giving any way to compute them.
Thus, avoiding them is necessary to maintain the character of type theory as a theory of computation.

Fortunately, constructive reasoning is not as hard as it may seem.
In some cases, simply by rephrasing some definitions, a theorem can be made constructive and its proof more elegant.
Moreover, in univalent foundations this seems to happen more often.
For instance:
\begin{enumerate}
\item In set-theoretic foundations, at various points in homotopy theory and category theory one needs the axiom of choice to perform transfinite constructions.
  But with higher inductive types, we can encode these constructions directly and constructively.
  In particular, none of the ``synthetic'' homotopy theory in \cref{cha:homotopy} requires \LEM{} or \choice{}.
\item In set-theoretic foundations, the statement ``every fully faithful and essentially surjective functor is an equivalence of categories'' is equiv\-a\-lent to the axiom of choice.
  But with the univalence axiom, it is just \emph{true}; see \cref{cha:category-theory}.
\item In set theory, various circumlocutions are required to obtain notions of ``cardinal number'' and ``ordinal number'' which canonically represent isomorphism classes of sets and well-ordered sets, respectively --- possibly involving the axiom of choice or the axiom of foundation.
  But with univalence and higher inductive types, we can obtain such representatives directly by truncating the universe; see \cref{cha:set-math}.
\item In set-theoretic foundations, the definition of the real numbers as equivalence classes of Cauchy sequences requires either the law of excluded middle or the axiom of (countable) choice to be well-behaved.
  But with higher inductive types, we can give a version of this definition which is well-behaved and avoids any choice principles; see \cref{cha:real-numbers}.
\end{enumerate}
Of course, these simplifications could as well be taken as evidence that the new methods will not, ultimately, prove to be really constructive.  However, we emphasize again that the reader does not have to care, or worry, about constructivity in order to read this book.  The point is that in all of the above examples, the version of the theory we give has independent advantages, whether or not \LEM{} and \choice{} are assumed to be available.  Constructivity, if attained, will be an added bonus.\index{constructivity}%

Given this discussion of adding new principles such as univalence, higher inductive types, \choice{}, and \LEM{}, one may wonder whether the resulting system remains consistent.
(One of the original virtues of type theory, relative to set theory, was that it can be seen to be consistent by proof-theoretic means).
As with any foundational system, consistency\index{consistency} is a relative question: ``consistent with respect to what?''
The short answer is that all of the constructions and axioms considered in this book have a model in the category of Kan\index{Kan complex} complexes, due to Voevodsky~\cite{klv:ssetmodel} (see~\cite{ls:hits} for higher inductive types).
Thus, they are known to be consistent relative to ZFC (with as many inaccessible cardinals
\index{inaccessible cardinal}\index{consistency}%
as we need nested univalent universes).
Giving a more traditionally type-theoretic account of this consistency is work in progress (see,
e.g.,~\cite{lh:canonicity,coquand2012constructive}).

We summarize the different points of view of the type-theoretic operations in \cref{tab:pov}.

\begin{table}[htb]
  \centering
  \OPTsmalltable
 \begin{tabular}{lllll}
    \toprule
       Types && Logic & Sets & Homotopy\\ \addlinespace[2pt]
    \midrule
       $A$ && proposition & set & space\\ \addlinespace[2pt]
       $a:A$ && proof & element & point \\ \addlinespace[2pt]
       $B(x)$ && predicate & family of sets & fibration \\ \addlinespace[2pt]
       $b(x) : B(x)$ && conditional proof & family of elements & section\\ \addlinespace[2pt]
       $\emptyt, \unit$ && $\bot, \top$ & $\emptyset, \{ \emptyset \}$ & $\emptyset, *$\\ \addlinespace[2pt]
       $A + B$ && $A\vee B$ & disjoint union & coproduct\\ \addlinespace[2pt]
       $A\times B$ && $A\wedge B$ & set of pairs & product space\\ \addlinespace[2pt]
       $A\to B$ && $A\Rightarrow B$ & set of functions & function space\\ \addlinespace[2pt]
       $\sm{x:A}B(x)$ &&  $\exists_{x:A}B(x)$ & disjoint sum & total space\\ \addlinespace[2pt]
       $\prd{x:A}B(x)$ &&  $\forall_{x:A}B(x)$ & product & space of sections\\ \addlinespace[2pt]
       $\mathsf{Id}_{A}$ && equality $=$ & $\setof{\pairr{x,x} | x\in A}$ & path space $A^I$ \\ \addlinespace[2pt]
    \bottomrule
  \end{tabular}
  \caption{Comparing points of view on type-theoretic operations}\label{tab:pov}
\end{table}

\index{mathematics!constructive|)}%

\subsection*{Open problems} 

\index{open!problem|(}%

For those interested in contributing to this new branch of mathematics, it may be encouraging to know that there are many interesting open questions.

\index{univalence axiom!constructivity of}%
Perhaps the most pressing of them is the ``constructivity'' of the Univalence Axiom, posed by Voevodsky in \cite{Universe-poly}.
The basic system of type theory follows the structure of Gentzen's natural deduction. Logical connectives are defined by their introduction rules, and have elimination rules justified by computation rules. Following this pattern, and using Tait's computability method, originally designed to analyse G\"odel's Dialectica interpretation, one can show the property of \emph{normalization} for type theory. This in turn implies important properties such as decidability of type-checking (a crucial property since type-checking corresponds to proof-checking, and one can argue that we should be able to ``recognize a proof when we see one''), and the so-called ``canonicity\index{canonicity} property'' that any closed term of the type of natural numbers reduces to a numeral. This last property, and the uniform structure of introduction/elimination rules, are lost when one extends type theory with an axiom, such as the axiom of function extensionality, or the univalence axiom. Voevodsky has formulated a precise mathematical conjecture connected to this question of canonicity for type theory extended with the axiom of Univalence: given a closed term of the type of natural numbers, is it always possible to find a numeral and a proof that this term is equal to this numeral, where this proof of equality may itself use the univalence axiom? More generally, an important issue is whether it is possible to provide a constructive justification of the univalence axiom.
What about if one adds other homotopically motivated constructions, like higher inductive types?
These questions remain open at the present time, although methods are currently being developed to try to find answers.

Another basic issue is the difficulty of working with types, such as the natural numbers, that are essentially sets (i.e., discrete spaces),
\index{discrete!space}%
containing only trivial paths.
At present, homotopy type theory can really only characterize spaces up to homotopy equivalence, which means that these ``discrete spaces'' may only be \emph{homotopy equivalent} to discrete spaces.
Type-theoretically, this means there are many paths that are equal to reflexivity, but not \emph{judgmentally} equal to it (see \cref{sec:types-vs-sets} for the meaning of ``judgmentally'').
While this homotopy-invariance has advantages, these ``meaningless'' identity terms do introduce needless complications into arguments and constructions, so it would be convenient to have a systematic way of eliminating or collapsing them.
% In some cases, the proliferation of such superfluous identity terms makes it very difficult or impossible to formulate what should be a straightforward concept, such as the definition of a (semi-)simplicial type.

A more specialized, but no less important, problem is the relation between homotopy type theory and the research on \emph{higher toposes}%
\index{.infinity1-topos@$(\infty,1)$-topos}
currently happening at the intersection of higher category theory and homotopy theory.
There is a growing conviction among those familiar with both subjects that they are intimately connected.
For instance, the notion of a univalent universe should coincide with that of an object classifier, while higher inductive types should be an ``elementary'' reflection of local presentability.
More generally, homotopy type theory should be the ``internal language'' of $(\infty,1)$-toposes, just as intuitionistic higher-order logic is the internal language of ordinary 1-toposes.
Despite this general consensus, however, details remain to be worked out --- in particular, questions of coherence and strictness remain to be addressed  --- and doing so will undoubtedly lead to further insights into both concepts.

\index{mathematics!formalized}%
But by far the largest field of work to be done is in the ongoing formalization of everyday mathematics in this new system.
Recent successes in formalizing some facts from basic homotopy theory and category theory have been encouraging; some of these are described in \cref{cha:homotopy,cha:category-theory}.
Obviously, however, much work remains to be done.

\index{open!problem|)}%

The homotopy type theory community maintains a web site and group blog at \url{http://homotopytypetheory.org}, as well as a discussion email list.
Newcomers are always welcome!


\subsection*{How to read this book}

This book is divided into two parts.
\cref{part:foundations}, ``Foundations'', develops the fundamental concepts of homotopy type theory.
This is the mathematical foundation on which the development of specific subjects is built, and which is required for the understanding of the univalent foundations approach. To a programmer, this is ``library code''.
Since univalent foundations is a new and different kind of mathematics, its basic notions take some getting used to; thus \cref{part:foundations} is fairly extensive.

\cref{part:mathematics}, ``Mathematics'', consists of four chapters that build on the basic notions of \cref{part:foundations} to exhibit some of the new things we can do with univalent foundations in four different areas of mathematics: homotopy theory (\cref{cha:homotopy}), category theory (\cref{cha:category-theory}), set theory (\cref{cha:set-math}), and real analysis (\cref{cha:real-numbers}).
The chapters in \cref{part:mathematics} are more or less independent of each other, although occasionally one will use a lemma proven in another.

A reader who wants to seriously understand univalent foundations, and be able to work in it, will eventually have to read and understand most of \cref{part:foundations}.
However, a reader who just wants to get a taste of univalent foundations and what it can do may understandably balk at having to work through over 200 pages before getting to the ``meat'' in \cref{part:mathematics}.
Fortunately, not all of \cref{part:foundations} is necessary in order to read the chapters in \cref{part:mathematics}.
Each chapter in \cref{part:mathematics} begins with a brief overview of its subject, what univalent foundations has to contribute to it, and the necessary background from \cref{part:foundations}, so the courageous reader can turn immediately to the appropriate chapter for their favorite subject.
For those who want to understand one or more chapters in \cref{part:mathematics} more deeply than this, but are not ready to read all of \cref{part:foundations}, we provide here a brief summary of \cref{part:foundations}, with remarks about which parts are necessary for which chapters in \cref{part:mathematics}.

\cref{cha:typetheory} is about the basic notions of type theory, prior to any homotopical interpretation.
A reader who is familiar with Martin-L\"of type theory can quickly skim it to pick up the particulars of the theory we are using.
However, readers without experience in type theory will need to read \cref{cha:typetheory}, as there are many subtle differences between type theory and other foundations such as set theory.

\cref{cha:basics} introduces the homotopical viewpoint on type theory, along with the basic notions supporting this view, and describes the homotopical behavior of each component of the type theory from \cref{cha:typetheory}.
It also introduces the \emph{univalence axiom} (\cref{sec:compute-universe}) --- the first of the two basic innovations of homotopy type theory.
Thus, it is quite basic and we encourage everyone to read it, especially \crefrange{sec:equality}{sec:basics-equivalences}.

\cref{cha:logic} describes how we represent logic in homotopy type theory, and its connection to classical logic as well as to constructive and intuitionistic logic.
Here we define the law of excluded middle, the axiom of choice, and the axiom of propositional resizing (although, for the most part, we do not need to assume any of these in the rest of the book), as well as the \emph{propositional truncation} which is essential for representing traditional logic.
This chapter is essential background for \cref{cha:set-math,cha:real-numbers}, less important for \cref{cha:category-theory}, and not so necessary for \cref{cha:homotopy}.

\cref{cha:equivalences,cha:induction} study two special topics in detail: equivalences (and related notions) and generalized inductive definitions.
While these are important subjects in their own rights and provide a deeper understanding of homotopy type theory, for the most part they are not necessary for \cref{part:mathematics}.
Only a few lemmas from \cref{cha:equivalences} are used here and there, while the general discussions in \cref{sec:bool-nat,sec:strictly-positive,sec:generalizations} are helpful for providing the intuition required for \cref{cha:hits}.
The generalized sorts of inductive definition discussed in \cref{sec:generalizations} are also used in a few places in \cref{cha:set-math,cha:real-numbers}.

\cref{cha:hits} introduces the second basic innovation of homotopy type theory --- \emph{higher inductive types} --- with many examples.
Higher inductive types are the primary object of study in \cref{cha:homotopy}, and some particular ones play important roles in \cref{cha:set-math,cha:real-numbers}.
They are not so necessary for \cref{cha:category-theory}, although one example is used in \cref{sec:rezk}.

Finally, \cref{cha:hlevels} discusses homotopy $n$-types and related notions such as $n$-connected types.
These notions are important for \cref{cha:homotopy}, but not so important in the rest of \cref{part:mathematics}, although the case $n=-1$ of some of the lemmas are used in \cref{sec:piw-pretopos}.

This completes \cref{part:foundations}.
As mentioned above, \cref{part:mathematics} consists of four largely unrelated chapters, each describing what univalent foundations has to offer to a particular subject.

Of the chapters in \cref{part:mathematics}, \cref{cha:homotopy} (Homotopy theory) is perhaps the most radical.
Univalent foundations has a very different ``synthetic'' approach to homotopy theory in which homotopy types are the basic objects (namely, the types) rather than being constructed using topological spaces or some other set-theoretic model.
This enables new styles of proof for classical theorems in algebraic topology, of which we present a sampling, from $\pi_1(\Sn^1)=\Z$ to the Freudenthal suspension theorem.

In \cref{cha:category-theory} (Category theory), we develop some basic (1-)category theory, adhering to the principle of the univalence axiom that \emph{equality is isomorphism}.
This has the pleasant effect of ensuring that all definitions and constructions are automatically invariant under equivalence of categories: indeed, equivalent categories are equal just as equivalent types are equal.
(It also has connections to higher category theory and higher topos theory.)

\cref{cha:set-math} (Set theory) studies sets in univalent foundations.
The category of sets has its usual properties, hence provides a foundation for any mathematics that doesn't need homotopical or higher-categorical structures.
We also observe that univalence makes cardinal and ordinal numbers a bit more pleasant, and that higher inductive types yield a cumulative hierarchy satisfying the usual axioms of Zermelo--Fraenkel set theory.

In \cref{cha:real-numbers} (Real numbers), we summarize the construction of Dedekind real numbers, and then observe that higher inductive types allow a definition of Cauchy real numbers that avoids some associated problems in constructive mathematics.
Then we sketch a similar approach to Conway's surreal numbers.

Each chapter in this book ends with a Notes section, which collects historical comments, references to the literature, and attributions of results, to the extent possible.
We have also included Exercises at the end of each chapter, to assist the reader in gaining familiarity with doing mathematics in univalent foundations.

Finally, recall that this book was written as a massively collaborative effort by a large number of people.
We have done our best to achieve consistency in terminology and notation, and to put the mathematics in a linear sequence that flows logically, but it is very likely that some imperfections remain.
We ask the reader's forgiveness for any such infelicities, and welcome suggestions for improvement of the next edition.


% Local Variables:
% TeX-master: "hott-online"
% End:

% !TEX root = runners-in-action.tex

\section{Algebraic effects, handlers, and runners}
\label{sect:runnersbyexample}

We begin with a short overview of the theory of algebraic effects and handlers, as well as
runners. To keep focus on how runners give rise to a programming
concept, we work naively in set theory. Nevertheless, we use
category-theoretic language as appropriate, to make it clear that there are no essential
obstacles to extending our work to other settings (we return to this point in \cref{sec:semantics-types}).

\subsection{Algebraic effects and handlers}
\label{sect:algebraiceffects}

There is by now no lack of material on the algebraic approach to structuring computational effects.
For an introductory treatment we refer to~\cite{Bauer:WhatIsAlgebraic}, while of course
also recommend the seminal papers by Plotkin and
Power~\cite{Plotkin:SemanticsForAlgOperations,Plotkin:NotionsOfComputation}. The brief
summary given here only recalls the essentials and introduces notation.

An \emph{(algebraic) signature} is given by a set $\sig$ of \emph{operation symbols},
and for each $\op \in \sig$ its \emph{operation signature}
%
$
  \op : A_\op \opto B_\op
$,
%
where  $A_\op$ and $B_\op$ are called the \emph{parameter} and \emph{arity} set.
A \emph{$\sig$-structure} $\M$ is given by a carrier set $\Mcarrier$, and for each
operation symbol $\op \in \sig$, a map $\op_{\M} : A_\op \times (B_\op \expto \Mcarrier) \to \Mcarrier$,
where~$\expto$ is set exponentiation. The \emph{free $\sig$-structure~$\Tree{\sig}{X}$}
over a set~$X$ is the set of well-founded trees generated inductively by
%
\begin{itemize}
\item $\retTree{x} \in \Tree{\sig}{X}$, for every $x \in X$, and
\item $\op(a, \kappa) \in \Tree{\sig}{X}$, for every $\op \in \sig$, $a \in A_\op$, and $\kappa : B_\op \to \Tree{\sig}{X}$.
\end{itemize}
%
We are abusing notation in a slight but standard way, by using $\op$ both as the
name of an operation and a tree-forming constructor.
%
The elements of $\Tree{\sig}{X}$ are called \emph{computation trees}: a leaf
$\retTree{x}$ represents a pure computation returning a value $x$, while
$\op(a, \kappa)$ represents an effectful computation that calls $\op$ with
parameter~$a$ and continuation~$\kappa$, which expects a result from~$B_\op$.

An \emph{algebraic theory $\Th = (\Thsig, \Theq)$} is given by a \emph{signature~$\Thsig$} and
a set of \emph{equations~$\Theq$}.
%
The equations $\Theq$ express computational behaviour via interactions between
operations, and
are written in a suitable formalism, e.g.,~\cite{Plotkin:HandlingEffects}. We
explain these by way of examples, as the precise details do not matter for our purposes.
Let $\Zero = \{\,\}$ be the empty set and $\One = \{\star\}$ the standard singleton.

\begin{example}
  \label{ex:state}
  Given a set $C$ of possible states, the theory of \emph{$C$-valued state} has two operations, whose
  somewhat unusual naming will become clear later on,
  %
  \begin{equation*}
    \siggetenv : \One \opto C,
    \qquad\qquad
    \sigsetenv : C \opto \One
  \end{equation*}
  %
  and the equations (where we elide appearances of $\star$):
  %
  \begin{gather*}
    \siggetenv (\lam{c} \sigsetenv(c, \kappa)) = \kappa,
    \qquad
    \sigsetenv (c, \siggetenv\, \kappa) = \sigsetenv(c, \kappa\, c), \\
    \sigsetenv (c, \sigsetenv(c', \kappa)) = \sigsetenv(c', \kappa).
  \end{gather*}
  %
  For example,
  the second equation states that reading state right after setting it to~$c$ gives precisely~$c$.
  The third equation states that $\sigsetenv$ overwrites the state.
\end{example}

\begin{example}
  \label{ex:exceptions}
  Given a set of exceptions $E$, the algebraic theory of \emph{$E$-many exceptions} is given by
  a single operation $\sigraise : E \opto \Zero$, and no equations.
\end{example}

A \emph{$\Th$-model}, also called a \emph{$\Th$-algebra}, is a $\Thsig$-structure which
satisfies the equations in $\Theq$. The \emph{free $\Th$-model} over a set~$X$ is constructed
as the quotient
%
\begin{equation*}
  \FreeAlg{\Th}{X} = \Tree{\Thsig}{X}/{\sim}
\end{equation*}
%
by the $\Thsig$-congruence $\sim$ generated by $\Theq$. Each $\op \in \Thsig$
is interpreted in the free model as the map
$(a, \kappa) \mapsto [\op(a, \kappa)]$, where $[{-}]$ is the $\sim$-equivalence class.

$\FreeAlg{\Th}{-}$ is the functor part of a \emph{monad} on sets, whose \emph{unit} at a
set~$X$ is
%
\begin{equation*}
  \xymatrix@C=3em@R=2em@M=0.5em{
    {X} \ar[r]^(0.35){\retTree{}}
    &
    {\Tree{\Thsig}{X}} \ar@{->>}[r]^{[{-}]}
    &
    {\FreeAlg{\Th}{X}.}
  }
\end{equation*}
%
The \emph{Kleisli extension} for this monad is then the operation which lifts any map \linebreak
$f : X \to \Tree{\Thsig}{Y}$ to the map $\lift{f} : \FreeAlg{\Thsig}{X} \to \FreeAlg{\Thsig}{Y}$,
given by
%
\begin{equation*}
  \lift{f}\,[\retTree{x}] \defeq f \, x,
  \qquad\qquad
  \lift{f}\,[\op(a, \kappa)] \defeq [\op(a, \lift{f} \circ \kappa)].
\end{equation*}
%
That is, $\lift{f}$ traverses a computation tree and replaces each leaf $\retTree{x}$
with $f\,x$.

The preceding construction of free models and the monad may be retro-fitted to an
algebraic signature~$\sig$, if we construe~$\sig$ as an algebraic theory with no
equations. In this case~$\sim$ is just equality, and so we may omit the quotient and the
pesky equivalence classes. Thus the carrier of the free $\sig$-model is the set of
well-founded trees $\Tree{\sig}{X}$, with the evident monad structure.

A fundamental insight of Plotkin and
Power~\cite{Plotkin:SemanticsForAlgOperations,Plotkin:NotionsOfComputation} was that many
computational effects may be adequately described by algebraic theories, with
the elements of free models corresponding to effectful computations. For example, the monads
induced by the theories from \cref{ex:state,ex:exceptions} are respectively isomorphic to the usual \emph{state monad}
$\St{C}\,X \defeq (C \Rightarrow X \times C)$ and the \emph{exceptions monad}
$\Exc{E}\,X \defeq X + E$.

Plotkin and Pretnar~\cite{Plotkin:HandlingEffects} further
observed that the universal property of free models may be used to model a
programming concept known as \emph{handlers}. Given a $\Th$-model $\M$ and a map
$f : X \to \Mcarrier$, the universal property of the free $\Th$-model gives us a
unique $\Th$-homomorphism $\freelift{f} : \FreeAlg{\Th}{X} \to \Mcarrier$ satisfying
%
\begin{equation*}
  \freelift{f} \, [\retTree{x}] = f\,x,
  \qquad\qquad
  \freelift{f} \, [\op(a, \kappa)] = \op_\M(a, \freelift{f} \circ \kappa).
\end{equation*}

A handler for a theory $\Th$ in a language such as \pl{Eff} amounts to a model 
$\M$ whose carrier $\Mcarrier$
is the carrier $\FreeAlg{\Th'}{Y}$ of the free model
for some other theory~$\Th'$, while the associated handling
construct is the induced $\Th$-homomorphism $\FreeAlg{\Th}{X} \to \FreeAlg{\Th'}{Y}$.
Thus handling transforms computations with effects~$\Th$ to computations with
effects~$\Th'$. There is however no restriction on how a handler implements an
operation, in particular, it may use its continuation in an arbitrary
fashion.
%
We shall put the universal property of free models to good use as well, while
making sure that the continuations are always used affinely.

\subsection{Runners}
\label{sect:purerunners}

Much like monads, handlers are useful for simulating computational effects, because they
allow us to transform $\Th$-computations to $\Th'$-computations. However, eventually there
has to be a ``top level'' where such transformations cease and actual computational
effects happen. For these we need another concept, known as
\emph{runners}~\cite{Uustalu:Runners}.
%
Runners are equivalent to the concept of
\emph{comodels}~\cite{Plotkin:TensorsOfModels,Power:Comodels}, which are ``just
models in the opposite category'', although one has to apply the motto
correctly by using powers and co-powers where seemingly exponentials and products would
do. 
Without getting into the intricacies, let us spell out the definition.

\begin{definition}
  A \emph{runner} $\R$ for a signature $\sig$ is given by a carrier set~$\Rcarrier$ together with, for
  each $\op \in \sig$, a
  \emph{co-operation~$\coop_{\R} : A_\op \to (\Rcarrier \expto B_\op \times \Rcarrier)$.}
\end{definition}

%
Runners are usually defined to have co-operations in the equivalent uncurried form
$\coop_\R : A_\op \times \Rcarrier \to B_\op \times \Rcarrier$, but that is less convenient for our purposes.

Runners may be defined more generally for theories $\Th$, rather than just signatures,
by requiring that the co-operations satisfy $\Theq$. We shall have no use for these,
although we expect no obstacles in incorporating them into our work.

A runner tells us what to do when an effectful computation reaches the top-level runtime
environment. Think of~$\Rcarrier$ as the set of configurations of the runtime environment. Given
the current configuration $c \in \Rcarrier$, the operation $\op(a, \kappa)$ is executed as the
corresponding co-operation $\coop_\R\,a\,c$ whose result $(b, c') \in B_\op \times \Rcarrier$ gives
the result of the operation $b$ and 
the next runtime configuration $c'$. The continuation $\kappa\,b$
then proceeds in runtime configuration~$c'$.

It is not too difficult to turn this idea into a mathematical model. For any
$X$, the co-operations induce a $\sig$-structure $\M$ with
$\Mcarrier \defeq \St{\Rcarrier} X = (\Rcarrier \expto X \times \Rcarrier)$ 
and operations $\op_\M : A_\op \times (B_\op \expto \St{\Rcarrier} X) \to \St{\Rcarrier} X$
given by
%
\begin{equation*}
  \op_\M (a, \kappa) \defeq \lam{c} \kappa\, (\pi_1 (\coop_\R\,a\,c))\, (\pi_2 (\coop_\R\,a\,c)).
\end{equation*}
%
We may then use the universal property of the free $\sig$-model to obtain a $\sig$-homomorphism
$\runh_X : \Tree{\sig}{X} \to \St{\Rcarrier} X$ satisfying the equations
%
\[
  \runh_X(\retTree{x}) = \lam{c} (x, c),
  \qquad\qquad
  \runh_X(\op(a, \kappa)) = \op_\M(a, \runh_X \circ \kappa).
\]
%
The map $\runh_X$ precisely captures the idea that a runner
\emph{runs computations} by transforming (static) computation trees into
state-passing maps. Note how in the above definition of $\op_\M$, the
continuation~$\kappa$ is used in a controlled way, as it appears precisely once
as the head of the outermost application. In terms of programming, this
corresponds to linear use in a tail-call position.

Runners are less ad-hoc than they may seem. First, notice that $\op_\M$ is just the
composition of the co-operation $\coop_\R$ with
the state monad's Kleisli extension of the continuation $\kappa$, and so is
the standard way of turning \emph{generic effects} into $\sig$-structures~\cite{Plotkin:AlgOperations}.
%
Second, the map $\runh_X$ is the component at $X$ of a monad morphism
$\runh : \Tree{\sig}{{-}} \to \St{\Rcarrier}$. Møgelberg \& Staton~\cite{Mogelberg:LinearUsageOfState}, as
well as Uustalu~\cite{Uustalu:Runners}, showed that the passage from a runner~$\R$ to the
corresponding monad morphism~$\runh$ forms a one-to-one correspondence between the former and the
latter.

As defined, runners are too restrictive a model of top-level computation, because the only
effect available to co-operations is state, but in practice the runtime
environment may also signal errors and perform other effects, by calling its own runtime
environment. We are led to the following generalisation.

\begin{definition}
  For a signature $\sig$ and monad $\T$, a \emph{$\T$-runner $\R$} for~$\sig$, 
  or just an \emph{effectful runner}, 
  is given by, for each $\op \in \sig$, a \emph{co-operation}
  $\coop_\R : A_\op \to \T B_\op$.
\end{definition}

The correspondence between runners and monad morphisms still holds.

\begin{proposition}
  \label{prop:monadmorphism}
  For a signature $\sig$ and a monad $\T$, the monad morphisms $\Tree{\sig}{{-}} \to \T$
  are in one-to-one correspondence with $\T$-runners for~$\sig$.
\end{proposition}

\begin{proof}
  This is an easy generalisation of the correspondence for 
  ordinary runners. Let us fix a signature $\sig$, and a monad $\T$ 
  with unit $\eta$ and Kleisli extension $\lift{{-}}$.

  Let $\R$ be a $\T$-runner for $\sig$. For any set $X$, $\R$ induces a $\sig$-structure
  $\M$ with $\Mcarrier \defeq \T X$ and 
  $\op_\M : A_\op \times (B_\op \expto \T X) \to \T X$ defined as
  %
  $
    \op_\M (a, \kappa) \defeq \lift{\kappa} (\coop_R\,a)
  $.
  %
  As before, the universal property of the free model $\Tree{\sig}{X}$ provides a unique
  $\sig$-homomorphism $\runh_X : \Tree{\sig}{X} \to \T X$, satisfying the equations
  %
  \begin{equation*}
    \runh_X (\retTree{x}) = \eta_X(x),
    \qquad\qquad
    \runh_X (\op(a, \kappa)) = \op_\M (a, \runh_X \circ \kappa).
  \end{equation*}
  %
  The maps $\runh_X$ collectively give us the desired monad morphism $\runh$ induced by $\R$.
  
  Conversely, given a monad morphism $\theta : \Tree{\sig}{{-}} \to \T$, we may recover a
  $\T$-runner~$\R$ for $\sig$ by defining the co-operations as
  $
    \coop_\R \, a \defeq \theta_{B_\op} (\op (a, \lam{b} \retTree{b}))
  $.
  It is not hard to check that we have described a one-to-one correspondence.
  \qed
\end{proof}



%%% Local Variables:
%%% mode: latex
%%% TeX-master: "runners-in-action"
%%% End:

% !TEX root = runners-in-action.tex

\section{Programming with runners}
\label{sec:programming-with-runners}

If ordinary runners are not general enough, the effectful ones are too general: 
parameterised by arbitrary monads $\T$, they do
not combine easily and they lack a clear notion of resource management. Thus,
we now engineer more specific monads whose associated runners can be turned into a
programming concept.
%
While we give up complete generality, the monads presented below are still quite
versatile, as they are parameterised by arbitrary algebraic signatures $\Sigma$,
and so are extensible and support various combinations of effects.

\subsection{The user and kernel monads}
\label{sec:user-kernel-monads}

Effectful source code running inside a runtime environment is just one example of a more
general phenomenon in which effectful computations are enveloped by a layer that provides
a supervised access to external resources: a user process is controlled by a kernel, a web
page by a browser, an operating system by hardware, or a virtual machine, etc. We shall
adopt the parlance of software systems, and refer to the two layers generically as the
\emph{user} and \emph{kernel} code.
%
Since the two kinds of code need not, and will not, use the same effects, each
will be described by its own algebraic theory and compute in its own monad.

We first address the kernel theory. 
Specifically, we look for an algebraic theory such that effectful runners for the induced monad
satisfy the following desiderata:
%
\begin{enumerate}
\item Runners support management and controlled finalisation of resources.
\item Runners may use further external resources.
\item Runners may signal failure caused by unavoidable circumstances.
\end{enumerate}

The totality of external resources
available to user code appears as a stateful external environment, even though it
has no direct access to it. Thus, kernel computations should carry state. We
achieve this by incorporating into the kernel theory the operations $\siggetenv$
and $\sigsetenv$, and equations for state from \cref{ex:state}.

Apart from managing state, kernel code should have access to further
effects, which may be true external effects, or some outer
layer of runners. In either case, we should allow the kernel code to call
operations from a given signature~$\sig$.

Because kernel computations ought to be able to signal failure, we should
include an exception mechanism. In practice, many programming languages and
systems have two flavours of exceptions, variously called recoverable and fatal,
checked and unchecked, exceptions and errors, etc. One kind, which we call just
\emph{exceptions}, is raised by kernel code when a situation requires special
attention by user code. The other kind, which we call \emph{signals},
indicates an unrecoverable condition that prevents normal execution of user
code. These correspond precisely to the two standard ways of combining
exceptions with state, namely the coproduct and the tensor of algebraic
theories~\cite{Hyland:CombiningEffects}. The coproduct simply adjoins exceptions
$\sigraise : E \leadsto \Zero$ from \cref{ex:exceptions} to the theory of
state, while the tensor extends the theory of state with signals
$\sigkill : S \leadsto \Zero$, together with equations
%
\begin{equation}
  \label{eq:kill-state}%
  \siggetenv(\lam{c} \sigkill\,s) = \sigkill\,s,
  \qquad\qquad
  \sigsetenv(c, \sigkill\,s) = \sigkill\,s.
\end{equation}
%
These equations say that a signal discards state, which makes it unrecoverable.

To summarise, the \emph{kernel theory} $\ThKK{C}{\sig}{E}{S}$ contains 
operations from a signature $\sig$, as well as state operations
$\siggetenv : \One \opto C$, $\sigsetenv : C \opto \One$, exceptions
$\sigraise : E \opto \Zero$, and signals $\sigkill : S \opto \Zero$, with equations for state
from \cref{ex:state}, equations~\eqref{eq:kill-state} relating state and
signals, and for each operation $\op \in \sig$, equations
%
\begin{align*}
  \siggetenv(\lam{c} \op(a, \kappa\,c)) &= \op(a, \lam{b} \siggetenv (\lam{c} \kappa\,c\,b)),\\
  \sigsetenv(c, \op(a, \kappa)) &= \op(a, \lam{b} \sigsetenv(c, \kappa\,b)),
\end{align*}
%
expressing that external operations do not interact with kernel state. 
It is not difficult to see that $\ThKK{C}{\sig}{E}{S}$ induces, up to
isomorphism, the \emph{kernel monad}
%
\begin{equation*}
  \KK{C}{\sig}{E}{S} X \quad\defeq\quad C \expto \Tree{\sig}{((X + E) \times C) + S}.
\end{equation*}

How about user code? It can of course call operations from a 
signature~$\sig$ (not necessarily the same as the kernel code), and because we
intend it to handle exceptions, it might as well have the ability to raise them.
However, user code knows nothing about signals and kernel state. Thus, we choose the \emph{user theory
  $\ThUU{\sig}{E}$} to be the algebraic theory with operations $\sig$, exceptions
$\sigraise : E \opto \Zero$, and no equations. This theory induces the \emph{user
  monad} $\UU{\sig}{E} X \defeq \Tree{\sig}{X + E}$.

\subsection{Runners as a programming construct}
\label{sec:runn-as-progr}

In this section, we turn the ideas presented so far into programming constructs.
We strive for a realistic result,
but when faced with several design options, we prefer simplicity and semantic
clarity. We focus here on translating the central concepts, and postpone
various details to \cref{sect:corecalculus}, where we present a full calculus.

We codify the idea of user and kernel computations by having syntactic
categories for each of them, as well as one for values. We use letters $M$, 
$N$ to indicate user computations, $K$, $L$ for kernel computations, 
and $V$, $W$ for values.

User and kernel code raise exceptions with operation $\tmkw{raise}$, and catch
them with exception handlers based on Benton and Kennedy's \emph{exceptional
  syntax}~\cite{Benton:ExceptionalSyntax},
%
\begin{equation*}
  \tmtry{M}{\{
    \tmreturn{x} \mapsto N,
    \ldots, \tmraise{e} \mapsto N_e, \ldots
  \}}, 
\end{equation*}
%
and analogously for kernel code. The familiar binding construct
%
$\tmlet{x}{M}{N}$
%
is simply shorthand for
%
$\tmtry{M}{\{\tmreturn{x} \mapsto N, \ldots, \tmraise{e} \mapsto \tmraise{e}, \ldots\}}$.

As a programming concept, a runner $R$ takes the form
%
\begin{equation*}
  \tmrunner{(\tm{op}\,x \mapsto K_{\tm{op}})_{\tm{op} \in \sig}}{C}, 
\end{equation*}
%
where each $K_\op$ is a kernel computation, with the variable $x$ bound in $K_{\tm{op}}$, so that
each clause $\tm{op} \, x \mapsto K_{\tm{op}}$ determines a co-operation for the
kernel monad. The subscript $C$ indicates the type of the state used by  
the kernel code $K_\op$.

The corresponding elimination form is a handling-like construct
%
\begin{equation}
  \label{eq:using}
  \tmrun{R}{V}{M}{F}, 
\end{equation}
%
which uses the co-operations of runner $R$ ``at'' initial kernel state~$V$ to
run user code~$M$, and finalises its return value, exceptions, and signals
with~$F$, see~\eqref{eq:finally-clause} below.
%
When user code $M$ calls an operation $\op$, the enveloping $\tmkw{run}$ construct runs the
corresponding co-operation $K_\op$ of $R$. While doing so, $K_\op$ might raise 
exceptions. But not every exception makes sense for every operation, and so
we assign to each operation $\op$ a set of exceptions $E_\op$ which the
co-operations implementing it may raise, by augmenting its operation signature with $E_\op$, as 
%
\begin{equation*}
  \op : \tysigop{A_\op}{B_\op}{E_\op}.
\end{equation*}
%
An exception raised by the co-operation $K_\op$ propagates back to the operation call in
the user code. Therefore, an operation call should have not only a continuation $x\,.\,M$
receiving a result, but also continuations $N_e$, one for each $e \in E_\op$,
%
\begin{equation*}
  \tmop{op}{}{V}{\tmcont x M}{\tmexccont N e {E_\op}}.
\end{equation*}
%
If $K_\op$ returns a value $b \in B_\op$, the execution proceeds 
as $M[b/x]$, and as $N_e$ if $K_\op$ raises an exception 
$e \in E_\op$. In examples, we use the generic versions of 
operations~\cite{Plotkin:AlgOperations}, written $\tmgeneff{op}{V}$,
which pass on return values and re-raise exceptions.

One can pass exceptions back to operation calls also in 
a language with handlers, such as \pl{Eff}, by 
changing the signatures of operations to
$A_\op \opto B_\op + E_\op$, and implementing 
the exception mechanism by hand, so that every
operation call is followed by a case distinction on $B_\op + E_\op$.
One is reminded of how operating system calls communicate 
errors back to user code as exceptional values.

A co-operation $K_\op$ may also send a signal, in which case the rest of the user code $M$
is skipped and the control proceeds directly to the corresponding case of the finalisation part~$F$ of the
$\tmkw{run}$ construct~\eqref{eq:using}, whose syntactic form is
%
\begin{equation}
  \label{eq:finally-clause}%
  \{ \tmreturn{x} \at c \mapsto N,
     \ldots, \tmraise{e} \at c \mapsto N_e,
     \ldots, \tmkill{s} \mapsto N_s,
     \ldots
  \}.
\end{equation}
%
Specifically, if $M$ returns a value $v$, then $N$ is evaluated with $x$ bound
to $v$ and $c$ to the final kernel state; if~$M$ raises an
exception~$e$ (either directly or indirectly via a co-operation of $R$), 
then $N_e$ is executed, again with $c$ bound to the final kernel state; and 
if a co-operation of $R$ sends a signal $s$, then $N_s$ is executed.

\begin{example}
  \label{ex:file-IO}%
  In anticipation of setting up the complete calculus we show how one can work with files.
  %
  The language implementors can provide an operation $\tm{open}$ which opens a
  file for writing and returns its file handle, an operation $\tm{close}$ which closes a
  file handle, and a runner $\mathsf{fileIO}$ that implements
  writing.
  Let us further suppose that $\mathsf{fileIO}$ may raise
  an exception $\mathsf{QuotaExceeded}$ if a write exceeds the user disk quota,
  and send a signal $\mathsf{IOError}$ 
  if an unrecoverable external error occurs.
  %
  The following code illustrates how to guarantee proper closing of the file:
  %
\begin{lstlisting}
using fileIO @ (open "hello.txt") run
  write "Hello, world."
finally {
  return x @ fh -> close fh,
  raise QuotaExceeded @ fh -> close fh,
  kill IOError -> return () }
\end{lstlisting}
  %
  Notice that the user code does not have direct access to the file handle.
  Instead, the runner holds it in its state, where it is available to the
  co-operation that implements $\tm{write}$. The finalisation block gets access to
  the file handle upon successful completion and raised exception, so it can close
  the file, but when a signal happens the finalisation cannot close the file,
  nor should it attempt to do so.

  We also mention that the code ``cheats'' by placing the call to $\tm{open}$ in a
  position where a value is expected. We should have $\tmkw{let}$-bound the file handle
  returned by $\tm{open}$ outside the $\tmkw{run}$ construct, which would make it clear that
  opening the file happens \emph{before} this construct (and that $\tm{open}$ is
  \emph{not} handled by the finalisation), but would also expose the file handle. Since
  there are clear advantages to keeping the file handle inaccessible, a realistic
  language should accept the above code and hoist computations from value positions
  automatically.
\end{example}

%%% Local Variables:
%%% mode: latex
%%% TeX-master: "runners-in-action"
%%% End:

% !TEX root = runners-in-action.tex


\newcommand{\tmcoop}[3]{\mathtt{\overline{#1}}~#2 \mapsto #3}

\section{A calculus for programming with runners}
\label{sect:corecalculus}

Inspired by the semantic notion of runners and the ideas of the previous
section, we now present a calculus for programming with co-operations and
runners, called $\lambdacoop$. It is a low-level fine-grain call-by-value
calculus~\cite{Levy:CBPV}, and as such could inspire an intermediate language
that a high-level language is compiled to.

\subsection{Types}
\label{sect:types}

\begin{figure}[tb]
  \parbox{\textwidth}{
  \centering
  \small
  \begin{align*}
  \text{Ground type $A$, $B$, $C$}
  \bnfis& \tybase      & &\text{base type} \\
  \bnfor& \tyunit       & &\text{unit type} \\
  \bnfor& \tyempty      & &\text{empty type} \\
  \bnfor& \typrod{A}{B} & &\text{product type} \\
  \bnfor& \tysum{A}{B}  & &\text{sum type}
  \\[1ex]
  \text{Constant signature:}
  \phantom{\bnfis}& \tmconst{f} : (A_1,\ldots,A_n) \to B
  \\[1ex]
  \text{Signature $\sig$}
  \bnfis& \{\op_1, \op_2, \ldots, \op_n\} \subset \Ops
  \\
  \text{Exception set $E$}
  \bnfis& \{e_1, e_2, \ldots, e_n\} \subset \Excs
  \\
  \text{Signal set $S$}
  \bnfis& \{s_1, s_2, \ldots, s_n\} \subset \Sigs
  \\[1ex]
  \text{Operation signature:}
  \phantom{\bnfis}& \op : \tysigop{A_\op}{B_\op}{E_\op}
  \\[1ex]
  \text{Value type $X$, $Y$, $Z$}
  \bnfis& A                                       & &\text{ground type} \\
  \bnfor& \typrod{X}{Y}                           & &\text{product type} \\
  \bnfor& \tysum{X}{Y}                            & &\text{sum type} \\
  \bnfor& \tyfun{X}{\tyuser{Y}{\Ueff}}           & &\text{user function type} \\
  \bnfor& \tyfunK{X}{\tykernel{Y}{\Keff}}         & &\text{kernel function type} \\
  \bnfor& \tyrunner{\sig}{\sig'}{S}{C}      & &\text{runner type}
  \\[1ex]
  \text{User (computation) type:}
  \phantom{\bnfor} &\tyuser{X}{\Ueff} \quad \text{where $\Ueff = (\sig, E)$}
  \\
  \text{Kernel (computation) type:}
  \phantom{\bnfor}& \tykernel{X}{\Keff} \quad \text{where $\Keff = (\sig, E, S, C)$}
  \end{align*}
  } 
  \caption{The types of $\lambdacoop$.}
  \label{fig:lambdacoop-types}
\end{figure}

The types of $\lambdacoop$ are shown in \cref{fig:lambdacoop-types}.
%
The \emph{ground types} contain \emph{base types}, and are closed under finite sums and
products. These are used in operation signatures and as types of kernel state. (Allowing
arbitrary types in either of these entails substantial complications that can be dealt
with but are tangential to our goals.) Ground types can also come with corresponding 
constant symbols~$\tmconst{f}$, each associated with a fixed \emph{constant signature}
$\tmconst{f} : (A_1,\ldots,A_n) \to B$.

We assume a supply of operation symbols $\Ops$, exception names $\Excs$, and
signal names $\Sigs$. Each operation symbol~$\op \in \Ops$ is equipped with an
\emph{operation signature} $\tysigop{A_\op}{B_\op}{E_\op}$, which specifies its
parameter type~$A_\op$ and arity type~$B_\op$, and the exceptions~$E_\op$ 
that the corresponding co-operations may raise in runners.

The \emph{value types} extend ground types with two 
function types, and a type of runners.
%
The \emph{user function type $\tyfun{X}{\tyuser{Y}{(\sig, E)}}$} classifies
functions taking arguments of type~$X$ to computations classified by the \emph{user
  (computation) type}~${\tyuser{Y}{(\sig, E)}}$, i.e., those that return values of
type~$Y$, and may call operations~$\sig$ and raise exceptions~$E$.
%
Similarly, the \emph{kernel function type
  $\tyfunK{X}{\tykernel{Y}{(\sig, E, S, C)}}$} classifies functions taking
arguments of type~$X$ to computations classified by the \emph{kernel
  (computation) type~$\tykernel{Y}{(\sig, E, S, C)}$}, i.e., those that return
values of type~$Y$, and may call operations~$\sig$, raise exceptions~$E$, send
signals~$S$, and use state of type~$C$. We note that the ingredients for user and kernel types
correspond precisely to the parameters of the user monad $\UU{\sig}{E}$ and the
kernel monad $\KK{C}{\sig}{E}{S}$ from \cref{sec:user-kernel-monads}.
%
Finally, the \emph{runner type $\tyrunner{\sig}{\sig'}{S}{C}$} classifies runners that
implement co-operations for the operations~$\sig$ as kernel computations which use 
operations~$\sig'$, send signals~$S$, and use state of type~$C$.


\subsection{Values and computations}
\label{sec:values-computations}


\begin{figure}[tp]
  \parbox{\textwidth}{
  \centering
  \small
  \abovedisplayskip=0pt
  \begin{align*}
  \intertext{\textbf{Values}}
  V, W
  \bnfis& x                                       & &\text{variable} \\
  \bnfor& \tmconst{f}(V_1,\ldots,V_n)                                       & &\text{ground constant} \\
  \bnfor& \tmunit                                 & &\text{unit} \\
  \bnfor& \tmpair{V}{W}                           & &\text{pair} \\
  \bnfor& \tminl[X,Y]{V} \bnfor \tminr[X,Y]{V}    & &\text{injection} \\
  \bnfor& \tmfun{x : X}{M}                        & &\text{user function} \\
  \bnfor& \tmfunK{x : X}{K}                       & &\text{kernel function} \\
  \bnfor& \tmrunner{(\tm{op}\,x \mapsto K_{\tm{op}})_{\tm{op} \in \sig}}{C}
                                                  & &\text{runner}
  \\[1ex]
  \intertext{\textbf{User computations}}
  M, N
  \bnfis& \tmreturn{V}                            & &\text{value} \\
  \bnfor& V\,W                                    & &\text{application} \\
  \bnfor& \tmtry{M}{
          \{ \tmreturn{x} \mapsto N,
             (\tmraise{e} \mapsto N_e)_{e \in E} \}
          }
                                                  & &\text{exception handler} \\
  \bnfor& \tmmatch{V}{\tmpair{x}{y} \mapsto M}    & &\text{product elimination} \\
  \bnfor& \tmmatch[X]{V}{}                        & &\text{empty elimination} \\
  \bnfor& \tmmatch{V}{\tminl{x} \mapsto M, \tminr{y} \mapsto N}
                                                  & &\text{sum elimination} \\
  \bnfor& \tmop{op}{X}{V}{\tmcont x M}{\tmexccont N e {E_\op}}
                                                  & &\text{operation call} \\
  \bnfor& \tmraise[X]{e}                          & &\text{raise exception} \\
  \bnfor& \tmrun{V}{W}{M}{F} 
                                                  & &\text{running user code} \\
  \bnfor& 
            \tmkernel{K}{W}{F}
                                                  & &\text{switch to kernel mode}
  \\[2ex]
  F \bnfis & \omit \rlap{$\{ \tmreturn{x} \at c \mapsto N, 
                    (\tmraise{e} \at c \mapsto N_e)_{e \in E},
                    (\tmkill{s} \mapsto N_s)_{s \in S} \}$}
  \\[1ex]
  \intertext{\textbf{Kernel computations}}
  K, L
  \bnfis& \tmreturn[C]{V}                         & &\text{value} \\
  \bnfor& V\,W                                    & &\text{application} \\
  \bnfor& \tmtry{K}{
          \{ \tmreturn{x} \mapsto L,
             (\tmraise{e} \mapsto L_e)_{e \in E} \}
          }
                                                  & &\text{exception handler} \\
  \bnfor& \tmmatch{V}{\tmpair{x}{y} \mapsto K}    & &\text{product elimination} \\
  \bnfor& \tmmatch[X \at C]{V}{}                  & &\text{empty elimination} \\
  \bnfor& \tmmatch{V}{\tminl{x} \mapsto K, \tminr{y} \mapsto L}
                                                  & &\text{sum elimination} \\
  \bnfor& \tmop{op}{X}{V}{\tmcont x K}{\tmexccont L e {E_\op}}
                                                  & &\text{operation call} \\
  \bnfor& \tmraise[X \at C]{e}                    & &\text{raise exception} \\
  \bnfor& \tmkill[X \at C]{s}                     & &\text{send signal} \\
  \bnfor& \tmgetenv[C]{\tmcont c K}               & &\text{get kernel state} \\
  \bnfor& \tmsetenv{V}{K}                         & &\text{set kernel state} \\
  \bnfor& \tmuser{M}{
          \begin{aligned}[t]
          \{ &\tmreturn{x} \mapsto K,
             (\tmraise{e} \mapsto L_e)_{e \in E} \}
          \end{aligned}
          }
                                                  & &\text{switch to user mode}
  \end{align*}
  } 
  \caption{Values, user computations, and kernel computations of $\lambdacoop$.}
  \label{fig:lambdacoop-terms}
\end{figure}

The syntax of terms is shown in \cref{fig:lambdacoop-terms}. The
usual fine-grain call-by-value stratification of terms into pure values and effectful
computations is present, except that we further distinguish between \emph{user} and
\emph{kernel} computations.

\subsubsection{Values}
\label{sec:values}

Among the values are variables, constants for ground types, and constructors 
for sums and products. There are two kinds of functions, for abstracting over user and
kernel computations. A \emph{runner} is a value of the form
%
\begin{equation*}
  \tmrunner{(\tm{op}\,x \mapsto K_{\tm{op}})_{\tm{op} \in \sig}}{C}.
\end{equation*}
%
It implements co-operations for operations $\tm{\op}$ as kernel
computations~$K_\op$, with $x$ bound in~$K_\op$. The type annotation~$C$
specifies the type of the state that~$K_\op$ uses.
Note that $C$ ranges over ground types, a restriction that allows us to define a naive
set-theoretic semantics.
%
We sometimes omit these type annotations.

\subsubsection{User and kernel computations}

The user and kernel computations both have pure computations, function application,
exception raising and handling, standard elimination forms, and operation calls.
Note that the typing annotations on some of these differ according to their mode.
For instance, a user operation call is annotated
with the result type~$X$, whereas the annotation $X \at C$ on a kernel operation call
also specifies the kernel state type~$C$.

The binding construct $\tmlet[X ! E]{x}{M}{N}$ is not part of the syntax,
but is an abbreviation for
%
$
  \tmtry{M}{
  \{ \tmreturn{x} \mapsto N,
     (\tmraise{e} \mapsto \tmraise[X]{e})_{e \in E}
   \}}
$,
%
and there is an analogous one for kernel computations. We often drop the 
annotation $X ! E$.

Some computations are specific to one or the other mode. Only the kernel mode
may send a signal with $\tmkill{\!}$, and manipulate state with
$\tmkw{getenv}$ and $\tmkw{setenv}$, but only the user mode has the 
$\tmkw{run}$ construct from \cref{sec:runn-as-progr}.
%
Finally, each mode has the ability to ``context switch'' to the other one.
The kernel computation
%
\begin{equation*}
\tmuser{M}{
   \{
     \tmreturn{x} \mapsto K,
     (\tmraise{e} \mapsto L_e)_{e \in E}
   \}
}
\end{equation*}
%
runs a user computation $M$ and handles the returned value and leftover
exceptions with kernel computations $K$ and $L_e$.
Conversely, the user computation
%
\begin{equation*}
\tmkernel{K}{W}{
  \{x \at c \mapsto M,
    (\tmraise{e} \at c \mapsto N_e)_{e \in E},
    (\tmkill{s} \mapsto N_s)_{s \in S}
  \}
}
\end{equation*}
%
runs kernel computation $K$ with initial state $W$, and handles the returned value,
and leftover exceptions and signals with user computations $M$, $N_e$, and $N_s$.

\subsection{Type system}
\label{sec:typesystem}

We equip $\lambdacoop$ with a type system akin to type and effect systems for
algebraic effects and handlers~\cite{Bauer:EffectSystem,Benton:ExceptionalSyntax,Kammar:Handlers}.
We are experimenting with resource control, so it makes sense for the type system
to tightly control resources. Consequently, our effect system
does not allow effects to be implicitly propagated outwards.

In \cref{sect:types}, we assumed that each operation $\op \in \Ops$ 
is equipped with some fixed operation signature
%
$
  \op : \tysigop{A_\op}{B_\op}{E_\op}
$.
%
We also assumed a fixed constant signature $\tmconst{f} : (A_1, \ldots, A_n) \to B$
for each ground constant $\tmconst{f}$.
%
We consider this information to be part of the type system and say no more about it.

Values, user computations, and kernel computations each have a corresponding
\emph{typing judgement} form and a \emph{subtyping relation}, given by 
%
\begin{align*}
  &\Gamma \types V : X,
& &\Gamma \types M : \tyuser{X}{\Ueff},
& &\Gamma \types K : \tykernel{X}{\Keff},\\
  &X \sub Y,
& &\tyuser{X}{\Ueff} \sub \tyuser{Y}{\Veff},
& &\tykernel{X}{\Keff} \sub \tykernel{Y}{\Leff}, 
\end{align*}
%
where $\Gamma$ is a \emph{typing context} $x_1 : X_1, \ldots, x_n : X_n$.
%
The effect information is an over-approximation, i.e., $M$ and $K$ employ \emph{at
  most} the effects described by $\Ueff$ and $\Keff$.
%
The complete rules for these judgements are given in the online appendix. %\cref{sec:typing-rules}. 
We comment here
only on the rules that are peculiar to~$\lambdacoop$, see \cref{fig:typing-selected}.

\begin{figure}[tp]
  \centering
  \small
  \begin{mathpar}
    %%% NOTE: these should be kept in sync with the rules given in
    %%%       tpying-rules.tex. The rules in that file are the official version,
    %%%       the ones here are the copies. If you delete a rule here, make sure
    %%%       to note so in typing-rules.tex.
    %%%
    %%%       Do not place a rule here without commenting on it in the text,
    %%%       that's bad manners.
    \coopinfer{Sub-Ground}{ }{A \sub A}
    
    \coopinfer{Sub-Runner}{
      \sig_1' \subseteq \sig_1 \\
      \sig_2 \subseteq \sig_2' \\
      S \subseteq S' \\
      C \equiv C'
    }{
      \tyrunner{\sig_1}{\sig_2}{S}{C} \sub \tyrunner{\sig_1'}{\sig_2'}{S'}{C'}
    }

    \coopinfer{Sub-Kernel}{
      X \sub X' \\
      \sig \subseteq \sig' \\
      E \subseteq E' \\
      S \subseteq S' \\
      C \equiv C'
    }{
      \tykernel{X}{(\sig, E, S, C)} \sub \tykernel{X'}{(\sig', E', S', C')}
    }

  \coopinfer{TyUser-Try}{
    \Gamma \types M : \tyuser{X}{(\sig,E)}
    \\
    \Gamma, x \of X \types N : \tyuser{Y}{(\sig,E')}
    \\
    \big(
      \Gamma \types N_e : \tyuser{Y}{(\sig,E')}
    \big)_{e \in E}
  }{
    \Gamma \types
    \tmtry{M}{
        \{ \tmreturn{x} \mapsto N,
           (\tmraise{e} \mapsto N_e)_{e \in E} \}
        }
    : \tyuser{Y}{(\sig,E')}
  }

  \coopinfer{TyUser-Run}{
    F \equiv
    \{ \tmreturn{x} \at c \mapsto N,
       (\tmraise{e} \at c \mapsto N_e)_{e \in E},
       (\tmkill{s} \mapsto N_s)_{s \in S}
    \}
    \\\\
    \Gamma \types V : \tyrunner{\sig}{\sig'}{S}{C} \\
    \Gamma \types W : C \\\\
    \Gamma \types M : \tyuser{X}{(\sig, E)} \\
    \Gamma, x \of X, c \of C \types N : \tyuser{Y}{(\sig', E')} \\
    \big(
       \Gamma, c \of C \types N_e : \tyuser{Y}{(\sig', E')}
    \big)_{e \in E} \\
    \big(
       \Gamma \types N_s : \tyuser{Y}{(\sig', E')}
    \big)_{s \in S} \\
  }{
    \Gamma \types \tmrun{V}{W}{M}{F} : \tyuser{Y}{(\sig', E')}
  }

  \coopinfer{TyUser-Op}{
    \Ueff \equiv (\sig,E) \\
    \op \in \sig \\
    \Gamma \types V : A_\op \\\\
    \Gamma, x \of B_\op \types M : \tyuser{X}{\Ueff} \\
    \big(
      \Gamma \vdash N_e : \tyuser{X}{\Ueff}
    \big)_{e \in E_\op}
  }{
    \Gamma \types \tmop{op}{X}{V}{\tmcont x M}{\tmexccont N e {E_\op}} : \tyuser{X}{\Ueff}
  }

  \coopinfer{TyKernel-Op}{
    \Keff \equiv (\sig, E, S, C) \\
    \op \in \sig \\
    \Gamma \types V : A_\op \\\\
    \Gamma, x \of B_\op \types K : \tykernel{X}{\Keff} \\
    \big(
      \Gamma \vdash L_e : \tykernel{X}{\Keff}
    \big)_{e \in E_\op}
  }{
    \Gamma \types \tmop{op}{X}{V}{\tmcont x K}{\tmexccont L e {E_\op}} : \tykernel{X}{\Keff}
  }

  \coopinfer{TyUser-Kernel}{
    F \equiv
    \{ \tmreturn{x} \at c \mapsto N,
       (\tmraise{e} \at c \mapsto N_e)_{e \in E},
       (\tmkill{s} \mapsto N_s)_{s \in S}
    \}
    \\\\
    \Gamma \types K : \tykernel{X}{(\sig, E, S, C)} \\
    \Gamma \types W : C \\
    \Gamma, x \of X, c \of C \types N : \tyuser{Y}{(\sig, E')} \\
    \big(
      \Gamma, c \of C \types N_e : \tyuser{Y}{(\sig, E')}
    \big)_{e \in E} \\
    \big(
      \Gamma \types N_s : \tyuser{Y}{(\sig, E')}
    \big)_{s \in S} \\
  }{
    \Gamma \types \tmkernel{K}{W}{F} : \tyuser{Y}{(\sig, E')}
  }

  \coopinfer{TyKernel-User}{
   \Keff \equiv (\sig, E', S, C) \\
   \Gamma \types M : \tyuser{X}{(\sig, E)} \\\\
   \Gamma, x \of X \types K : \tykernel{Y}{\Keff} \\
   \big(
     \Gamma \types L_e : \tykernel{Y}{\Keff}
   \big)_{e \in E}
  }{
    \Gamma \types
    \tmuser{M}{
      \{ \tmreturn{x} \mapsto K,
         (\tmraise{e} \mapsto L_e)_{e \in E}
      \}
    }
    : \tykernel{Y}{\Keff}
  }
  \end{mathpar}
  \caption{Selected typing and subtyping rules.}
  \label{fig:typing-selected}
\end{figure}

Subtyping of ground types \textsc{Sub-Ground} is trivial, as it relates only equal types.
Subtyping of runners \textsc{Sub-Runner} and kernel computations
\textsc{Sub-Kernel} requires equality of the kernel state types~$C$ and~$C'$
because state is used invariantly in the kernel monad.
We leave it for future work to replace ${C \equiv C'}$ with a
\emph{lens}~\cite{Foster:Lenses} from~$C'$ to~$C$, i.e., maps $C' \to C$ and ${C' \times C \to C'}$
satisfying state equations analogous to \cref{ex:state}. It has been
observed~\cite{OConnor:Lens,Power:Comodels} that such a lens in fact amounts to
an ordinary runner for $C$-valued state.

The rules \textsc{TyUser-Op} and \textsc{TyKernel-Op} govern operation calls, where 
we have a success continuation which receives a value returned by a 
co-operation, and exceptional continuations which receive exceptions raised by co-operations.

The rule \textsc{TyUser-Run} requires that the runner $V$ implements \emph{all} the
operations $M$ can use, meaning that operations are \emph{not} implicitly propagated outside a $\tmkw{run}$ block (which is different from how handlers are sometimes implemented). Of course, the co-operations of the runner may call further external operations, as recorded by the signature~$\sig'$. Similarly, we require the finally block~$F$ to intercept all exceptions and signals that might be produced by the co-operations of $V$ or the user code $M$.
%
Such strict control is exercised throughout. For example, in 
\textsc{TyUser-Run}, \textsc{TyUser-Kernel}, and \textsc{TyKernel-User} we catch all the exceptions and signals that the code might produce.
%
One should judiciously relax these requirements in a language that is presented to
the programmer, and allow re-raising and re-sending clauses to be automatically inserted.

%%% Local Variables:
%%% mode: latex
%%% TeX-master: "runners-in-action"
%%% End:

% !TEX root = runners-in-action.tex

\subsection{Equational theory}
\label{sect:eqtheory}

We present $\lambdacoop$ as an \emph{equational calculus}, i.e., the interactions between
its components are described by equations. Such a presentation makes it easy to reason
about program equivalence.
%
There are three equality judgements
%
\begin{equation*}
\Gamma \types V \equiv W : X, 
\qquad
\Gamma \types M \equiv N : \tyuser{X}{\Ueff}, 
\qquad
\Gamma \types K \equiv L  : \tyuser{X}{\Keff}.
\end{equation*}
%
It is presupposed that we only compare well-typed expressions with the indicated types.
For the most part, the context and the type annotation on judgements will play no significant role,
and so we shall drop them whenever possible.

We comment on the computational equations for constructs characteristic
of~$\lambdacoop$, and refer the reader to the online appendix for other equations.
%
When read left-to-right, these equations explain the operational meaning of programs.

Of the three equations for $\tmkw{run}$, the first two specify that returned values and
raised exceptions are handled by the corresponding clauses,
%
\begin{align*}
  \tmrun{V}{W}{(\tmreturn{V'})}{F} &\equiv N[V'/x, W/c], 
  \\
  \tmrun{V}{W}{(\tmraise[X]{e})}{F} &\equiv N_{e}[W/c],
\end{align*}
%
where
%
$F \defeq \{\tmreturn{x} \at c \mapsto N,
   (\tmraise{e} \at c \mapsto N_e)_{e \in E},
   (\tmkill{s} \mapsto N_s)_{s \in S}
\}$.
%
The third equation below relates running an operation $\op$ with executing the corresponding co-operation~$K_\op$, 
where $R$ stands for the runner
%
$\tmrunner{(\tm{op}\,x \mapsto K_{\tm{op}})_{\tm{op} \in \sig}}{C}$:
%
\begin{multline*}
  \tmrun{R}{W}{(
    \tmop{op}{X}{V}{\tmcont x M}{\tmexccont {N'} {e'} {E_\op}}
    )}{F} \equiv {}
  \\
  \begin{aligned}[t]
     &\tmkernel{K_\op[V/x]}{W}{} \\
     &\qquad\big\{
         \begin{aligned}[t]
           &\tmreturn{x} \at c' \mapsto (\tmrun{R}{c'}{M}{F}), \\
           &\left(
               \tmraise{e'} \at c' \mapsto (\tmrun{R}{c'}{N'_{e'}}{F})
             \right)_{e' \in E_\op},\\
           &\left(
               \tmkill{s} \mapsto N_s
             \right)_{s \in S} \big\}
         \end{aligned}
 \end{aligned}
\end{multline*}
%
Because $K_\op$ is kernel code, it is executed in kernel mode, whose
$\tmkw{finally}$ clauses specify what happens afterwards: if $K_\op$ returns a value, or
raises an exception, execution continues with a suitable continuation, with~$R$
wrapped around it; and if $K_\op$ sends a signal, the corresponding finalisation code from $F$ is
evaluated.

The next bundle describes how kernel code is executed within user code:
%
\begin{align*}
  \tmkernel{(\tmreturn[C]{V})}{W}{F} &\equiv N[V/x, W/c], \\
  \tmkernel{(\tmraise[X \at C]{e})}{W}{F} &\equiv N_{e}[W/c], \\
  \tmkernel{(\tmkill[X \at C]{s})}{W}{F} &\equiv N_{s}, \\
  \tmkernel{(\tmgetenv[C]{\tmcont c K})}{W}{F} &\equiv \tmkernel{K[W/c]}{W}{F}, \\
  \tmkernel{(\tmsetenv{V}{K})}{W}{F} &\equiv \tmkernel{K}{V}{F}.
\end{align*}
%
We also have an equation stating that an operation called in kernel mode propagates out to
user mode, with its continuations wrapped in kernel mode:
%
\begin{multline*}
  \tmkernel{\tmop{op}{X}{V}{\tmcont x K}{\tmexccont L {e'} E}}{W}{F}
  \equiv {} \\
  \tmop{op}{X}{V}{\tmcont x {\tmkernel{K}{W}{F}}}{
    \left(
      \tmkernel{L_{e'}}{W}{F}
    \right)_{e' \in E}}.
\end{multline*}
%
Similar equations govern execution of user computations in kernel mode.

The remaining equations include standard $\beta\eta$-equations for
exception handling~\cite{Benton:ExceptionalSyntax}, deconstruction of products and sums,
algebraicity equations for operations~\cite{Pretnar:Thesis}, and the equations of kernel theory from \cref{sec:user-kernel-monads}, describing how $\tmkw{getenv}$ and $\tmkw{setenv}$ work, and how they interact with signals and other operations.

%%% Local Variables:
%%% mode: latex
%%% TeX-master: "runners-in-action"
%%% End:

\section{Denotational semantics}
\label{sec:denotation}

In this section we assign mathematical meaning to the constituent parts of Clerical.
%
The types are interpreted by  the expected sets:
%
\begin{align*}
\sem{\dZ} &\defeq \IZ &
\sem{\dB} &\defeq \{\semff, \semtt\} &
\sem{\dR} &\defeq \IR &
\sem{\dU} &\defeq \{\star\} \enspace .
\end{align*}
%
Typing contexts are interpreted by cartesian product:
%
\begin{align*}
  \sem{x_1 \of \tau_1, \ldots, x_m \of \tau_m} &\defeq
  \sem{\tau_1} \times \cdots \times \sem{\tau_m} \enspace .
%  &
%  \sem{\Gamma; \Delta} = \sem{\Gamma} \times \sem{\Delta}
\end{align*}
%
The denotation $\sem{\Gamma}$ of a read-only context $\Gamma$ is thought of as an \defemph{environment} specifying values of variables, whereas the denotation of a read-write context $\Gamma ; \Delta$ has two components, the environment $\sem{\Gamma}$ and the \defemph{state} $\sem{\Delta}$.

The meaning of a well-typed pure expression $\Gamma \rotypes e : \tau$ and
general expression $\Gamma; \Delta \rwtypes c : \tau$ 
 will be maps:
%
\begin{align*}
  \sem{\Gamma  \rotypes e : \tau}  & : \sem{\Gamma} \to \PP{\sem{\tau}}, \\
  \sem{\Gamma; \Delta \rwtypes c : \tau}  & : \sem{\Gamma} \to (\sem{\Delta} \to \PP{\sem{\Delta} \times \sem{\tau}}),
\end{align*}
%
where $\PP{S}$ is  a \defemph{powerdomain}, a collection of sets representing the possible sets of outcomes of nondeterministic computations that return values from $S$. Due to the distinction between error and non-termination, already mentioned in Section~\ref{sec:overview}, 
Clerical requires a rather specific form of powerdomain. In order to motivate it, we discuss the distinction between
error and non-termination in more detail.

Mathematically, we would like to  consider a well-behaved (deterministic) expression $e$ of type $\dR$ as defining a real number value $r$. Computationally, however, the best that $e$ will be able to do is to produce, on demand, an approximation of $r$ to within any specified precision. In the ideal case, given precision $\epsilon > 0$, the evaluation of $e$ will determine in finite time some rational approximation $q$ such that $|q - r| < \epsilon$. In Clerical, not every  deterministic expression of type $\dR$ achieves this ideal.

Some expressions simply give rise to non-terminating computations that never provide any approximating information. Others, may appear to provide approximating information, but do so in a way that is either incomplete or inconsistent.\footnote{Incompleteness may  arise, for example, if approximating values are only computed for $\epsilon$ that are not too small. Similarly, inconsistency can occur if two different $\epsilon_1, \epsilon_2$ result in putative approximations $q_1,q_2$ with $|q_1 - q_2| \geq \epsilon_1 + \epsilon_2$.} Our semantics of Clerical distinguishes between these two eventualities. Expressions that produce incomplete or inconsistent approximating information are considered \defemph{erroneous}, and the semantics for Clerical will ensure that no such expression is ever executed within a program with valid semantics. The motivation for this is to avoid any situation in which faulty approximations can provide misleading information. In contrast, \defemph{non-terminating} expressions are considered harmless in the sense that they cannot be a source of incorrect information. As is standard in denotational semantics, such expressions are assigned the special denotation $\bot$. 
These need to be distinguished from erroneous ones, because, as discussed in Section~\ref{sec:overview},  non-terminating expressions 
have an essential role to play  when programming in Clerical.


For any set~$S$,  define $\liftnoerr{S} \defeq S + \{\bot\}$, where~$\bot$ represents non-termination, as discussed above. 
Although $S + \{\bot\}$ is strictly speaking a coproduct (sum) of two sets, in practice we shall only use it
in instances in which $\bot \notin S$. This allows us to represent $\liftnoerr{S}$ as the (disjoint) union 
$S \cup \{\bot\}$.
Define:
\begin{equation*}
  \PP{S} \defeq
  \{ X \subseteq \liftnoerr{S} \such
       \text{$X$ infinite} \Rightarrow \bot \in X
  \},
\end{equation*}
%
A set $X \in  \PP{S}$ represents the results of a nondeterministic Clerical computation in the following way. Firstly, $X = \emptyset$ in the case that the computation is \emph{erroneous} (see the discussion above), in which case no result value is relevant. 
%
The case $\bot \in X$ applies if the nondeterministic computation has at least one non-terminating branch, in which case $X \setminus \{\bot\}$ is the set of all result values returned by terminating nondeterministic branches. If instead $\bot \notin X$ then $X$ represents the set of possible results of a necessarily terminating nondeterministic computation. Since Clerical has only  finite nondeterministic branching, such a set is necessarily finite.

We shall implicitly make use of  the fact that powerdomain $\Pstar$ carries the structure of a monad on the category of sets. For the purposes of this paper, we never need to directly refer to the abstract structure of the monad. However, the following maps associated with this structure will be useful. Firstly, for any $S$, there is a map $x \mapsto \pure{x} : S \to \PP{S}$, that maps any $x \in S$ to the singleton $\pure{x}$ representing the deterministic computation that returns $x$ as its result. 
Secondly, for any function $f : S \to  \PP{T}$ we define $\lift{f} : \PP{S} \to  \PP{T}$ by:
\[
\lift{f}(X) = 
\begin{cases} 
\emptyset & \text{if $\some{x \in X}\ f(x)= \emptyset$,}\\
\textstyle
\{\bot \such \bot \in X\} \cup \bigcup_{x \in X \setminus \{\bot\}} f(x) & \text{otherwise.}
\end{cases}
\]
It is easily checked that this indeed defines a set in $\PP{T}$. The idea behind the definition is that 
$\lift{f}(X)$ models a sequencing of nondeterministic computations: first execute the nondeterministic computation whose result is represented by $X$, then for each potential value $x \in X$
run the nondeterministic computation modelled by $f(x)$ to obtain potential return values.
This idea motivates the alternative notation
\[
\PPlet{x}{X}\, f(x)
\]
for  $\lift{f}(X)$, which we shall often use.

The reason behind the first clause in the definition of 
$\lift{f}(X)$ is that a computation is considered illegitimate if an error occurs along any possible nondeterministic branch. In such a case the entire computation is given the error denotation $\emptyset$.

\begin{figure}
  \begin{mdframed}
  \centering
  \small
\begin{align*}
  \sem{\Gamma \rotypes c : \tau} \, \gamma = {}&
     \PPlet
        {(v, ())}
        {\sem{\Gamma; \emptyctx \rwtypes c : \tau} \, \gamma \, ()}
        {\pure{v}}
  \\
  \sem{x_1 \of \tau_1; \ldots, x_n \of \tau_n \rotypes x_i : \tau_i} \, \gamma
  = {}& \pure{\gamma_i}
  \\
  \sem{\Gamma \rotypes \cfalse : \dB} \, \gamma = {}& \pure{\semff} \\
  \sem{\Gamma \rotypes \ctrue : \dB} \, \gamma = {}& \pure{\semtt} \\
  \sem{\Gamma \rotypes \numeral{k} : \dZ} \, \gamma = {}& \pure{k} \\
  \sem{\Gamma \rotypes \cskip : \dU} \, \gamma = {}& \pure{\semuu} \\
  \sem{\Gamma \rotypes \ccoerce{e} : \dR} \, \gamma = {}& \PPlet
        {z}
        {\sem{\Gamma \rotypes e : \dZ}\,\gamma}
        {\pure{\inclZ{z}}}
  \\
  \sem{\Gamma \rotypes e_1 \iop e_2 : \dZ} \, \gamma = {}&
    \PPlet
    {x}
    {\sem{\Gamma \rotypes e_1 : \dZ}\,\gamma} 
 \\
 & \PPlet
    {y}
    {\sem{\Gamma \rotypes e_2 : \dZ}\,\gamma}
    \pure{x \op y}
 \\
 \sem{\Gamma \rotypes e_1 \rop e_2 : \dR} \, \gamma = {}& 
   \PPlet
    {x}
    {\sem{\Gamma \rotypes e_1 : \dR}\,\gamma}
\\
& \PPlet
    {y}
    {\sem{\Gamma \rotypes e_2 : \dR}\,\gamma}
    \pure{x \op y}
 \\
 \sem{\Gamma \rotypes \cinv{e} : \dR} \, \gamma = {}&
 \PPlet
    {x}
    {\sem{\Gamma \rotypes e : \dR}\,\gamma}
    \begin{cases} 
    \pure{x^{-1}} & \text{if $x \neq 0$} \\
    \PPbot  & \text{if $x=0$}
    \end{cases}
 \\
  \sem{\Gamma \rotypes e_1 = e_2 : \dZ} \, \gamma = {}&
    \PPlet
    {x}
    {\sem{\Gamma \rotypes e_1 : \dZ}\,\gamma}
\\
& 
\PPlet
    {y}
    {\sem{\Gamma \rotypes e_2 : \dZ}\,\gamma}
    \begin{cases} 
    \pure{\semtt} & \text{if $x = y$} \\
    \pure{\semff} & \text{if $x \neq y$}
    \end{cases}
\\
  \sem{\Gamma \rotypes e_1 < e_2 : \dZ} \, \gamma = {}&
    \PPlet
    {x}
    {\sem{\Gamma \rotypes e_1 : \dZ}\,\gamma}
\\
& \PPlet
    {y}
    {\sem{\Gamma \rotypes e_2 : \dZ}\,\gamma}
    \begin{cases} 
    \pure{\semtt} & \text{if $x < y$} \\
    \pure{\semff} & \text{if $x \geq y$}
    \end{cases}
\\
   \sem{\Gamma \rotypes e_1 \rlt e_2 : \dR} \, \gamma = {}&
    \PPlet
    {x}
    {\sem{\Gamma \rotypes e_1 : \dR}\,\gamma} \\
    &\PPlet
    {y}
    {\sem{\Gamma \rotypes e_2 : \dR}\,\gamma}
    \begin{cases}
    \pure{\semtt} & \text{if $x < y$} \\
    \pure{\semff} & \text{if $x > y$} \\
    \PPbot &  \text{if $x = y$} 
    \end{cases}
  \\
  \sem{\Gamma \rotypes (\clim{x}{e}) : \dR}\,\gamma = {}&
  \begin{cases}
    \pure{t} &
      \begin{aligned}[t]
      &\text{if $t \in \IR$ and $\all{k \in \IZ}$} \\
      &\text{  $\sem{\Gamma, x \of \dZ \rotypes e : \dR} (\gamma, k) \subseteq \IR$ and} \\
      &\text{  $\all{u \in \sem{\Gamma, x \of \dZ \rotypes e : \dR} (\gamma, k)} |u - t| < 2^{-k}$,}
      \end{aligned}
    \\
    \PPerr   & \text{if no such $t \in \IR$ exists.}
  \end{cases}
\end{align*}
\end{mdframed}
\caption{Denotational semantics of pure expressions}
\label{figure:ro-denotations}
\end{figure}


\begin{figure}
  \begin{mdframed}
  \centering
  \small
\begin{align*}
\sem{\Gamma; \Delta \rwtypes e : \tau} \, \gamma \, \delta &=
  \begin{aligned}[t]
   &\PPlet
     {v}
     {\sem{\Gamma, \Delta \rotypes e : \tau} \, (\gamma, \delta)} \\
             &
        \quad
     \pure{(\delta, v)}
    \end{aligned}
\\
\sem{\Gamma; \Delta \rwtypes c_1 ; c_2 : \tau} \, \gamma \, \delta &=
   \begin{aligned}[t]
   &\PPlet
     {(\delta', \semuu)}
     {\sem{\Gamma; \Delta \rwtypes c_1 : \dU} \, \gamma \, \delta} 
   \\
   &  
     \quad\sem{\Gamma; \Delta \rwtypes c_2 : \tau} \, \gamma \, \delta'
       \end{aligned}
\\
\sem{\Gamma ; \Delta \rwtypes (\cvar{x}{e} c) : \tau} \, \gamma \, \delta &=
  \begin{aligned}[t]
   &\PPlet
     {v}
     {\sem{\Gamma, \Delta \rotypes e : \sigma} \, (\gamma, \delta)}
   \\
   &
    \PPletx
     {((\delta', v'), v'')}
     \\
     &
     \qquad\PPinx{\sem{\Gamma; \Delta, x \of \sigma \rwtypes c : \tau} \, \gamma \, (\delta,v)}
   \\
    &
     \quad\pure{(\delta', v'')}
  \end{aligned}
\\
\sem{\Gamma ; \Delta \rwtypes (\clet{x}{e}) : \dU} \, \gamma \, \delta &=
  \begin{aligned}[t]
  &\PPlet
    {v}
    {\sem{\Gamma, \Delta \rotypes e : \tau} \, (\gamma, \delta)}
        \\
        &
        \quad
    \pure{(\delta[x :=v]), \semuu)}
    \end{aligned}
\\
\sem{\Gamma ; \Delta \rwtypes (\cif e \cthen c_1 \celse c_2 \cend) : \tau} \, \gamma \, \delta &=
  \begin{aligned}[t]
  &\PPlet
    {b}
    {\sem{\Gamma, \Delta \rotypes e : \dU} \, (\gamma, \delta)}
    \\
    & \quad \begin{cases}
    \sem{\Gamma; \Delta \rwtypes c_1 : \tau} \, \gamma \, \delta & \text{if $b = \semtt$} \\
     \sem{\Gamma; \Delta \rwtypes c_2 : \tau} \, \gamma \, \delta & \text{if $b = \semff$}
     \end{cases}
  \end{aligned}
\end{align*}
%
\end{mdframed}
\caption{Denotational semantics of general expressions, excluding $\mathtt{case}$ and $\mathtt{while}$}
\label{figure:rw-denotations}
\end{figure}

\Cref{figure:ro-denotations} assigns denotational semantics to pure expressions.
One point that deserves explanation is the denotation of $\cinv{e}$, when~$e$ is a real expression representing~$0$. Since there is no appropriate real value to be given, the denotation could be chosen to be 
either $\PPerr$ or $\PPbot$. We choose the latter, as it reflects the fact that an algorithm for calculating reciprocal will run forever, given a representation of the real number~$0$, without ever returning any erroneous approximation to a result value.
Similarly, $\PPbot$ is given as the denotation of  $e_1 \rlt e_2$ when $e_1,e_2$ are two real expressions representing equal numbers, reflecting the fact that an algorithm trying to distinguish between the two numbers will run forever when given equal inputs.

The most complex definition in \cref{figure:ro-denotations} is the semantics of the  limit operation $\clim{x}{e}$.
%
In this definition, note that there is at most one $t \in \IR$ satisfying the first condition.
Its existence places strong requirements on the expression $e$, which must represent a sequence of sets of real numbers, such that every real number~$u$ in the $k$-th set lies within a distance of $2^{-k}$ of~$t$.
That is, every choice of a real number from every set furnishes a Cauchy sequence rapidly converging to a common limit~$t$.
The use of a sequence of sets allows the behaviour of~$e$ to be nondeterministic, but this nondeterminism is highly constrained by the common limit requirement. Furthermore, $e$ is neither allowed to diverge nor be erroneous. If any of the conditions required for~$t$ to exist fails, then the $\clim{x}{e}$ computation is declared erroneous.
This is appropriate because, in an algorithmic implementation of the limit operation, the source of error may occur deep in the computation (e.g., only at some high value for the integer~$k$) meaning that the algorithm may, before the error transpires, return erroneous information in the form of approximating values to a non-existent limit.

\Cref{figure:rw-denotations} assigns semantics to
several of the general expression  constructors: sequencing $c_1 ; c_2$,
local variable declarations $\cvar{x}{e} c$,
assignments $\clet{x_{i}}{e}$,
conditionals $\cif e \cthen c_1 \celse c_2 \cend$, and expressions $e$ qua commands.

Next, let us define the semantics of guarded choice
%
\begin{equation*}
  \Gamma ; \Delta \rwtypes (\ccase e_1 \To c_1 \such \cdots \such e_n \To c_n \cend) : \tau.
\end{equation*}
%
Using the abbreviations
%
$\sem{e_i} = \sem{\Gamma, \Delta \rotypes e_i : \dB}$
and
$\sem{c_i} = \sem{\Gamma; \Delta \rwtypes c_i : \tau}$
%
we set
%
\begin{align*}
  &\sem{\Gamma ; \Delta \rwtypes (\ccase e_1 \To c_1 \such \cdots \such e_n \To c_n \cend) : \tau} \, \gamma \, \delta = 
  \\
  &\quad
  \begin{cases}
    \emptyset \qquad\text{if $\some{i} \sem{e_i} \, (\gamma, \delta) = \emptyset \lor (\semtt \in \sem{e_i} \, (\gamma, \delta) \land \sem{c_i} \, \gamma \, \delta = \emptyset)$}
    \\
    S
    \cup \{ \bot \such \all{i} \sem{e_i} (\gamma, \delta) \neq \pure{\semtt} \}
    \qquad\text{otherwise}
  \end{cases}
  \\
  &\quad\text{where $S = \bigcup \left\{\sem{c_i} \, \gamma \, \delta \such
                 1 \leq i \leq n \land
                \semtt \in \sem{e_i} (\gamma, \delta) \right\}.$}
\end{align*}
The idea behind this definition is as follows. 
The $n$ (potentially nondeterministic) guard expressions $e_1, \dots, e_n$ are evaluated in parallel. 
Ignoring, for the moment, the possibility that one of these expressions might be erroneous, suppose
that one of them, 
$e_i$ say, evaluates to $\semtt$. If this occurs, then the parallel evaluation of the other guards is terminated and the continuation
$c_i$ is executed. Note that the choice of $i$ here is potentially nondeterministic.
If instead none of the guards evaluates to  $\semtt$, then none of the continuations is triggered and we consider this as amounting to nontermination (no error is caused), so we include 
$\bot$ in the denotation of the case expression. Note that this $\bot$ can arise as a result of \emph{bona fide} nontermination
(none of the guards terminates) or of deadlock (all guards terminate with $\semff$).
In the case that any of the $e_i$ guards causes an error (i.e., has denotation $\emptyset$), or if some $e_i$ has a possible evaluation to $\semtt$ that triggers the execution of an erroneous continuation command $c_i$, then the case statement is itself given the error denotation $\PPerr$.
Some of the subtleties of guarded choice are illustrated through the examples in
Section~\ref{sec:boolean-ops} below. 



It remains to define the semantics of while loops.
As usual, the meaning of a while loop is required to be invariant under unrolling; i.e.,
\begin{align*}
& \sem{\Gamma;\Delta \rwtypes \cwhile e \cdo c \cend :\dU}\,\gamma\,\delta  = 
\\
& \qquad \sem{\Gamma;\Delta \rwtypes\cif e \cthen (c\, ; \, \cwhile e \cdo c \cend) \celse \cskip \cend :\dU}\,\gamma\,\delta
\end{align*}
That is, we want the value $\sem{\Gamma ; \Delta \rwtypes (\cwhile e \cdo c \cend) : \dU} \, \gamma $ to be a fixed point of the map 
%
\begin{equation}
  \label{eq:def-W}
  \begin{aligned}
    W_{\gamma} &:  \PP{\sem{\Delta} \times \{\semuu\}}^{\sem{\Delta}}
    \to \PP{\sem{\Delta} \times \{\semuu\}}^{\sem{\Delta}}
    \\
    W_\gamma&(h) \, \delta 
    \defeq \PPlet{b}{\sem{e}\,(\gamma,\delta)}
      \begin{cases}
        \lift{(h \circ \pi_1)}(\sem{\Gamma; \Delta \rwtypes c : \dU}\,\gamma\,\delta) & \text{if $b = \semtt$} \\
        \pure{(\delta, \semuu)} & \text{if $b = \semff$}
      \end{cases}
  \end{aligned}
\end{equation}
%
where $\pi_1$ is the projecting isomorphism from $\sem{\Delta} \times \{\semuu\}$ to $\sem{\Delta}$.

As is standard, we find the appropriate fixed point  by showing that $\PP{S}$ carries the structure of a domain ($\omega$-complete partial order with least element) and that $W_\gamma$ is an $\omega$-continuous function with respect to the induced pointwise order on the function space 
$\PP{\sem{\Delta} \times \{\semuu\}}^{\sem{\Delta}}$.
This allows the definition of $\sem{\Gamma ; \Delta \rwtypes (\cwhile e \cdo c \cend) : \dU} \, \gamma $ as the least fixed point $\mathsf{LFP}(W_\gamma)$ of $W_\gamma$:
%
\begin{equation}
\label{equation:lfp}
  \sem{\Gamma ; \Delta \rwtypes (\cwhile e \cdo c \cend) : \dU} \, \gamma  \defeq
  \mathsf{LFP}(W_\gamma) \enspace .
\end{equation}
%


The required partial order on $\PP{S}$ is that of the Plotkin powerdomain~\cite{plotkin1976powerdomain} modified to take account of our use of the error set $\PPerr$:
%
\begin{equation*}
  X \PPleq Y ~ \defiff~
  X = Y  ~\lor ~
  (\bot \in X \, \land\, (Y = \PPerr \lor (X \!\setminus\! \PPbot \subseteq Y))).
\end{equation*}
%
For $X, Y$ other than the error set $\PPerr$, the above coincides with the usual Egli-Milner order of the Plotkin powerdomain.
The positioning of~$\PPerr$ within the partial order is motivated by the following considerations.
%
The denotational semantics of $\cwhile e \cdo c \cend$ in environment~$\gamma$ is approximated by
applying $W_\gamma$ iteratively to the constantly-bottom function  $\delta \mapsto \PPbot$, yielding an $\omega$-chain
%
\begin{equation*}
  (\delta \mapsto \PPbot) \PPleq W_\gamma(\delta \mapsto \PPbot) \PPleq W_\gamma^2(\delta \mapsto \PPbot) \PPleq \dots
\end{equation*}
%
with each new approximation corresponding to one further level of unrolling of the loop. 
The presence of~$\bot$ in any $W_\gamma^n(\delta \mapsto \PPbot)\,\delta_0$ can indicate that
further unfoldings are needed to determine
$\sem{\Gamma ; \Delta \rwtypes (\cwhile e \cdo c \cend) : \dU} \, \gamma \, \delta_0$.
It is possible that some such further unfolding will result in an erroneous computation, at which point the denotational semantics will assume the value~$\PPerr$. For this reason it is necessary to have $X \PPleq \PPerr$, whenever $\bot \in X$.
In the case that $\PPbot \notin W_\gamma^n(\delta \mapsto \PPbot)\,\delta_0$, it holds that 
$\sem{\Gamma ; \Delta \rwtypes (\cwhile e \cdo c \cend) : \dU} \, \gamma \, \delta_0 = 
W_\gamma^n(\delta \mapsto \PPbot)\,\delta_0$, i.e., the semantics is fully determined at iteration $n$, and does not change under further iterations.
Thus nonempty sets $X$ with $\bot \notin X$ do not approximate $\PPerr$, i.e., $X \not\PPleq \PPerr$.

\begin{proposition} 
\label{prop:domain}
For any set $S$, it holds that $(\PP{S}, \PPleq)$ is an $\omega$-complete partial order with least element.
\end{proposition}

\begin{proof}
The least element is $\PPbot$.
%
For $\omega$-completeness, suppose that
%
\begin{equation*}
  X_0 \PPleq X_1 \PPleq X_2 \PPleq \dots
\end{equation*}
%
is an $\omega$-chain.
In the case that every $X_n$ contains $\bot$, it is easy to check that the supremum is $ \dsup{i} X_i \defeq \bigcup_i X_i$.
If instead there exists $n$ such that $\bot \notin X_n$ then necessarily $X_m = X_n$ for all $m \geq n$, so 
the supremum is $\dsup{i} X_i \defeq X_n$. (In the case that no $X_n$ is~$\PPerr$, the above coincides with the Plotkin powerdomain.)
\end{proof}

In the proof of the following proposition we use the \defemph{strict union} operation on $\PP{S}$, which models nondeterministic choice:
\[
X \uplus Y = \begin{cases}
\emptyset & \text{if $X = \emptyset$ or $Y = \emptyset$,} \\
X \cup Y & \text{otherwise.}
\end{cases}
\]

\begin{proposition} 
\label{prop:continuity}
The function $W_\gamma$ defined by \eqref{eq:def-W} is continuous with respect to the
pointwise partial order on $\PP{\sem{\Delta} \times \{\semuu\}}^{\sem{\Delta}}$.
\end{proposition}
%

\begin{proof}
%
The function $W_\gamma$ is the composition of several maps, two of which need their continuity checked.
%
The first one is monadic evaluation
%
\begin{align*}
   \PP{T}^S \times \PP{S} &\to \PP{T}  \\
  (g, X) &\mapsto \lift{g}(X)
\end{align*}
%
Monotonicity is straightforward.
%
To establish continuity in~$X$, consider an $\omega$-chain $X_0 \PPleq X_1 \PPleq X_2 \PPleq \dots$.
If every $X_n$ contains $\bot$, then 
\begin{multline*}
\textstyle
\lift{g}(\dsup{i}X_i) =
\lift{g}\left(\bigcup_{i}X_i\right) =
\{\bot\} \cup \bigcup \left\{g(x) \such x \in \left(\bigcup_i X_i\right) - \{\bot\}\right\}
\\
\textstyle
= \{\bot\} \cup \bigcup_i \{g(x) \such x \in X_i - \{\bot\}\}
= \dsup{i}  \lift{g}(X_i).
\end{multline*}
%
Otherwise the chain is eventually constant and~$g$ preserves its supremum because it is monotone.

To establish continuity in $g$, consider an $\omega$-chain $g_0 \PPleq g_1 \PPleq g_2 \PPleq \dots$ with respect to the pointwise order on $\PP{T}^S$.
%
If there are $k \in \IN$ and $x \in X$ such that $g_k(x) = \PPerr$ then $\lift{(\dsup{i} g_i)}(X) = \PPerr = \dsup{i} \lift{g_i}(X)$.
Thus it remains to prove that the sets
%
\begin{equation*}
  \textstyle
  \lift{(\dsup{i} g_i)}(X) =
      \{\bot \such \bot \in X\} \cup 
      \bigcup_{x \in X \setminus \{\bot\}} \dsup{i} g_i(x) \eqdef A
\end{equation*}
%
and
%
\begin{equation*}
  \textstyle
  \dsup{i} \lift{g_i}(X) =
  \dsup{i} \left(
      \{\bot \such \bot \in X\} \cup 
      \bigcup_{x \in X \setminus \{\bot\}} g_i(x)
   \right) \eqdef B.
\end{equation*}
%
are equal, assuming $g_i(x) \neq \PPerr$ for all $i \in \IN$ and $x \in X$.
%
Clearly, $\bot \in A \liff \bot \in B$.
%
To show that $A \setminus \{\bot\} = B \setminus \{\bot\}$, we first note that, for all $k \in \IN$, $x \in X$ and $y \neq \bot$,
%
\begin{equation*}
  y \in \dsup{i} g_i(x)
  \iff
  \some{i \in \IN} y \in g_i(x) \setminus \{\bot\},
\end{equation*}
%
thanks to the standing assumption that~$\PPerr$ does not appear in the supremum.
%
Now, if $y \in A \setminus \{\bot\}$ then $y \in \dsup{i} g_i(x)$ for some $x \in X \setminus \{\bot\}$,
hence $y \in g_j(x)$ for some $j \in \IN$, and so $y \in B$.
%
Conversely, if $y \in B \setminus \{\bot\}$ then $y \in g_j(x)$ for some $j \in \IN$ and $x \in X \setminus \{\bot\}$, hence $y \in \dsup{i} g_i(x)$, and so $y \in A$.

The other function partaking in $W_\gamma$ is
\begin{align*}
  C &: \PP{\{\semff, \semtt\}} \times \PP{S} \times \PP{S} \to \PP{S}
  \\
  C &: (X,Y,Z) \mapsto \PPlet{b}{X} 
      \begin{cases} 
        Y & \text{if $b = \semtt$} \\
        Z & \text{if $b = \semff$} 
      \end{cases}
\end{align*}
%
We only verify continuity in $Y$, which is done by case analysis on~$X$:
%
\begin{itemize}
\item If $\semtt \notin X$ then, $C$ is constant in~$Y$.
\item If $X = \pure{\semtt}$ then $C(\pure{\semtt},Y,Z) = Y$ is a projection.
\item If $X = \{\semtt, \semff\}$ then $C(Y) = Y \uplus Z$.
\item If $X = \{\semtt,\bot\}$ then $C(Y) = Y\cup \{\bot \such Y \neq \PPerr\}$.
\item If $X = \{\semtt, \semff,\bot\}$ then  $C(Y) = (Y \uplus Z) \cup \{\bot \such (Y \uplus Z) \neq \PPerr\}$.\end{itemize}
In each case, continuity in $Y$ is easy to show.
\end{proof}

It follows from \cref{prop:domain,prop:continuity} that the semantic definition of while commands in~\eqref{equation:lfp} is well-defined.


\subsection{Semantics of first-order functions}
\label{sec:semant-first-order}

We briefly address the denotational semantics of first-order functions from \cref{sec:first-order-func}.
The denotation of a function
%
\begin{equation*}
  \cfunction f {x_1 \of \tau_1, \ldots, x_n \of \tau_n} {\sigma} {e}
\end{equation*}
%
is just the denotation of its body,
%
\begin{equation*}
  \sem{f} \defeq
  \sem{x_1 \of \tau_1, \ldots, x_n \of \tau_n \rotypes e : \sigma} :
  \sem{\tau_1} \times \cdots \times \sem{\tau_n} \to \PP{\sem{\sigma}}.
\end{equation*}
%
We need to be a bit careful about the denotation of a function call $f(e_1, \ldots, e_n)$ because the arguments
$e_1, \ldots, e_n$ may diverge or yield nondeterministic values, so it matters if and when they are evaluated.
%
We opt for the call-by-value evaluation strategy that is most commonly seen in imperative languages.

Given sets~$S$ and~$T$, define the \defemph{monadic pairing}
%
$\PPtuple{{-}, {-}} : \PP{S} \times \PP{T} \to \PP{S \times T}$
%
by
%
\begin{equation*}
  \PPtuple{X, Y} \defeq (\PPlet{x}{X} \PPlet{y}{Y} \{(x, y)\}).
\end{equation*}
%
Note that the binding order does not matter, i.e., $(\PPlet{x}{X} \PPlet{y}{Y} \cdots) = (\PPlet{y}{Y} \PPlet{x}{X} \cdots)\})$ because $\Pstar$ is a commutative monad.
%
Monadic pairing is easily extended from pairs to $n$-tuples for arbitrary~$n$.

The denotation of a function call is application adorned with the monad structure:
%
\begin{align*}
  \sem{\Gamma \rotypes f(e_1, \ldots, e_n) : \sigma} \gamma \defeq &
  \lift{\sem{f}} \PPtuple{\sem{e_1} \gamma, \ldots, \sem{e_n} \gamma}.
\end{align*}


\section{Nondeterminism and parallelism}
\label{sec:boolean-ops}

The guarded case construct of Clerical requires parallel evaluation of the guards
leading to potential nondeterminism.  As a result, several basic operations involving nondeterminism and parallel evaluation are definable in Clerical.

A simple binary nondeterministic choice between two pure expressions
$\Gamma \rotypes e_1 : \tau$ and $\Gamma \rotypes e_2 : \tau$ is implemented
by
\[\Gamma \rotypes (\ccase \ctrue \To e_1 \such \ctrue \To e_2 \cend) : \tau \]
This has the derived semantics
\[
\sem{\Gamma \rotypes (\ccase \ctrue \To e_1 \such \ctrue \To e_2 \cend) : \tau }\, \gamma
~ = ~ \sem{\Gamma \rotypes e_1 : \tau } \, \gamma\,
\uplus \,
\sem{\Gamma \rotypes e_2  : \tau }\,\gamma
\]
using the strict union operation defined above Proposition~\ref{prop:continuity}.

It is also possible to implement McCarthy's \defemph{ambiguous choice} between 
$\Gamma \rotypes e_1 : \tau$ and $\Gamma \rotypes e_2 : \tau$, by:
\[\Gamma \rotypes (\ccase (\cvar{x}{e_1} \ctrue) \To e_1 \such 
(\cvar{x}{e_2} \ctrue) \To e_2 \cend) : \tau \]
Writing $\Gamma \rotypes e_1 \,\mathtt{amb}\, e_2 : \tau$ for the above, 
we have
%
\begin{multline*}
\sem{\Gamma \rotypes  e_1 \,\mathtt{amb}\, e_2 : \tau }\, \gamma = \\
  (\sem{\Gamma \rotypes e_1 : \tau } \, \gamma\, \uplus \, \sem{\Gamma \rotypes e_2  : \tau }\,\gamma)
  \setminus
  \{ \bot \mid \bot \notin (\sem{\Gamma \rotypes e_1  : \tau }\,\gamma 
\,\cap \,\sem{\Gamma \rotypes e_2  : \tau}\,\gamma)\}\;.
\end{multline*}

A well-known issue with ambiguous choice is that it is not monotonic with respect to any reasonable powerdomain partial ordering \cite{LEVY2007221}, meaning that it does not have a domain-theoretic semantics. Indeed, such a failure of monotonicity holds with respect to 
the ordering $\PPleq$ we have defined on our powerdomain $\PP{-}$. It follows that the denotational semantics of Clerical expressions does not, in general, act monotonically in the semantics of subexpressions. We present a  simple example of this phenomen.


Consider the expression below, which is parametric in the subexpression
$\rotypes e  : \dB$:
%
\begin{lstlisting}
case
| e => while true do skip end
| true => skip
end
\end{lstlisting}
%
In the case that $\sem{e} =  \pure{\bot}$, the denotation of the whole expression is $\{\semuu\} $. If instead $\sem{e} =  \pure{\semtt}$, then the denotation of the expression is $\{\semuu, \bot\}$. This breaks monotonicity because
$\pure{\bot} \PPleq  \pure{\semtt}$ in $\PP{\{\semff, \semtt\}}$, but $\{\semuu\} \not \PPleq \{\semuu, \bot\}$  in $\PP{\{\semuu\}}$.
Similarly, considering the case in which $e$ is an erroneous expression (i.e., $\sem{e} = \PPerr$),
we have $\pure{\bot} \PPleq  \PPerr$, but $\{\semuu\} \not \PPleq \PPerr$.


Given the non-monotonicity properties illustrated above, it is perhaps surprising that it is possible to give Clerical a denotational semantics involving a domain-theoretic fixed-point argument to establish the semantics of while loops. The reason for this is that the operator $W_\gamma$, used in defining the semantics of while loops, is defined purely as a combination of conditional statements and sequencing, and does not involve the problematic non-monotonic aspects of Clerical. Indeed, as Proposition~\ref{prop:continuity} establishes,  the particular operator $W_\gamma$ is always continuous, hence \emph{a fortiori}  monotone. 

Clerical, as we have defined it, does not include any primitive operator for manipulating booleans. This does not limit expressivity since logical operations are definable using the conditional construct.
For example, negation of a boolean expression $b$ is defined by
\[
\neg b \defeq (\cif b \cthen \cfalse \celse  \ctrue  \cend)
\]
The most concise way of defining the disjunction of two boolean typed expressions $b_1$ and $b_2$ is by $\cif b_1 \cthen \ctrue \celse  b_2 \cend$. This is asymmetric: if 
$\sem{b_1} = \pure{\semtt}$ and $\sem{b_2} = \pure {\bot}$ then the disjunction has denotation $\pure{\semtt}$, whereas if
$\sem{b_1} = \pure{\bot}$ and $\sem{b_2} = \pure {\semtt}$ then 
the resulting denotation is $\pure{\bot}$. It is similarly possible to define a symmetric strict disjunction by 
\[\cif b_1 \cthen (\cif b_2 \cthen \ctrue \celse \ctrue \cend) \celse  b_2 \cend\]
More interestingly, the guarded case construct can be used to define Plotkin's \defemph{parallel or} \cite{plotkin1977lcf}, another symmetric version of  disjunction which requires parallel evaluation, by
\begin{equation*}
b_1 \cdisj b_2 \defeq
(\ccase
      b_1 \To \ctrue
\such b_2 \To \ctrue
\such \neg b_1 \To \neg b_2
\cend)
\end{equation*}
From a logical perspective, when applied to deterministic expressions, $b_1 \cdisj b_2$ computes the disjunction of $b_1$ and $b_2$ from Kleene (and Priest) 3-valued logic.

Since Clerical is a nondeterministic language, the defined logical operators all have an induced effect on nondeterministic and erroneous expressions. For example, the derived full semantics for parallel or is:
\begin{align*}
\sem{\Gamma \rotypes b_1 \cdisj b_2 : \dB }(\gamma) &= \bigcup_{\substack{v_1 \in \sem{\Gamma \rotypes b_1:\dB}(\gamma)\\v_2 \in \sem{\Gamma \rotypes b_2 : \dB}(\gamma)}}\begin{cases}
\{\semtt\}&\text{if }v_1 = \semtt \lor v_2 = \semtt,\\
\{\semff\}&\text{if }v_1 = \semff \land v_2 = \semff,\\
\{\bot\}&\text{otherwise}.
\end{cases}
\end{align*}



%%% Local Variables:
%%% mode: latex
%%% TeX-master: "clerical"
%%% End:

\chapter{Example Sheets}

%% TJWP removed plain tex commands
%\magnification\magstephalf
%\font\bb=msbm10 scaled 1200
%\font\small=cmr8
%\font\little=cmmi7
%%%%%%%%%%%%%%%%%%%%%%%

%% TJWP include page numbers
%\nopagenumbers
\def \pr{\partial}
\def\pmb#1{\setbox0=\hbox{#1}%
 \kern-.025em\copy0\kern-\wd0
 \kern.05em\copy0\kern-\wd0
 \kern-.025em\raise.0433em\box0 }
\def \bE{{\pmb {${\cal E}$}}}
\noindent

%\centerline{{\bf Example Sheet 1}}
%\vskip 10pt
\section{Example Sheet 1}

\noindent
{\bf 1.} Explain why 
\vskip 0.3cm
(i) GR effects are important for neutron stars but not for white
dwarfs
\vskip 0.3cm
(ii) inverse beta-decay becomes energetically favourable for
densities higher than those in white dwarfs.
\vskip 10 pt
\noindent
{\bf 2.} Use Newtonian theory to derive the Newtonian pressure
support equation
$$
P'(r) \equiv {dP\over dr} = -{Gm\rho\over r^2}\ ,
$$
where
$$
m=4\pi \int_0^r \tilde r^2 \rho(\tilde r) d\tilde r\ ,
$$
for a spherically-symmetric and static star with pressure $P(r)$ and 
density $\rho(r)$. Show that
%% TJWP no eqalign in latex  \eqalign{
\bean
\int_0^r P(\tilde r)\tilde r^3 d\tilde r & = & 
{P(r)r^4\over 4} -{1\over 4}\int_0^r P'(\tilde r) \tilde r^4 d\tilde 
r \\
 & = & {Gm^2(r)\over 32\pi} + {P(r)r^4\over 4}\ .
\eean
Assuming that $P'\le 0$, with $P=0$ at the star's surface, show that
$$
{d\over dr}\left[\left(\int_0^r P(\tilde r)\tilde r^3 d\tilde
r\right)^{3/4}\right] \le {3\sqrt{2}\over 4} P^{3/4} r^2\ .
$$
Assuming the bound
$$
P\ {\buildrel <\over\sim}\ (\hbar c) n_e^{4/3}\ ,
$$
where $n_e(r)$ is the electron number density, show that the total
mass, $M$, of the star satisfies
$$
M\ {\buildrel <\over\sim}\ \left({hc\over
G}\right)^{3/2}\left({\mu_e\over m_N}\right)^2 
$$
where $m_N$ is the nucleon mass and $\mu_e$ is the number of
electrons per nucleon. Why is it reasonable to bound the pressure as
you have done? Compare your bound with Chandresekhar's limit.

\vskip 10 pt
\noindent
{\bf 3.} A particle orbits a Schwarzschild black hole with
non-zero angular momentum per unit mass $h$. Given that $\sigma=0$ for
a massless particle and $\sigma=1$ for a massive particle, show that
the orbit satisfies  
$$ 
{d^2 u\over d\phi^2} + u = {M\sigma \over
h^2} + 3Mu^2 
$$
where $u=1/r$ and $\phi$ is the azimuthal angle. Verify that this
equation is solved by
$$
u= {1\over 6M} + {2\omega^2\over 3M} -{2\omega^2\over M\cosh ^2
(\omega\phi)}\ ,
$$
where $\omega$ is given by
$$
4\omega^2 = \pm\sqrt{\left(1 - {12M^2\sigma\over h^2}\right)}\ .
$$
where $\sigma=1$ for a massive particle and $\sigma=0$ for a massless
particle. Interpret these orbits in terms of the effective potential.
Comment on the cases $\omega^2=1/4$, $\omega^2=1/8$ and $\omega^2=0$.

\vskip 10 pt
\noindent
{\bf 4.} A photon is emitted outward from a point P outside
a Schwarzschild black hole with radial coordinate r in the range
$2M<r<3M$. Show that if the photon is to reach infinity the angle its
initial direction makes with the radial direction (as determined by 
a stationary observer at P) cannot exceed
$$
{\rm arcsin} \sqrt{{27M^2\over r^2}\left( 1-{2M\over r}\right)}\ .
$$

\vskip 10 pt
\noindent
{\bf 5.} Show that in region II of the Kruskal manifold one may
regard $r$ as a time coordinate and introduce a new spatial
coordinate $x$ such that
$$
ds^2 = -{dr^2\over \left({2M\over r} -1\right)} +\left({2M\over
r}-1\right)dx^2 + r^2d\Omega^2\ .
$$
Hence show that {\it every} timelike
curve in region II intersects the singularity at $r=0$ within a proper
time no greater than $\pi M$. For what curves is this bound attained?
Compare your result with the time taken for the collapse of a ball of
pressure free matter of the same gravitational mass $M$. Calculate
the binding energy of such a ball of dust as a fraction of its
(conserved) rest mass.
\vskip 10 pt
\noindent
{\bf 6.} Using the map
$$
(t,x,y,z) \mapsto X= \pmatrix{t+z & x+iy\cr x-iy & t-z}\ ,
$$
show that Minkowski spacetime may be identified with the space of
Hermitian $2\times 2$ matrices $X$ with metric
$$
ds^2 = -\det (dX)\ .
$$
Using the Cayley map $X\mapsto U={1+iX\over1-iX}$, show further that
Minkowski spacetime may be identified with the space of unitary
$2\times 2$ matrices $U$ for which $\det (1+U)\ne0$. Now show that any
$2\times 2$ unitary matrix $U$ may be expressed uniquely in terms of
a real number $\tau$ and two complex numbers $\alpha$, $\beta$, as 
$$
U=e^{i\tau}\pmatrix{\alpha & \beta\cr -\bar\beta & \bar\alpha}
$$
where the parameters $(\tau,\alpha,\beta)$ satisfy $|\alpha|^2
+|\beta|^2 =1$, and are subject to the identification 
$$
(\tau,\alpha,\beta) \sim (\tau +\pi, -\alpha, -\beta)\ .
$$
Using the relation
$$
(1+U)dX = -2i dU(1+U)^{-1}\ ,
$$
deduce that 
$$
ds^2 = {1\over (\cos \tau + {\cal R}e\,\alpha)^2}\left(-d\tau^2
+|d\alpha|^2 + |d\beta|^2\right) 
$$
is the metric on Minkowski spacetime and hence conclude that the
conformal compactification of Minkowski spacetime may be identified
with the space of unitary $2\times 2$ matrices, i.e the group
$U(2)$. Explain how $U(2)$ may be identified with a portion of the
Einstein static universe $S^3\times {\bb R}$.

\vfill\eject

%\centerline {{\bf Example Sheet 2}}
%\vskip 10 pt
\section{Example Sheet 2}

\noindent
{\bf 1.} Let $\zeta$ be a Killing vector field. Prove that
$$
D_\sigma D_\mu \zeta_\nu = R_{\nu\mu\sigma}{}^\lambda \zeta_\lambda\ ,
$$
where $R_{\nu\mu\sigma\lambda}$ is the Riemann tensor, defined by
$[D_\mu,D_\nu]\, v_\rho = R_{\mu\nu\rho}{}^\sigma v_\sigma$ for
arbitrary vector field $v$.

\vskip 10 pt
\noindent
{\bf 2.}  A conformal Killing vector is one for which  
$$
({\cal L}_\xi g)_{\mu\nu} = \Omega^2 g_{\mu\nu}\ .
$$
for some non-zero function $\Omega$.
Given that $\xi$ is a Killing vector of $ds^2$, show that it is a
conformal Killing vector of the conformally-equivalent
metric $\Lambda^2 ds^2$ for arbitrary (non-vanishing) conformal factor
$\Lambda$.
 
Show that the action for a {\it massless} particle,
$$ S[x,e]={1\over2}\int d\lambda\, e^{-1}\dot x^\mu\dot x^\nu
g_{\mu\nu}(x)\ , $$ is invariant, to first order in the constant
$\alpha$, under the transformation 
$$
x^\mu \rightarrow x^\mu + \alpha \xi^\mu(x) \qquad \qquad
e \rightarrow e + {1\over4}\alpha\, e g^{\mu\nu}({\cal L}_\xi
g)_{\mu\nu}
$$
if $\xi =\xi^\mu \partial_\mu$ is a conformal Killing vector. Show
that $\xi$ is the operator corresponding to the conserved charge
implied by Noether's theorem.



\vskip 10 pt
\noindent
{\bf 3.} Show that the extreme RN metric in isotropic coordinates is
$$
ds^2 = -\left(1+{M\over\rho}\right)^{-2}dt^2 + 
\left(1+{M\over\rho}\right)^{2}\left(d\rho^2 + \rho^2 d\Omega^2\right)
\qquad (\dagger)
$$
Verify that $\rho=0$ is at infinite proper distance from any finite
$\rho$ along any curve of constant $t$. Verify also that
$|t|\rightarrow \infty$ as $\rho\rightarrow 0$ along any timelike or
null curve but that a timelike or null ingoing radial geodesic reaches
$\rho=0$ for {\it finite} affine parameter. By introducing a null
coordinate to replace $\rho$ show that $\rho=0$ is merely a coordinate
singularity and hence that the metric ($\dagger$) is geodesically
incomplete. What happens to the particles that reach $\rho=0$?
Illustrate your answers using a Penrose diagram.

\vskip 10 pt
 \noindent
{\bf 4.} The action for a particle of mass $m$ and charge $q$ is
$$
S[x,e] =\int d\lambda\,\left[{1\over2} e^{-1}\dot x^\mu\dot x^\nu
g_{\mu\nu}(x) -{1\over2}m^2 e -q\, \dot x^\mu A_\mu(x)\right]\qquad
\qquad (*) 
$$
where $A_\mu$ is the electromagnetic 4-potential. Show
that if
$$
({\cal L}_\xi A)_\mu \equiv \xi^\nu\partial_\nu A_\mu +
(\partial_\mu \xi^\nu) A_\nu =0
$$
for Killing vector $\xi$, then $S$ is invariant, to first-order in
$\xi$, under the transformation $x^\mu\rightarrow x^\mu
+\alpha\xi^\mu(x)$. Verify that the corresponding Noether charge
$$
-\xi^\mu \left(m u_\mu -q A_\mu\right)\ ,
$$ 
where $u^\mu$ is the particle's 4-velocity, is a constant of the 
motion. Verify for the Reissner-Nordstrom solution of the vacuum
Einstein-Maxwell equations, with mass $M$ and charge $Q$, that ${\cal
L}_k A =0$ for $k={\partial\over \partial t}$ and hence deduce, for
$m\ne 0$, that  
$$
\left(1-{2M\over r} + {Q^2\over r^2}\right) {dt\over d\tau} =
\varepsilon -{q\over m}{Q\over r}\ ,
$$
where $\tau$ is the particle's proper time and $\varepsilon$ is the
energy per unit mass. Show that the trajectories $r(t)$ of
massive particles with zero angular momentum satisfy
$$
({dr\over d\tau})^2 = (\varepsilon^2 -1) + \left(1-\varepsilon
{qQ\over mM}\right){2M\over r} +\left(\left({q\over m}\right)^2
-1\right){Q^2\over r^2}\ .
$$
Give a physical interpretation of this result for the special case
for which $q^2=m^2$, $qQ=mM$, and $\varepsilon=1$.


\vskip 10 pt
\noindent
{\bf 5.} Show that the action
$$
S[p,x,e]=\int d\lambda \big\{ p_\mu \dot x^\mu
-{1\over2}e\, \big[g^{\mu\nu}(x)p_\mu p_\nu + m^2\big]\big\}
$$
for a point particle of mass $m$ is equivalent, for $q=0$, to the action
of Q.4. Show that $S$ is invariant to
first order in $\alpha$ under the transformation
$$
\delta x^\mu =\alpha K^{\mu\nu}p_\nu \qquad \delta p_\mu
=-{1\over2}\alpha\, p_\rho p_\sigma \partial_\mu K^{\rho\sigma}
$$
for any symmetric tensor $K_{\mu\nu}$ obeying the {\it Killing
tensor} condition
$$
D_{(\rho} K_{\mu\nu)}=0\ .
$$
Show that the corresponding Noether charge is proportional to
$K^{\mu\nu}p_\mu p_\nu$ and verify that it is a constant of the
motion. A trivial example is $K_{\mu\nu}=g_{\mu\nu}$; what is
the corresponding constant of the motion? Show that
$\xi_\mu\xi_\nu$ is a Killing tensor if $\xi$ is a Killing vector. [A
Killing tensor that cannot be constructed from the metric and Killing
vectors is said to be irreducible. In a general axisymmetric
metric there are no such tensors, and so only three constants of the
motion, but for geodesics of the Kerr-Newman metric there is a 
`fourth constant' of the motion corresponding to an
irreducible Killing tensor.]
%$$
%\eqalign{
%K_{\mu\nu}dx^\mu dx^\nu = - &{\Sigma a^2\cos^2\theta\over \Delta}dr^2
%+ {\Delta a^2 \cos^2\theta\over \Sigma}(dt -a\sin^2\theta\,
%d\phi)^2\cr & +{r^2\sin^2\theta\over\Sigma}[adt - (r^2 + a^2)d\phi]^2
%+ r^2\Sigma
% d\theta^2\ .}
%$$

\vskip 10 pt 
\noindent 
{\bf 6.} By replacing the time coordinate $t$
by one of the radial null coordinates
$$
u= t+ {M\over \lambda} \qquad v= t- {M\over \lambda}
$$
show that the singularity at $\lambda=0$ of the Robinson-Bertotti (RB)
metric
$$
ds^2 = -\lambda^2 dt^2 + M^2 \left({d\lambda\over \lambda}\right)^2 +
M^2 d\Omega^2 
$$
is merely a coordinate singularity.
Show also that $\lambda=0$ is a degenerate Killing Horizon with respect
to $\partial\over \partial t$. By
introducing the new coordinates $(U,V)$, defined by
$$
u= \tan \left({U\over 2}\right)\qquad v=-\cot \left({V\over2}\right)
$$
obtain the maximal analytic extension of the RB metric and deduce its
Penrose diagram.

\vfill\eject

%\centerline {{\bf Example Sheet 3}}
%\vskip 10pt
\section{Example Sheet 3}

\noindent 
{\bf 1.} Let $\varepsilon$ and $h$ be the energy and and angular
momentum per unit mass of a zero charge particle in free fall
within the equatorial plane, i.e on a timelike ($\sigma=1$) or null
($\sigma=0$) geodesic with $\theta=\pi/2$, of a Kerr-Newman black hole.
Show that the particle's Boyer-Lindquist radial coordinate $r$
satisfies   
$$ 
\left({dr\over d\lambda}\right)^2 =\varepsilon^2 - V_{eff}(r)\ ,
$$ 
where $\lambda$ is an affine parameter, and the effective potential
$V_{eff}$ is given by $$ 
V_{eff} =
\left(1-{2M\over r}+{e^2\over r^2}\right)\left(\sigma + {h^2\over
r^2}\right) + {2a\varepsilon h\over r^3}\left( 2M -{e^2\over r}\right)
+{a^2\over r^2}\left[\sigma-\varepsilon^2\left(1+{2M\over r}-{e^2\over
r^2}\right)\right] \ .
$$

\vskip 10 pt 
\noindent 
{\bf 2.} Show that the surface gravity of the event horizon of a Kerr
black hole of mass $M$ and angular momentum $J$ is given by
$$
\kappa = {\sqrt{M^4 -J^2}\over 2M( M^2 + \sqrt{M^4 -J^2})}\ .
$$

\vskip 10 pt 
\noindent 
{\bf 3.} A particle at fixed $r$ and $\theta$ in a stationary
spacetime, with metric $ds^2= g_{\mu\nu}(r,\theta)dx^\mu dx^\nu$, has
angular velocity $\Omega= {d\phi\over dt}$ with respect to infinity.
Show that $\Omega(r,\theta)$ must satisfy
$$
g_{tt} + 2g_{t\phi}\Omega + g_{\phi\phi}\Omega^2 \le 0
$$
and hence deduce that
$$
{\cal D}\equiv g_{t\phi}^2 - g_{tt}g_{\phi\phi} \ge 0
$$
Show that ${\cal D}=\Delta (r) \sin^2\theta$ for the Kerr-Newman 
metric in Boyer-Lindquist coordinates, where $\Delta= r^2-2Mr+a^2+e^2$.
What happens if $(r,\theta)$ are such that ${\cal D}<0$? For what
values of $(r,\theta)$ can $\Omega$ vanish? Given that $r_{\pm}$ are the 
roots of $\Delta$, show that when ${\cal D}=0$  
$$ 
\Omega = {a\over r_{\pm}^2 +a^2} \ .
$$
\vskip 10pt
\noindent
{\bf 4.} Show that the area of the event horizon of a Kerr-Newman
black hole is
$$
A= 8\pi\big[ M^2 - {e^2\over2} + \sqrt{ M^4 -e^2 M^2 -J^2 }\,\big]\ .
$$

\vskip 10 pt
\noindent
{\bf 5.} A perfect fluid has stress tensor
$$
T_{\mu\nu} = (\rho + P)u_\mu u_\nu + P g_{\mu\nu}\ ,
$$
where $\rho$ is the density and $P(\rho)$ the pressure. State the
dominant energy condition for $T_{\mu\nu}$ and show that
for a perfect fluid in Minkowski spacetime this condition
is equivalent to  
$$ 
\rho\ge |P|\ .
$$
Show that the same condition arises from the requirement of
causality, i.e. that the speed of sound, $\sqrt{|dP/d\rho|}$, not
exceed that of light, together with the fact that the pressure 
vanishes in the vacuum.
\vskip 10pt
\noindent
{\bf 6.} The vacuum Einstein-Maxwell equations are
$$
G_{\mu\nu}= 8\pi T_{\mu\nu}(F) \qquad D_\mu F^{\mu\nu}=0
$$
where $F_{\mu\nu}= \partial_{\mu}A_{\nu}- \partial_\nu A_\mu$, and
$$
T_{\mu\nu}(F)= {1\over 4\pi}\big(F_{\mu}{}^\lambda F_{\nu\lambda}
-{1\over4}g_{\mu\nu}F^{\alpha\beta}F_{\alpha\beta}\big)\ .
$$
Asymptotically-flat solutions are stationary and axisymmetric if
the metric admits Killing vectors $k$ and $m$ that can be taken to be
$k={\partial\over\partial t}$ and $m={\partial\over\partial \phi}$ near
infinity, and if (for some choice of electromagnetic gauge)
$$
{\cal L}_k A={\cal L}_m A=0\ ,
$$
where the Lie derivative of $A$ with respect to a vector $\xi$, 
${\cal L}_\xi A$, is as defined in Q.4 of Example Sheet 2. The event
horizon of such a solution is necessarily a Killing horizon of $\xi =
k+\Omega_H m$, for some constant $\Omega_H$. What is the physical
interpretation of $\Omega_H$? What is its value for the Kerr-Newman
solution? The co-rotating electric potential is defined by 
$$ 
\Phi = \xi^\mu A_\mu \ .
$$
Use the fact that $R_{\mu\nu}\xi^\mu\xi^\nu=0$ on a Killing horizon
to show that $\Phi$ is constant on the horizon. In particular, show
that for a choice of the electromagnetic gauge for which $\Phi=0$ at
infinity, 
$$
\Phi_H= {Qr_+\over r_+^2 +a^2}
$$
for a charged rotating black hole, where $r_+= M+\sqrt{M^2-Q^2-a^2}$.

\vskip 10pt
\noindent
{\bf 7.} Let $({\cal M},g,A)$ be an asymptotically flat, stationary,
axisymmetric, solution of the Einstein-Maxwell equations of Q.6 and
let $\Sigma$ be a spacelike hypersurface with one boundary at spatial
infinity and an internal boundary, $H$, at the event horizon of a black
hole of charge $Q$. Show that    
$$ 
-2\int_\Sigma dS_\mu T^\mu{}_\nu(F)\xi^\nu = \Phi_H Q  
$$
where $\Phi_H$ is the co-rotating electric potential on the horizon. 
Use this result to deduce that the mass $M$ of a charged rotating black
hole is given by 
$$
M= {\kappa A\over 4\pi} + 2\Omega_H J + \Phi_H Q\ .
$$
where $J$ is the total angular momentum.
Use this formula for $M$ to deduce the first law of black hole
mechanics for charged rotating black holes: 
$$
dM= {\kappa \over 8\pi}dA + \Omega_H dJ + \Phi_H dQ\ .
$$
[Hint: ${\cal L}_\xi (F^{\mu\nu}A_\nu)=0$ ]

\vfill\eject


%\centerline{{\bf Example Sheet 4}}
%\vskip 10pt
\section{Example Sheet 4}

\noindent
{\bf 1.} Use the Komar integral,
$$
J= {1\over 16\pi G}\oint_\infty dS_{\mu\nu}D^\mu m^\nu\ ,
$$
for the total angular momentum of an asymptotically-flat axisymmetric
spacetime (with Killing vector $m$) to verify that $J=Ma$ for the
Kerr-Newman solution with parameter $a$.

\vskip 10pt
\noindent
{\bf 2.} Let $l$ and $n$ be two linearly independent vectors and
$\hat B$ a second rank tensor such that
$$
\hat B_\mu{}^\nu l_\nu =\hat B_\mu{}^\nu n_\nu =0\ .
$$
Given that $\eta^{(i)}$ $(i=1,2)$ are two further linearly
independent vectors, show that
$$
\varepsilon^{\mu\nu\rho\sigma}l_\mu n_\nu \hat B_\rho{}^\lambda
\big(\eta^{(1)}_\lambda\eta^{(2)}_\sigma - 
\eta^{(1)}_\sigma\eta^{(2)}_\lambda\big) =  \theta\, 
\varepsilon^{\mu\nu\rho\sigma} l_\mu n_\nu
\eta_\rho^{(1)}\eta_\sigma^{(2)}\ .
$$
where $\theta= \hat B_\alpha{}^\alpha$.

\vskip 10pt
\noindent
{\bf 3.} Let ${\cal N}$ be a Killing horizon of a Killing vector field
$\xi$, with surface gravity $\kappa$. Explain why, for any third-rank
totally-antisymmetric tensor $A$, the scalar 
$\Psi = A^{\mu\nu\rho}(\xi_\mu D_\nu\xi_\rho)$ vanishes on ${\cal N}$.
Use this to show that
$$
(\xi_{[\rho}D_{\sigma]} \xi_\nu)(D^\nu\xi^\mu) =\kappa
\xi_{[\rho}D_{\sigma]}\xi^\mu \qquad ({\rm on}\ {\cal N})\ ,\qquad (*)
$$
where the square brackets indicate antisymmetrization on the enclosed
indices.

>From the fact that $\Psi$ vanishes on ${\cal N}$ it follows that its
derivative on ${\cal N}$ is normal to ${\cal N}$, and hence that
$\xi_{[\mu}\partial_{\nu]}\Psi=0$ on ${\cal N}$. Use this fact and the
Killing vector lemma of Q.II.1 to deduce that, on ${\cal N}$,
$$
(\xi_\nu R_{\sigma\rho[\beta}{}^\lambda\xi_{\alpha]}
+\xi_\rho R_{\nu\sigma[\beta}{}^\lambda\xi_{\alpha]}
+\xi_\sigma R_{\rho\nu[\beta}{}^\lambda\xi_{\alpha]})\xi_\lambda\ .
$$
Contract on $\rho$ and $\alpha$ and use the fact that $\xi^2=0$ on
${\cal N}$ to show that
$$
\xi^\nu\xi_{[\rho}R_{\sigma]\nu\mu}{}^\lambda\xi_\lambda =
-\xi_\mu\xi_{[\rho}R_{\sigma]}{}^\lambda\xi_\lambda
\qquad ({\rm on}\ {\cal N})\ ,\qquad (\dagger)
$$
where $R_{\mu\nu}$ is the Ricci tensor.

For any vector $v$ the scalar $\Phi=(\xi\cdot D\xi -\kappa\xi)\cdot v$
vanishes on ${\cal N}$. It follows that
$\xi_{[\mu}\partial_{\nu]}\Phi|_{\cal N} =0$. Show that this fact, the
result (*) derived above and the Killing vector lemma imply
that, on ${\cal N}$, 
%% $$ \eqalign{
\bean
\xi^\mu\xi_{[\rho}\partial_{\sigma]}\kappa & = & \xi^\nu
R_{\mu\nu[\sigma}{}^\lambda\xi_{\rho]}\xi_\lambda \\
& = &\xi^\nu\xi_{[\rho}R_{\sigma]\nu\mu}{}^\lambda\xi_\lambda\ ,
\eean
%% } $$
where the second line is a consequence of the cyclic identity
satisfied by the Riemann tensor. Now use $(\dagger)$ to show that, on
${\cal N}$,
%$$ \eqalign{
\bea
\xi^\mu\xi_{[\rho}\partial_{\sigma]}\kappa & = &
\xi_{[\sigma}R_{\rho]}{}^\lambda\xi_\lambda \\
&= & 8\pi G\, \xi_{[\sigma}T_{\rho]}{}^\lambda\xi_\lambda\ ,
\eea
%} $$
where the second line follows on using the Einstein equations. Hence
deduce the zeroth law of black hole mechanics: that, provided the
matter stress tensor satisfies the dominant energy condition, the
surface gravity of any Killing vector field $\xi$ is constant on
each connected component of its Killing horizon (in particular, on the
event horizon of a stationary spacetime). 
% TJWP removed \vfill\eject

\vskip 10pt
\noindent
{\bf 4.} A scalar field $\Phi$ in the Kruskal spacetime satisfies
the Klein-Gordon equation
$$
D^2\Phi -m^2\Phi =0\ .
$$
Given that, in static Schwarzshild coordinates, $\Phi$ takes the form
$$
\Phi = R_\ell(r) e^{-i\omega t} Y_{\ell}(\theta,\phi)
$$
where $Y_{\ell\, m}$ is a spherical harmonic, find the radial equation
satisfied by $R_\ell(r)$. Show that near the horizon at $r=2M$,
$\Phi\sim e^{\pm i\omega r^*}$, where $r^*$ is the Regge-Wheeler radial
coordinate. Verify that ingoing waves are analytic, in Kruskal
coordinates, on the future horizon, ${\cal H}^+$, but not, in general,
on the past horizon, ${\cal H}^-$, and conversely for outgoing waves.

Given that both $m$ and $\omega$ vanish, show that
$$
R_\ell = A_\ell P_\ell(z) + B_\ell Q_\ell(z)
$$
for constants $A_\ell,\, B_\ell$, where $z=(r-M)/M$, $P_\ell(z)$ is a
Legendre Polynomial and $Q_\ell(z)$ a linearly-independent solution.
Hence show that there are no {\it non-constant} solutions that are both
regular on the horizon, ${\cal H}= {\cal H}^+ \cup {\cal H}^-$, and
bounded at infinity.

\vskip 10 pt
\noindent
{\bf 5.} Use the fact that a Schwarzschild black hole radiates at the
Hawking temperature
$$
T_H ={1\over 8\pi M}
$$
(in units for which $\hbar$, $G$, $c$, and Bolzmann's constant all
equal $1$) to show that the thermal equilibrium of a black hole with an
infinite reservoir of radiation at temperature $T_H$ is unstable.

A finite reservoir of radiation of volume $V$ at temperature $T$ has
an energy, $E_{res}$ and entropy, $S_{res}$ given by
$$
E_{res} = \sigma VT^4 \qquad S_{res} ={4\over3}\sigma VT^3
$$
where $\sigma$ is a constant. A Schwarzschild black hole of mass $M$ is
placed in the reservoir. Assuming that the black hole has entropy
$$
S_{BH} =4\pi M^2\ ,
$$
show that the total entropy $S= S_{BH}+S_{res}$ is extremized 
for fixed total energy $E= M+E_{res}$, when $T=T_H$, Show that the
extremum is a maximum if and only if $V<V_c$, where the critical value
of $V$ is $$
V_c = {2^{20}\pi^4E^5\over 5^5\sigma }
$$
What happens as $V$ passes from $V<V_c$ to $V>V_c$, or
vice-versa?


\vskip 10pt
\noindent
{\bf 6.} The specific heat of a charged black hole of mass $M$, at
fixed charge $Q$, is 
$$
C\equiv T_H {\partial S_{BH}\over \partial
T_H}\bigg|_Q \ ,
$$
where $T_H$ is its Hawking temperature and $S_{BH}$ its
entropy. Assuming that the entropy of a black hole is given by $S_{BH}=
{1\over4}A$, where $A$ is the area of the event horizon, show that the
specific heat of a Reissner-Nordstrom black hole is
$$
C= {2S_{BH}\sqrt{M^2-Q^2}\over (M-2\sqrt{M^2-Q^2})}\ .
$$
Hence show that $C^{-1}$ changes sign when $M$ passes through
${2|Q|\over\sqrt{3}}$. 

Repeat Q.5 for a Reissner-Nordstrom black hole.
Specifically, show that the critical reservoir volume, $V_c$, is
infinite for $|Q|\le M \le {2|Q|\over\sqrt{3}}$. Why is this result
to be expected from your previous result for $C$? 



%\end










\section{Implementation}
\label{sec:implementation}

To experiment with \eff we implemented a prototype interpreter whose main evaluation loop
is essentially the same as the denotational semantics described in
Section~\ref{sec:semantics}. Apart from inessential improvements, such as recursive type
definitions, inclusion of $\kord{for}$ and $\kord{while}$ loops, and pattern matching, the
implemented language differs from the one presented here in two ways that make it usable:
we implemented Hindley-Milner style type inference with parametric
polymorphism~\cite{milner78atheory}, and in the concrete syntax we hid the distinction
between expressions and computations. We briefly discuss each.

There are no surprises in the passage from monomorphic type checking to parametric
polymorphism and type inference. The infamous value
restriction~\cite{wright95simpleimperative} is straightforward because the distinction
between expressions and computations corresponds exactly to the distinction between
nonexpansive and expansive terms. In fact, it may be worth investing in an effect system
that would relax the value restriction on those computations
that can safely be presumed pure. Because $\new E$ is a computation, effect instances are
\emph{not} polymorphic, which is in agreement with ML-style references being
non-polymorphic.
%, see Section~\ref{sec:state}.
% There is nothing in this section about value restriction

A syntactic division between pure expressions and possibly effectful computations is
annoying because even something as simple as $f \, x \, y$ has to be written as
$\letin{g = f\, x} g \, y$, and having to insert $\kord{val}$ in the right places is no
fun either. Therefore, the concrete syntax allows the programmer to arbitrarily mix expressions
and computations, and a desugaring phase appropriately separates the two.

The desugaring process is fairly simple. It inserts a $\kord{val}$ when an expression
stands where a computation is expected. And if a computation stands where an expression is expected, the computation is hoisted to an enclosing $\kord{let}$
statement. Because several computations may be hoisted from a single expression, the
question arises how to order the corresponding $\kord{let}$ statements. For example, $(f\, x, g\, y)$ can be desugared as
%
\begin{equation*}
  \begin{aligned}
    &\letin{a = f \, x} \\ &\quad \letin{b = g \, y} (a, b)
  \end{aligned}
  \qquad\text{or}\qquad
  \begin{aligned}
    &\letin{b = g \, y} \\ &\quad \letin{a = f \, x} (a, b)
  \end{aligned}
\end{equation*}
%
The order of $f \, x$ and $g \, y$ matters when both computations cause computational
effects. The desugaring phase avoids making a decision by using the \emph{simultaneous}
$\kord{let}$ statement
%
\begin{equation*}
  \kpre{let} a = f \, x \kop{and} b = g \, y \kop{in} (a, b)
\end{equation*}
%
which leaves open the possibility of various compiler optimizations. The prototype simply
evaluates simultaneous bindings in the order they are given, and a command-line
option enables sequencing warnings about possible unexpected order of effects.
%
It could be argued that the warnings should actually be errors, but we allow some slack until
we have an effect system that can detects harmless instances of simultaneous binding.

For one-off handlers, \eff provides an inline syntax so that one can write

\begin{center}
\begin{tabular}{ccc}
\begin{source}
handle
  $c$
with
  $\cdots$
\end{source}&
\quad instead of \quad \hbox{}&
\begin{source}
with
  handler
    $\cdots$
handle
  $c$
\end{source}
\end{tabular}
\end{center}
%
Additionally, the \inline{val} and \inline{finally} clauses may be omitted, in which case
they are assumed to be the identities.

%%% Local Variables: 
%%% mode: latex
%%% TeX-master: "eff"
%%% End: 

% !TEX root = runners-in-action.tex

\section{Related work}
\label{sect:relatedwork}

Comodels and (ordinary) runners have been used as a natural
model of stateful top-level behaviour.
%
For instance, Plotkin and Power~\cite{Plotkin:TensorsOfModels} have given a treatment of operational 
semantics using the tensor product of a model and a comodel.
%
Recently, Katsumata, Rivas, and Uustalu have generalised this interaction of models and comodels
to monads and comonads~\cite{Katsumata:InteractionLaws}.
%
An early version of \pl{Eff}~\cite{Bauer:AlgebraicEffects} implemented \emph{resources},
which were a kind of stateful runners, although they lacked satisfactory theory.
%
Uustalu~\cite{Uustalu:Runners} has pointed out that runners are the additional
structure that one has to impose on state
to run algebraic effects statefully.
%
Møgelberg and Staton's~\cite{Mogelberg:LinearUsageOfState} linear-use state-passing
translation also relies on equipping the state with a comodel
structure for the effects at hand. Our runners arise
when their setup is specialised to a certain Kleisli adjunction.

Our use of kernel state is analogous to the use
of parameters in parameter-passing 
handlers~\cite{Plotkin:HandlingEffects}: their 
$\tmkw{return}$ clause also provides a form of finalisation, as the 
final value of the parameter is available.
There is however no guarantee of
finalisation happening because handlers need not use the
continuation linearly.

The need to tame the excessive generality of handlers, and willingness to
give it up in exchange for efficiency and predictability, has recently
been recognised by \pl{Multicore OCaml}'s implementors, who
have observed that in practice most handlers resume 
continuations precisely once~\cite{Dolan:MulticoreOCaml}. In exchange for impressive efficiency, they require
continuations to be used linearly by default, whereas discarding and copying
must be done explicitly, incurring additional cost.
%
Leijen~\cite{Leijen:Finalisation} has extended 
handlers in \pl{Koka} with a $\tmkw{finally}$ clause, whose 
semantics ensures that finalisation happens whenever a handler
discards its continuation.
Leijen also added an $\tmkw{initially}$ clause to
parameter-passing handlers, which is used to compute the initial value of the parameter before handling, but that 
gets executed again every time the handler resumes its continuation.


%%% Local Variables:
%%% mode: latex
%%% TeX-master: "runners-in-action"
%%% End:

\section{Directions and Conclusions}
\label{sec:conclusions}

We have presented our initial exploration of \DV\ across a wide range of tasks and domains, providing supporting evidence to the claim that \DV's abilities are comparable to human-level for many of them. This conclusion is consistent with the findings by OpenAI presented in \cite{gpt4}. A primary goal of our experiments is to give a preliminary assessment of \DV's {\em intelligence}, which is an arduous task given the lack of formal definition for this concept, especially for artificial systems. We hope that our exploration provides a useful and necessary first step to appreciate the remarkable capabilities and challenges of {\DV}, and that it opens up new opportunities for developing more formal and comprehensive methods for testing and analyzing future AI systems with such broad intelligence. The capabilities of the model, which have been demonstrated above, both in terms of depth and generality, suggest that the machine learning community needs to move beyond classical benchmarking via structured datasets and tasks, and that the evaluation of the capabilities and cognitive abilities of those new models have become much closer in essence to the task of evaluating those of a human rather than those of a narrow AI model. We hope our investigation stimulates further research on {\DV} and similar systems, both in terms of exploring new applications and domains, and in terms of understanding the mechanisms and principles that underlie their intelligence.
\newline

The central claim of our work is that \DV\ attains a form of \emph{general} intelligence, indeed showing {\em sparks of artificial general intelligence}. This is demonstrated by its core mental capabilities (such as reasoning, creativity, and deduction), its range of topics on which it has gained expertise (such as literature, medicine, and coding), and the variety of tasks it is able to perform (e.g., playing games, using tools, explaining itself, ...). A lot remains to be done to create a system that could qualify as a complete AGI. We conclude this paper by discussing several immediate next steps, regarding defining AGI itself, building some of missing components in LLMs for AGI, as well as gaining better understanding into the origin of the intelligence displayed by the recent LLMs.

\subsection{Definitions of intelligence, AI, and AGI} \label{sec:otherdefinitions}
In this paper, we have used the 1994 definition of intelligence by a group of psychologists \cite{gottfredson1997mainstream} as a guiding framework to explore \DV's artificial intelligence. This definition captures some important aspects of intelligence, such as reasoning, problem-solving, and abstraction, but it is also vague and incomplete. It does not specify how to measure or compare these abilities. Moreover, it may not reflect the specific challenges and opportunities of artificial systems, which may have different goals and constraints than natural ones. Therefore, we acknowledge that this definition is not the final word on intelligence, but rather a useful starting point for our investigation. There is a rich and ongoing literature that attempts to propose more formal and comprehensive definitions of intelligence, artificial intelligence, and artificial general intelligence \cite{goertzel2014artificial, chollet2019measure}, but none of them is without problems or controversies.
For instance, Legg and Hutter \cite{legg2008machine} propose a goal-oriented definition of artificial general intelligence: Intelligence measures an agent’s ability to achieve goals in a wide range of environments. However, this definition does not necessarily capture the full spectrum of intelligence, as it excludes passive or reactive systems that can perform complex tasks or answer questions without any intrinsic motivation or goal. One could imagine as an artificial general intelligence, a brilliant oracle, for example, that has no agency or preferences, but can provide accurate and useful information on any topic or domain. Moreover, the definition around achieving goals in a wide range of environments also implies a certain degree of universality or optimality, which may not be realistic (certainly human intelligence is in no way universal or optimal). The need to recognize the importance of priors (as opposed to {\em universality}) was emphasized in the definition put forward by Chollet in \cite{chollet2019measure} which centers intelligence around skill-acquisition efficiency, or in other words puts the emphasis on a single component of the 1994 definition: learning from experience (which also happens to be one of the key weaknesses of LLMs). Another candidate definition of artificial general intelligence from Legg and Hutter \cite{legg2007universal} is: a system that can do anything a human can do. However, this definition is also problematic, as it assumes that there is a single standard or measure of human intelligence or ability, which is clearly not the case. Humans have different skills, talents, preferences, and limitations, and there is no human that can do everything that any other human can do. Furthermore, this definition also implies a certain anthropocentric bias, which may not be appropriate or relevant for artificial systems. While we do not adopt any of those definitions in the paper, we recognize that they provide important angles on intelligence. For example, whether intelligence can be achieved without any agency or intrinsic motivation is an important philosophical question. Equipping LLMs with agency and intrinsic motivation is a fascinating and important direction for future work. With this direction of work, great care would have to be taken on alignment and safety per a system's abilities to take autonomous actions in the world and to perform autonomous self-improvement via cycles of learning. We discuss a few other crucial missing components of LLMs next.

\subsection{On the path to more general artificial intelligence}
%We have provided evidence supporting that claim that {\DV} performance on a wide range of tasks is comparable to human-level abilities. We have argued that the model attains a form of \emph{general} intelligence in terms of core mental capabilities (such as reasoning, creativity, and deduction), in terms of the range of topics on which is has gained expertise (such as literature, medicine, and coding), and in terms of the variety of tasks it is able to perform (e.g., playing games, using tools, explaining itself, ...). We have also shown that {\DV} can generate and understand content that combines different topics, skills, and modalities, demonstrating its flexibility and creativity and that, despite being trained purely on text, it demonstrates remarkable capabilities in a variety of modalities such as vision. We have compared {\DV}'s performance to those of previous large language models (LLMs), most notably ChatGPT \cite{gpt3}, and we have found that {\DV} is far superior in terms of generality, creativity, and closeness to human-level intelligence. 

%As we allude to in the title of the paper, this work explores a ``first contact" with {\DV} and its potential descendants, rather than a comprehensive evaluation of the model's intelligence. We hope that our exploration provides a useful and necessary first step to appreciate the remarkable capabilities and challenges of {\DV}, and that it opens up new opportunities for developing more formal and comprehensive methods for testing and analyzing future AGI systems. The capabilities of the model, which have been demonstrated above, both in terms of depth and generality, suggest that the machine learning community needs to move beyond classical benchmarking via structured datasets and tasks, and that the evaluation of the capabilities and cognitive abilities of those new models have become much closer in essence to the task of evaluating those of a human rather than those of a narrow AI model. We hope our investigation stimulates further research on {\DV} and similar systems, both in terms of exploring new applications and domains, and in terms of understanding the mechanisms and principles that underlie their intelligence.

%We have also identified some of the main drawbacks of \DV, and we have discussed how they might be addressed in future work. These drawbacks include:
Some of the areas where \DV\ (and LLMs more generally) should be improved to achieve more general intelligence include (note that many of them are interconnected):
\begin{itemize}
    \item \textbf{Confidence calibration:} The model has trouble knowing when it should be confident and when it is just guessing. It both makes up facts that have not appeared in its training data, and also exhibits inconsistencies between the generated content and the prompt, which we referred to as {\em open-domain} and {\em closed-domain} hallucination in Figure \ref{fig:hallucination}. These hallucinations can be stated in a confident and persuasive manner that can be difficult to detect. Thus, such generations can lead to errors, and also to confusion and mistrust. While hallucination is a good thing when generating creative content, reliance on factual claims made by a model with hallucinations can be costly, especially for uses in high-stakes domains such as healthcare. There are several complementary ways to attempt to address hallucinations. One way is to improve the calibration of the model (either via prompting or fine-tuning) so that it either abstains from answering when it is unlikely to be correct or provides some other indicator of confidence that can be used downstream. Another approach, that is suitable for mitigating open-domain hallucination, is to insert information that the model lacks into the prompt, for example by allowing the model to make calls to external sources of information, such as a search engine as in Section \ref{sec:affordances}. For closed-domain hallucination the use of additional model computation through post-hoc checks is also promising, see Figure \ref{fig:hallucination} for an example. Finally, building the user experience of an application with the possibility of hallucinations in mind can also be part of an effective mitigation strategy. %Other directions include developing and refining mechanisms that endow systems with well-calibrated likelihoods that its generations are grounded, or, more directly, the likelihood that it is hallucinating versus relying upon and communicating content that it has learned from its training data. 
    \item \textbf{Long-term memory:} The model's context is very limited (currently 8000 tokens, but not scalable in terms of computation), it operates in a ``stateless" fashion and there is no obvious way to teach the model new facts. In fact, it is not even clear whether the model is able to perform tasks which require an evolving memory and context, such as reading a book, with the task of following the plot and understanding references to prior chapters over the course of reading.
    \item \textbf{Continual learning:} The model lacks the ability to update itself or adapt to a changing environment. The model is fixed once it is trained, and there is no mechanism for incorporating new information or feedback from the user or the world. One can fine-tune the model on new data, but this can cause degradation of performance or overfitting. Given the potential lag between cycles of training, the system will often be out of date when it comes to events, information, and knowledge that came into being after the latest cycle of training.
    \item \textbf{Personalization:} Some of the applications require the model to be tailored to a specific organization or end user. The system may need to acquire knowledge about the workings of an organization or the preferences of an individual. And in many cases, the system would need to adapt in a personalized manner over periods of time with specific changes linked to the dynamics of people and organizations. For example, in an educational setting, there would be an expectation of the need for the system to understand particular learning styles as well as to adapt over time to a student's progress with comprehension and prowess. The model does not have any way to incorporate such personalized information into its responses, except by using meta-prompts, which are both limited and inefficient. 
    \item \textbf{Planning and conceptual leaps:} As suggested by the examples in Section \ref{sec:limitations}, the model exhibits difficulties in performing tasks that require planning ahead or that require a ``Eureka idea" constituting a discontinuous conceptual leap in the progress towards completing a task. In other words, the model does not perform well on tasks that require the sort of conceptual leaps of the form that often typifies human genius.  
    \item \textbf{Transparency, interpretability and consistency:} Not only does the model hallucinate, make up facts and produce inconsistent content, but it seems that the model has no way of verifying whether or not the content that it produces is consistent with the training data, or whether it's self-consistent. While the model is often able to provide high-quality post-hoc explanations for its decisions (as demonstrated in Section \ref{sec:explainability}), using explanations to verify the process that led to a certain decision or conclusion only works when that process is accurately modeled and a sufficiently powerful explanation process is also accurately modeled (Section \ref{sec:explainability}). Both of these conditions are hard to verify, and when they fail there are inconsistencies between the model's decisions and its explanations. Since the model does not have a clear sense of its own limitations it makes it hard to establish trust or collaboration with the user without extensive experimentation in a narrow domain.
    \item \textbf{Cognitive fallacies and irrationality:} The model seems to exhibit some of some of the limitations of human knowledge and reasoning, such as cognitive biases and irrationality (such as biases of confirmation, anchoring, and base-rate neglect) and statistical fallacies. The model may inherit some of the biases, prejudices, or errors that are present in its training data, which may reflect the distribution of opinions or perspectives linked to subsets of the population or  larger common views and assessments. 
     \item \textbf{Challenges with sensitivity to inputs:} The model's responses can be very sensitive to details of the framing or wording of prompts and their sequencing in a session. Such non-robustness suggests that significant effort and experimentation is often required with engineering prompts and their sequencing and that uses in the absence of such investments of time and effort by people can lead to suboptimal and non-aligned inferences and results. 
\end{itemize}

A limitation of our exploration is the absence of a clear distinction between drawbacks founded in the way that the reinforcement learning step (RLHF) was carried out, versus drawbacks which are fundamentally inherent in the larger architecture and methodology. For example, it is not clear to what extent the hallucination problem can be addressed via a refined reinforcement learning step or via a focused effort to introduce new forms of calibration about the likelihoods of the veracity of alternative inferences that the system can compute and consider in its generations (see also \cite{gpt4} for more discussion on this). To draw an analogy to humans, cognitive biases and irrational thinking may be based in artifacts of our culture as well as to limitations in our cognitive capabilities. Pursuing better understandings of the sources and potential solutions to challenges of hallucination in \DV, will benefit from studies that compare several versions of the RL stage over the same architecture.
\newline

A broader question on the identified limitations is: which of the aforementioned drawbacks can be mitigated within the scope of next word prediction? Is it simply the case that a bigger model and more data will fix those issues, or does the architecture need to be modified, extended, or reformulated? Potential extensions to next word prediction include the following:

\begin{itemize}
    \item External calls by the model to components and tools such as a calculator, a database search or code execution, as suggested in Section \ref{sec:affordances}. 
    \item A richer, more complex ``slow-thinking" deeper mechanism that oversees the ``fast-thinking" mechanism of next word prediction. Such an approach could allow the model to perform long-term planning, exploration, or verification, and to maintain a working memory or a plan of action. The slow-thinking mechanism would use the next word prediction model as a subroutine, but it would also have access to external sources of information or feedback, and it would be able to revise or correct the outputs of the fast-thinking mechanism.
    \item Integration of long-term memory as an inherent part of the architecture, perhaps in the sense that both the input and output of the model will include, in addition to the tokens representing the text, a vector which represents the context.
    \item Going beyond single-word prediction: Replacing the sequence of tokens by a hierarchical structure, where higher-level parts of the text such as sentences, paragraphs or ideas are represented in the embedding and where the content is generated in a top-down manner. It is unclear whether richer predictions about the sequencing and interdependency of such higher-level concepts might emerge from large-scale compute and data centered on a next-word--prediction paradigm.
\end{itemize}

%In conclusion, we have worked to demonstrate that {\DV} has remarkable capabilities that challenge many of the recent assumptions and expectations within the AI community. We have also shown that {\DV} is by no means a perfect or complete AGI system, and that it has many limitations and biases that need to be addressed and understood. We hope that our exploration will inspire and inform further research on {\DV} and similar systems, both in terms of exploring their potential applications and domains, and in terms of understanding foundational mechanisms and potentially with identifying principles of intelligence. We believe that {\DV} represents a paradigm shift in the field of computer science and beyond, and that the model and its capabilities frame new questions, possibilities, and horizons for the field and for the advancement of human capabilities and well-being.

\subsection{What is actually happening?} \label{sec:whatsgoingon}
Our study of {\DV} is entirely phenomenological: We have focused on the surprising things that {\DV} can do, but we do not address the fundamental questions of why and how it achieves such remarkable intelligence. How does it reason, plan, and create? Why does it exhibit such general and flexible intelligence when it is at its core merely the combination of simple algorithmic components---gradient descent and large-scale transformers with extremely large amounts of data? These questions are part of the mystery and fascination of LLMs, which challenge our understanding of learning and cognition, fuel our curiosity, and motivate deeper research. Key directions include ongoing research on the phenomenon of emergence in LLMs (see \cite{wei2022emergent} for a recent survey). Yet, despite intense interest in questions about the capabilities of LLMs, progress to date has been quite limited with only toy models where some phenomenon of emergence is proved \cite{barak2022hidden, ahn2022learning,jelassi2022vision}. One general hypothesis \cite{olah2020zoom} is that the large amount of data (especially the  diversity of the content) forces neural networks to learn generic and useful ``neural circuits'', such as the ones discovered in \cite{olsson2022context, zhang2022unveiling, liu2022transformers}, while the large size of models provide enough redundancy and diversity for the neural circuits to specialize and fine-tune to specific tasks. 
%The Mixture of Experts (MoE) layers in modern LLMs can also contribute to the generality of the model~\cite{chen2022towards}. 
Proving these hypotheses for large-scale models remains a challenge, and, moreover, it is all but certain that the conjecture is only part of the answer. On another direction of thinking, the huge size of the model could have several other benefits, such as making gradient descent more effective by connecting different minima \cite{venturi2019spurious} or by simply enabling smooth fitting of high-dimensional data \cite{pmlr-v49-eldan16, NEURIPS2021_f197002b}. Overall, elucidating the nature and mechanisms of AI systems such as {\DV} is a formidable challenge that has suddenly become important and urgent.
\newline

\paragraph{Acknowledgements.} We thank OpenAI for creating such a marvelous tool and giving us early access to experience it. We also thank Miles Brundage at OpenAI, and the numerous people at Microsoft, who have provided thoughtful feedback on this work.


%% BIBLIOGRAPHY
\bibliography{references}



%% CC-BY PARAGRAPH REQUIRED BY THE PUBLISHER
\vfill

{\small\medskip\noindent{\bf Open Access} This chapter is licensed under the terms of the Creative Commons\break Attribution 4.0 International License (\url{http://creativecommons.org/licenses/by/4.0/}), which permits use, sharing, adaptation, distribution and reproduction in any medium or format, as long as you give appropriate credit to the original author(s) and the source, provide a link to the Creative Commons license and indicate if changes were made.}

{\small \spaceskip .28em plus .1em minus .1em The images or other third party material in this chapter are included in the chapter's Creative Commons license, unless indicated otherwise in a credit line to the material.~If material is not included in the chapter's Creative Commons license and your intended\break use is not permitted by statutory regulation or exceeds the permitted use, you will need to obtain permission directly from the copyright holder.}

\medskip\noindent\includegraphics{cc_by_4-0.eps}



% APPENDIX
\appendix
\makeatletter\def\@seccntformat#1{\appendixname~\csname the#1\endcsname: }\makeatother
% !TEX root = runners-in-action.tex

\section{Typing rules of $\lambdacoop$}
\label{sec:typing-rules}

In this appendix we give the complete typing rules for $\lambdacoop$. We refer to
\cref{fig:lambdacoop-types,fig:lambdacoop-terms} for the syntax of types, values, 
and user and kernel computations.
%
For each operation symbol $\op \in \Ops$, we assume a given and fixed operation signature
%
\begin{equation*}
  \op : \tysigop{A_\op}{B_\op}{E_\op},
\end{equation*}
%
and for each ground constant $\tmconst{f}$, we assume 
a signature $\tmconst{f} : (A_1,\ldots,A_n) \to B$, 
both of which the typing rules refer to without further ado.
%
Values, and user and kernel computations each have a typing and a subtyping judgement of the
form
%
\begin{align*}
  &\Gamma \types V : X,
& &\Gamma \types M : \tyuser{X}{\Ueff},
& &\Gamma \types K : \tykernel{X}{\Keff},\\
  &X \sub Y,
& &\tyuser{X}{\Ueff} \sub \tyuser{Y}{\Veff},
& &\tykernel{X}{\Keff} \sub \tykernel{Y}{\Leff}.
\end{align*}
%
where $\Gamma$ is the customary typing context assigning value types to variables. 
The subtyping rules are given in \cref{fig:subtyping}, and the 
typing rules in \cref{fig:typing-values,fig:typing-user,fig:typing-kernel}.


%%%%%%%%%%%%%%%%%%%%%% SUBTYPING RULES %%%%%%%%%%%%%%%%%%%%%%%%%%%%%%%%%%
\begin{figure}[p]
  \small
  \begin{mathpar}
    \coopinfer{Sub-Ground}{ }{A \sub A}

    \coopinfer{Sub-Product}{
      X \sub X' \\
      Y \sub Y'
    }{
      X \times Y \sub X' \times Y'
    }

    \coopinfer{Sub-Sum}{
      X \sub X' \\
      Y \sub Y'
    }{
      X + Y \sub X' + Y'
    }

    \coopinfer{Sub-UserFun}{
      X' \sub X \\
      \tyuser{Y}{\Ueff} \sub \tyuser{Y'}{\Ueff'}
    }{
      \tyfun{X}{\tyuser{Y}{\Ueff}} \sub \tyfun{X'}{\tyuser{Y'}{\Ueff'}}
    }

    \coopinfer{Sub-KernelFun}{
      X' \sub X \\
      \tykernel{Y}{\Keff} \sub \tykernel{Y'}{\Keff'}
    }{
      \tyfunK{X}{\tykernel{Y}{\Keff}} \sub \tyfunK{X'}{\tykernel{Y'}{\Keff'}}
    }

    %%% SELECTED RULE, KEEP IN SYNC WITH fig:typing-selected
    \coopinfer{Sub-Runner}{
      \sig_1' \subseteq \sig_1 \\
      \sig_2 \subseteq \sig_2' \\
      S \subseteq S' \\
      C \equiv C'
    }{
      \tyrunner{\sig_1}{\sig_2}{S}{C} \sub \tyrunner{\sig_1'}{\sig_2'}{S'}{C'}
    }

    \coopinfer{Sub-User}{
      X \sub X' \\
      \sig \subseteq \sig' \\
      E \subseteq E'
    }{
      \tyuser{X}{(\sig, E)} \sub \tyuser{X'}{(\sig', E')}
    }

    %%% SELECTED RULE, KEEP IN SYNC WITH fig:typing-selected
    \coopinfer{Sub-Kernel}{
      X \sub X' \\
      \sig \subseteq \sig' \\
      E \subseteq E' \\\\
      S \subseteq S' \\
      C \equiv C'
    }{
      \tykernel{X}{(\sig, E, S, C)} \sub \tykernel{X'}{(\sig', E', S', C')}
    }

    \coopinfer{Subsume-Value}{
      \Gamma \types V : X \\
      X \sub X'
    }{
      \Gamma \types V : X'
    }

    \coopinfer{Subsume-User}{
      \Gamma \types M : \tyuser{X}{\Ueff} \\
      \tyuser{X}{\Ueff} \sub \tyuser{X'}{\Ueff'}
    }{
      \Gamma \types M : \tyuser{X'}{\Ueff'}
    }

    \coopinfer{Subsume-Kernel}{
      \Gamma \types K : \tykernel{X}{\Keff} \\
      \tykernel{X}{\Keff} \sub \tykernel{X'}{\Keff'}
    }{
      \Gamma \types M : \tykernel{X'}{\Keff'}
    }
  \end{mathpar}
  \caption{Subtyping and subsumption rules.}
  \label{fig:subtyping}
\end{figure}

%%%%%%%%%%%%%%%%% VALUE RULES %%%%%%%%%%%%%%%%%%%%
\begin{figure}[p]
  \centering
  \small
\begin{mathpar}

  \coopinfer{TyValue-Var}{
    \Gamma(x) \equiv X
  }{
    \Gamma \types x : X
  }
  
  \coopinfer{TyValue-Const}{
    (\Gamma \types V_i : A_i)_{1 \leq i \leq n}
  }{
    \Gamma \types \tmconst{f}(V_1, \ldots , V_n) : B
  }

  \coopinfer{TyValue-Unit}{
  }{
    \Gamma \types \tmunit : \tyunit
  }

  \coopinfer{TyValue-Pair}{
    \Gamma \types V : X \\
    \Gamma \types W : Y
  }{
    \Gamma \types \tmpair{V}{W} : \typrod{X}{Y}
  }

  \coopinfer{TyValue-Inl}{
    \Gamma \types V : X
  }{
    \Gamma \types \tminl[X,Y]{V} : X + Y
  }

  \coopinfer{TyValue-Inr}{
    \Gamma \types W : Y
  }{
    \Gamma \types \tminr[X,Y]{W} : X + Y
  }

  \coopinfer{TyValue-UserFun}{
    \Gamma, x \of X \types M : \tyuser{Y}{\Ueff}
  }{
    \Gamma \types \tmfun{x : X}{M} : \tyfun{X}{\tyuser{Y}{\Ueff}}
  }

  \coopinfer{TyValue-KernelFun}{
    \Gamma, x \of X \types K : \tykernel{Y}{\Keff}
  }{
    \Gamma \types \tmfunK{x : X}{K} : \tyfunK{X}{\tykernel{Y}{\Keff}}
  }

  \coopinfer{TyValue-Runner}{
    \big(
      \Gamma, x \of A_\op \types K_\op : \tykernel{B_\op}{(\sig', E_\op, S, C)}
    \big)_{\op \in \sig}
  }{
    \Gamma \types
    \tmrunner{(\tm{op}\,x \mapsto K_{\tm{op}})_{\tm{op} \in \Sigma}}{C} :
    \tyrunner{\sig}{\sig'}{S}{C}
  }
\end{mathpar}
  \caption{Value typing rules.}
  \label{fig:typing-values}
\end{figure}

%%%%%%%%%%%%%%%%% USER RULES %%%%%%%%%%%%%%%%%%%%
\begin{figure}[tp]
  \centering
  \small
\begin{mathpar}
  \coopinfer{TyUser-Return}{
    \Gamma \types V : X
  }{
    \Gamma \types \tmreturn{V} : \tyuser{X}{\Ueff}
  }

  \coopinfer{TyUser-Apply}{
    \Gamma \types V : \tyfun{X}{\tyuser{Y}{\Ueff}} \\
    \Gamma \types W : X
  }{
    \Gamma \types \tmapp{V}{W} : \tyuser{Y}{\Ueff}
  }

  %%% SELECTED RULE, KEEP IN SYNC WITH fig:typing-selected
  \coopinfer{TyUser-Try}{
    \Gamma \types M : \tyuser{X}{(\sig,E)}
    \\
    \Gamma, x \of X \types N : \tyuser{Y}{(\sig,E')}
    \\
    \big(
      \Gamma \types N_e : \tyuser{Y}{(\sig,E')}
    \big)_{e \in E}
  }{
    \Gamma \types
    \tmtry{M}{
        \{ \tmreturn{x} \mapsto N,
           (\tmraise{e} \mapsto N_e)_{e \in E} \}
        }
    : \tyuser{Y}{(\sig,E')}
  }

  \coopinfer{TyUser-MatchPair}{
    \Gamma \types V : \typrod{X}{Y} \\
    \Gamma, x \of X, y \of Y \types M : \tyuser{Z}{\Ueff}
  }{
    \Gamma \types \tmmatch{V}{\tmpair{x}{y} \mapsto M} : \tyuser{Z}{\Ueff}
  }

  \coopinfer{TyUser-MatchEmpty}{
    \Gamma \types V : \tyempty
  }{
    \Gamma \types \tmmatch[Z]{V}{} : \tyuser{Z}{\Ueff}
  }

  \coopinfer{TyUser-MatchSum}{
    \Gamma \types V : X + Y \\
    \Gamma, x \of X \types M : \tyuser{Z}{\Ueff} \\
    \Gamma, y \of Y \types N : \tyuser{Z}{\Ueff} \\
  }{
    \Gamma \types \tmmatch{V}{\tminl{x} \mapsto M, \tminr{y} \mapsto N} : \tyuser{Z}{\Ueff}
  }

  %%% SELECTED RULE, KEEP IN SYNC WITH fig:typing-selected
  \coopinfer{TyUser-Op}{
    \Ueff \equiv (\sig,E) \\
    \op \in \sig \\
    \Gamma \types V : A_\op \\\\
    \Gamma, x \of B_\op \types M : \tyuser{X}{\Ueff} \\
    \big(
      \Gamma \vdash N_e : \tyuser{X}{\Ueff}
    \big)_{e \in E_\op}
  }{
    \Gamma \types \tmop{op}{X}{V}{\tmcont x M}{\tmexccont N e {E_\op}} : \tyuser{X}{\Ueff}
  }

  \coopinfer{TyUser-Raise}{
     e \in E
  }{
     \Gamma \types \tmraise[X]{e} : \tyuser{X}{(\sig, E)}
  }

  %%% SELECTED RULE, KEEP IN SYNC WITH fig:typing-selected
  \coopinfer{TyUser-Run}{
    F \equiv
    \{ \tmreturn{x} \at c \mapsto N,
       (\tmraise{e} \at c \mapsto N_e)_{e \in E},
       (\tmkill{s} \mapsto N_s)_{s \in S}
    \}
    \\\\
    \Gamma \types V : \tyrunner{\sig}{\sig'}{S}{C} \\
    \Gamma \types W : C \\\\
    \Gamma \types M : \tyuser{X}{(\sig, E)} \\
    \Gamma, x \of X, c \of C \types N : \tyuser{Y}{(\sig', E')} \\
    \big(
       \Gamma, c \of C \types N_e : \tyuser{Y}{(\sig', E')}
    \big)_{e \in E} \\
    \big(
       \Gamma \types N_s : \tyuser{Y}{(\sig', E')}
    \big)_{s \in S} \\
  }{
    \Gamma \types \tmrun{V}{W}{M}{F} : \tyuser{Y}{(\sig', E')}
  }

  %%% SELECTED RULE, KEEP IN SYNC WITH fig:typing-selected
  \coopinfer{TyUser-Kernel}{
    F \equiv
    \{ \tmreturn{x} \at c \mapsto N,
       (\tmraise{e} \at c \mapsto N_e)_{e \in E},
       (\tmkill{s} \mapsto N_s)_{s \in S}
    \}
    \\\\
    \Gamma \types K : \tykernel{X}{(\sig, E, S, C)} \\
    \Gamma \types W : C \\
    \Gamma, x \of X, c \of C \types N : \tyuser{Y}{(\sig, E')} \\
    \big(
      \Gamma, c \of C \types N_e : \tyuser{Y}{(\sig, E')}
    \big)_{e \in E} \\
    \big(
      \Gamma \types N_s : \tyuser{Y}{(\sig, E')}
    \big)_{s \in S} \\
  }{
    \Gamma \types \tmkernel{K}{W}{F} : \tyuser{Y}{(\sig, E')}
  }
\end{mathpar}
  \caption{User typing rules.}
  \label{fig:typing-user}
\end{figure}


%%%%%%%%%%%%%%%%% KERNEL RULES %%%%%%%%%%%%%%%%%%%%
\begin{figure}[tp]
  \centering
  \small
\begin{mathpar}

  \coopinfer{TyKernel-Return}{
    \Gamma \types V : X
  }{
    \Gamma \types \tmreturn[C]{V} : \tykernel{X}{(\sig, E, S, C)}
  }

  \coopinfer{TyKernel-Apply}{
    \Gamma \types V : \tyfun{X}{\tykernel{Y}{\Keff}} \\
    \Gamma \types W : X
  }{
    \Gamma \types \tmapp{V}{W} : \tykernel{Y}{\Keff}
  }

  \coopinfer{TyKernel-Try}{
    \Gamma \types K : \tykernel{X}{(\sig, E, S, C)}
    \\
    \Gamma, x \of X \types L : \tykernel{Y}{(\sig, E', S, C)}
    \\
    \big(
      \Gamma \types L_e : \tykernel{Y}{(\sig, E', S, C)}
    \big)_{e \in E}
  }{
    \Gamma \types
    \tmtry{K}{
        \{ \tmreturn{x} \mapsto L,
           (\tmraise{e} \mapsto L_e)_{e \in E} \}
        }
    : \tykernel{Y}{(\sig, E', S, C)}
  }

  \coopinfer{TyKernel-MatchPair}{
    \Gamma \types V : \typrod{X}{Y} \\
    \Gamma, x \of X, y \of Y \types K : \tykernel{Z}{\Keff}
  }{
    \Gamma \types \tmmatch{V}{\tmpair{x}{y} \mapsto K} : \tykernel{Z}{\Keff}
  }

  \coopinfer{TyKernel-MatchEmpty}{
    \Gamma \types V : \tyempty
  }{
    \Gamma \types \tmmatch[Z \at C]{V}{} : \tykernel{Z}{(\sig, E, S, C)}
  }

  \coopinfer{TyKernel-MatchSum}{
    \Gamma \types V : X + Y \\
    \Gamma, x \of X \types K : \tykernel{Z}{\Keff} \\
    \Gamma, y \of Y \types L : \tykernel{Z}{\Keff} \\
  }{
    \Gamma \types \tmmatch{V}{\tminl{x} \mapsto K, \tminr{y} \mapsto L} : \tykernel{Z}{\Keff}
  }

  %%% SELECTED RULE, KEEP IN SYNC WITH fig:typing-selected
  \coopinfer{TyKernel-Op}{
    \Keff \equiv (\sig, E, S, C) \\
    \op \in \sig \\
    \Gamma \types V : A_\op \\\\
    \Gamma, x \of B_\op \types K : \tykernel{X}{\Keff} \\
    \big(
      \Gamma \vdash L_e : \tykernel{X}{\Keff}
    \big)_{e \in E_\op}
  }{
    \Gamma \types \tmop{op}{X}{V}{\tmcont x K}{\tmexccont L e {E_\op}} : \tykernel{X}{\Keff}
  }

  \coopinfer{TyKernel-Raise}{
     e \in E
  }{
     \Gamma \types \tmraise[X \at C]{e} : \tykernel{X}{(\sig, E, S, C)}
  }

  \coopinfer{TyKernel-Kill}{
     s \in S
  }{
     \Gamma \types \tmkill[X \at C]{s} : \tykernel{X}{(\sig, E, S, C)}
  }

  \coopinfer{TyKernel-Getenv}{
    \Gamma, c \of C \types K : \tykernel{X}{(\sig, E, S, C)}
  }{
    \Gamma \types \tmgetenv[C]{\tmcont c K} : \tykernel{X}{(\sig, E, S, C)}
  }

  \coopinfer{TyKernel-Setenv}{
    \Gamma \types V : C \\
    \Gamma \types K : \tykernel{X}{(\sig, E, S, C)}
  }{
    \Gamma \types \tmsetenv{V}{K} : \tykernel{X}{(\sig, E, S, C)}
  }

  %%% SELECTED RULE, KEEP IN SYNC WITH fig:typing-selected
  \coopinfer{TyKernel-User}{
   \Keff \equiv (\sig, E', S, C) \\
   \Gamma \types M : \tyuser{X}{(\sig, E)} \\\\
   \Gamma, x \of X \types K : \tykernel{Y}{\Keff} \\
   \big(
     \Gamma \types L_e : \tykernel{Y}{\Keff}
   \big)_{e \in E}
  }{
    \Gamma \types
    \tmuser{M}{
      \{ \tmreturn{x} \mapsto K,
         (\tmraise{e} \mapsto L_e)_{e \in E}
      \}
    }
    : \tykernel{Y}{\Keff}
  }
\end{mathpar}
  \caption{Kernel typing rules.}
  \label{fig:typing-kernel}
\end{figure}


%%% Local Variables:
%%% mode: latex
%%% TeX-master: "runners-in-action"
%%% End:

% !TEX root = runners-in-action.tex

\section{Equational theory of $\lambdacoop$}
\label{sec:appendix-equational-rules}

Values, user and kernel computations each have an equality judgement
\begin{equation*}
\Gamma \types V \equiv W : X
\qquad
\Gamma \types M \equiv N : \tyuser{X}{\Ueff}
\qquad
\Gamma \types K \equiv L  : \tyuser{X}{\Keff}.
\end{equation*}
%
It is presupposed that we only compare well-typed expressions with the indicated types.
For the most part, the context and the type annotation will play no part in the equation,
and so we shall drop them when no confusion can arise.

The \emph{computational equations} are displayed in \cref{fig:computational-equations-user,fig:computational-equations-kernel}.
These can be read left-to-right as evaluation rules that explain the operational meaning
of computations. The remaining equations are displayed in \cref{fig:other-equations}.
%
We omit standard equations which specify how substitution is performed, as well
as equations stating that equality is a congruence with respect to all the term
formers.


%%%% User computation rules
\begin{figure}[tbp]
  \centering
  \parbox{\textwidth}{
  \small
  %
  \mathtoolsset{original-shortintertext=false,below-shortintertext-sep=0pt,above-shortintertext-sep=0pt}
  \begin{align*}
    \tmapp{(\tmfun{x \of X}{M})}{V} &\equiv M[V/x]
    \\[1ex]
    \tmtry{(\tmreturn{V})}{H} &\equiv N[V/x]
    \\
    \tmtry{(\tmraise[X]{e})}{H} &\equiv N_{e}
    \\
    \tmtry{(
      \tmop{op}{X}{V}{\tmcont x M}{\tmexccont {N'} {e'} {E_\op}}
    )}{H} &\equiv \\
    \tmop{op}{X}{V}{\tmcont x {\tmtry{&M}{H}}}{\left(\tmtry{N'_{e'}}{H}\right)_{e' \in E_\op}}
    \\[1ex]
    \tmmatch{\tmpair{V}{W}}{\tmpair{x}{y} \mapsto M} &\equiv
    M[V/x, W/y]
    \\
    \tmmatch[X]{V}{} &\equiv
    N
    \\
    \mathllap{
      \tmmatch{(\tminl[X,Y]{V})}{\tminl{x} \mapsto M, \tminr{y} \mapsto N} 
    } &\equiv
    M[V/x]
    \\
    \mathllap{
      \tmmatch{(\tminr[X,Y]{W})}{\tminl{x} \mapsto M, \tminr{y} \mapsto N}
    } &\equiv
    N[W/y]
    \\[1ex]
    \tmrun{V}{W}{(\tmreturn{V'})}{F} &\equiv N[V'/x, W/c]
    \\
    \tmrun{V}{W}{(\tmraise[X]{e})}{F} &\equiv N_{e}[W/c]
    \\
    \omit\rlap{$
    \tmkw{using}\; R \at W \;\tmkw{run}\;
      \tmop{op}{X}{V}{\tmcont x M}{\tmexccont {N'} {e'} {E_\op}}
      \;\tmkw{finally}\; F \equiv $}
    \\
    &
    \quad\tmkernel{K_\op[V/x]}{W}{F'}
    \\
    \shortintertext{\hfil \parbox{0.6\textwidth}{
    \abovedisplayskip=0pt
    \belowdisplayskip=0pt
    \begin{equation*}
      \text{where\ } F' \defeq 
      \begin{aligned}[t]
        \{
        &\tmreturn{x} \at c' \mapsto (\tmrun{R}{c'}{M}{F}), \\
        &\left(
            \tmraise{e'} \at c' \mapsto (\tmrun{R}{c'}{N'_{e'}}{F})
         \right)_{e' \in E_\op},\\
        &\left(
            \tmkill{s} \mapsto N_s
          \right)_{s \in S}
      \}
      \end{aligned}
    \end{equation*}
    }}
  \\
    \tmkernel{(\tmreturn[C]{V})}{W}{F} &\equiv N[V/x, W/c] \\
    \tmkernel{(\tmraise[X \at C]{e})}{W}{F} &\equiv N_{e}[W/c] \\
    \tmkernel{(\tmkill[X \at C]{s})}{W}{F} &\equiv N_{s} \\
    \tmkernel{(\tmgetenv[C]{\tmcont c K})}{W}{F} &\equiv \tmkernel{K[W/c]}{W}{F} \\
    \tmkernel{(\tmsetenv{V}{K})}{W}{F} &\equiv \tmkernel{K}{V}{F}
    \\
    \omit\rlap{$
       \tmkernel{\tmop{op}{X}{V}{\tmcont x K}{\tmexccont L {e'} {E_\op}}}{W}{F}
      \equiv $}
    \\
      \tmop{op}{X}{V}{\tmcont x {\tmkernel{K}{W}{&F}}}{
      \left(
         \tmkernel{L_{e'}}{W}{F}
      \right)_{e' \in E_\op}}
  \end{align*}
  %
  Abbreviations:
  %
  \begin{align*}
    F &\defeq
       \{ \tmreturn{x} \at c \mapsto N,
       (\tmraise{e} \at c \mapsto N_e)_{e \in E},
       (\tmkill{s} \mapsto N_s)_{s \in S}
       \}
    \\
    H &\defeq
       \{ \tmreturn{x} \mapsto N,
          (\tmraise{e} \mapsto N_e)_{e \in E}
       \}
    \\
    R &\defeq \tmrunner{(\tm{op}\,x \mapsto K_{\tm{op}})_{\tm{op} \in \sig}}{C}
  \end{align*}
  } % parbox
  \caption{Computational equations (user mode).}
  \label{fig:computational-equations-user}
\end{figure}

%%%% Kernel computation rules
\begin{figure}[tb]
  \centering
  \parbox{\textwidth}{
  \small
  %
  \begin{align*}
    \tmapp{(\tmfunK{x \of X}{K})}{V} &\equiv K[V/x]
    \\[1ex]
    \tmtry{(\tmreturn{V})}{G} &\equiv L[V/x]
    \\
    \tmtry{(\tmraise[X \at C]{e})}{G} &\equiv L_{e}
    \\
    \tmtry{(\tmkill[X \at C]{s})}{G} &\equiv \tmkill[X \at C]{s}
    \\
    \tmtry{(
      \tmop{op}{X}{V}{\tmcont x K}{\tmexccont {L'} {e'} {E_\op}}
    )}{G} &\equiv \\
    \tmop{op}{X}{V}{\tmcont x {\tmtry{&K}{G}}}{\left(\tmtry{L'_{e'}}{G}\right)_{e' \in E_\op}}
    \\
    \tmtry{(\tmgetenv[C]{\tmcont c K})}{G} &\equiv \tmgetenv[C]{\tmcont c {\tmtry{K}{G}}}
    \\
    \tmtry{(\tmsetenv{V}{K})}{G} &\equiv \tmsetenv{V}{\tmtry{K}{G}}
    \\[1ex]
    \tmmatch{\tmpair{V}{W}}{\tmpair{x}{y} \mapsto K} &\equiv
    K[V/x, W/y]
    \\
    \tmmatch[X \at C]{V}{} &\equiv
    K
    \\
    \tmmatch{(\tminl[X,Y]{V})}{\tminl{x} \mapsto K, \tminr{y} \mapsto L} &\equiv
    K[V/x]
    \\
    \tmmatch{(\tminr[X,Y]{W})}{\tminl{x} \mapsto K, \tminr{y} \mapsto L} &\equiv
    L[W/y]
    \\[1ex]
    \tmuser{(\tmreturn{V})}{G} &\equiv L[V/x]
    \\
    \tmuser{(\tmraise[X]{e})}{G} &\equiv L_e
    \\
    \tmuser{(
      \tmop{op}{X}{V}{\tmcont x M}{\tmexccont {N'} {e'} {E_\op}}
    )}{G} &\equiv \\
    \tmop{op}{X}{V}{\tmcont x {\tmuser{&M}{G}}}{\left(\tmuser{N'_{e'}}{G} \right)_{e' \in E_\op}}
  \end{align*}
  %
  Abbreviation: $G \defeq \{ \tmreturn{x} \mapsto L, (\tmraise{e} \mapsto L_e)_{e \in E} \}$
  } % parbox
  \caption{Computational equations (kernel mode).}
  \label{fig:computational-equations-kernel}
\end{figure}


\begin{figure}[tb]
  \centering
  \parbox{\textwidth}{
  \small
  \begin{gather*}
    V \equiv \tmunit : \tyunit \qquad\qquad
    \tmfun{x \of A}{\tmapp{V}{x}} \equiv V \qquad\qquad
    \tmfunK{x \of A}{\tmapp{V}{x}} \equiv V
    \\
    \tmtry{M}{
       \{ \tmreturn{x} \mapsto \tmreturn{x},
          (\tmraise{e} \mapsto \tmraise[X]{e})_{e \in E}
       \}
    }
    \equiv M
    \\
    \tmtry{K}{
       \{ \tmreturn{x} \mapsto \tmreturn{x},
          (\tmraise{e} \mapsto \tmraise[X \at C]{e})_{e \in E}
       \}
    }
    \equiv K
    \\[1ex]
  \begin{aligned}
  \tmgetenv[C]{\tmcont c {\tmsetenv{c}{K}}} &\equiv K
  \\
  \tmsetenv{V}{\tmgetenv[C]{\tmcont c K}} &\equiv
  \tmsetenv{V}{K[V/c]}
  \\
  \tmsetenv{V}{\tmsetenv{W}{K}} &\equiv \tmsetenv{W}{K}
  \\
  \tmgetenv[C]{\tmcont c {\tmkill[X \at C]{s}}} &\equiv
  \tmkill[X \at C]{s}
  \\
  \tmsetenv{V}{\tmkill[X \at C]{s}} &\equiv
  \tmkill[X \at C]{s}
  \\
  \tmgetenv[C]{\tmcont c {\tmop{op}{X}{V}{\tmcont x K}{\tmexccont L e {E_\op}}}}
  &\equiv \\
  \tmop{op}{X}{V}{\tmcont x {\tmgetenv[C]{&\tmcont c K}}}{\left(\tmgetenv[C]{\tmcont c {L_e}}\right)_{e \in E_\op}}
  \\
  \tmsetenv{V}{\tmop{op}{X}{V}{\tmcont x K}{\tmexccont L e {E_\op}}}
  &\equiv \\
  \tmop{op}{X}{&V}{\tmcont x {\tmsetenv{V}{K}}}{\left(\tmsetenv{V}{L_e}\right)_{e \in E_\op}}
  \end{aligned}
  \end{gather*}
  } % parbox
  \caption{Other equations (for $\eta$-expansion and the kernel theory from \cref{sec:user-kernel-monads}).}
  \label{fig:other-equations}
\end{figure}

%%% Local Variables:
%%% mode: latex
%%% TeX-master: "runners-in-action"
%%% End:


\end{document}


